\section{Projective limits of preschemes}
\label{section:IV.8}

\subsection{Introduction}
\label{subsection:IV.8.1}

\oldpage[IV-3]{5}

\begin{env}[8.1.1]
\label{IV.8.1.1}
In this paragraph, we will systematically study the following situation.
Let $\mathrm{I}$ be a filtered preordered increasing set, $(\mathrm{A}_\lambda,\varphi_{\mu\lambda})$ an inductive system of rings having $\mathrm{I}$ for index set, $\mathrm{A}=\varinjlim\mathrm{A}_\lambda$ its inductive limit.
For every $\alpha\in\mathrm{I}$ and every $\mathrm{A}_\alpha$-prescheme $\mathrm{X}_\alpha$, consider the $\mathrm{A}_\lambda$-preschemes $\mathrm{X}_\lambda=\mathrm{X}_\alpha\otimes_{\mathrm{A}_\alpha}\mathrm{A}_\lambda$ for $\lambda\geqslant\alpha$, and the $\mathrm{A}$-prescheme $\mathrm{X}=\mathrm{X}_\alpha\otimes_{\mathrm{A}_\alpha}\mathrm{A}$; it is clear that the preschemes $\mathrm{X}_\lambda$ (for $\lambda\geqslant\alpha$) form a \emph{projective system}, and we will see \sref{IV.8.2.5} that $\mathrm{X}$ is a \emph{projective limit} of this system in the category of preschemes.
We propose to find conditions on $\mathrm{X}_\alpha$ or on the $\mathrm{A}_\lambda$ allowing us to prove properties of the following type:
for $\mathrm{X}$ to have a property $\mathbf{P}$ (for example, the property of being proper over $\mathrm{S}=\operatorname{Spec}(\mathrm{A})$, or irreducible, or connected, etc.), it is necessary and sufficient that there exist an index $\mu\geqslant\alpha$ such that, for every $\lambda\geqslant\mu$, $\mathrm{X}_\lambda$ has (with respect to $\mathrm{S}_\lambda=\operatorname{Spec}(\mathrm{A}_\lambda)$, as the case may be) the same property $\mathbf{P}$.
We will obtain analogous statements for properties of $\sh{O}_\mathrm{X}$-Modules, of $\mathrm{A}$-morphisms of $\mathrm{A}$-preschemes, etc.
We will also show \sref{IV.8.9.1} that the datum of an $\mathrm{A}$-prescheme of finite presentation \sref{IV.1.6.1} is essentially equivalent to the datum of an $\mathrm{A}_\lambda$-prescheme of finite presentation $\mathrm{X}_\lambda$ for some sufficiently large $\lambda$, $\mathrm{X}$ being then \emph{isomorphic to} $\mathrm{X}_\lambda\otimes_{\mathrm{A}_\lambda}\mathrm{A}$.
There are also analogous statements for $\sh{O}_\mathrm{X}$-Modules of finite presentation, their homomorphisms, the $\mathrm{A}$-morphisms of $\mathrm{A}$-preschemes of finite presentation, etc.
\end{env}

\begin{env}[8.1.2]
\label{IV.8.1.2}
The utility of such results will appear, for example, in the following questions:
\begin{enumerate}[label=\emph{\alph*)}]
  \item
  Let $\mathrm{Y}$ be a prescheme, $y$ a point of $\mathrm{Y}$, $(\mathrm{U}_\lambda)$ the projective system of filtered decreasing open affine neighbourhoods of $y$ in $\mathrm{Y}$; if $\mathrm{A}_\lambda$ is the ring of $\mathrm{U}_\lambda$, the $\mathrm{A}_\lambda$ form a filtered increasing inductive system, whose inductive limit $\mathrm{A}$ is the local ring $\sh{O}_y$; moreover, if we denote by $\mathfrak{p}_\lambda$ the prime ideal $\mathfrak{i}_y$ in the ring $\mathrm{A}_\lambda$, the inductive system $(\mathrm{A}_\lambda)$ is cofinal in every inductive system $(\mathrm{A}_\alpha)_f$, where $f$ ranges over $\mathrm{A}_\alpha-\mathfrak{p}_\alpha$
\oldpage[IV-3]{6}
(for a fixed $\alpha$), since the $\mathrm{D}(f)$ form a fundamental system of neighbourhoods of $y$ in $\mathrm{U}_\alpha$, hence in $\mathrm{Y}$.
The results of the present paragraph will therefore imply that the algebraic geometry of $\sh{O}_y$-preschemes of finite presentation (and the theory of Modules of finite presentation on these preschemes) is essentially equivalent to the algebraic geometry of preschemes of finite presentation on ``sufficiently small'' open neighbourhoods of $y$.
Thus, the statement \sref{IV.8.10.5},~(xii) implies that if a morphism of finite presentation $f:\mathrm{X}\to\mathrm{Y}$ is such that $\mathrm{X}\times_\mathrm{Y}\operatorname{Spec}(\sh{O}_y)$ is proper over $\operatorname{Spec}(\sh{O}_y)$, it is necessary and sufficient that there exist an open neighbourhood $\mathrm{U}$ of $y$ in $\mathrm{Y}$ such that $f^{-1}(\mathrm{U})$ is proper over $\mathrm{U}$.

An important particular case, classical in a certain sense, is the one where $\mathrm{Y}$ is integral and where $y$ is its generic point, so that $\sh{O}_y$ is nothing other than \emph{the field} $\mathrm{R}(\mathrm{Y})=\mathrm{K}$ \emph{of rational functions on} $\mathrm{Y}$.
The results of the present paragraph then amount to interpreting the algebraic geometry over $\mathrm{K}$ in terms of algebraic geometry over ``sufficiently small'' non-empty open subsets of $\mathrm{Y}$, that is to say, intuitively, in terms of ``families'' of geometric objects indexed by the points of such an open set.
This point of view has moreover been commonly used for a long time, not only in Algebraic Geometry over algebraically closed fields, but also in the arithmetic study of varieties defined over a \emph{number field} $\mathrm{K}$ (a finite extension of $\mathbf{Q}$), considering the latter as the field of fractions of its ring of integers $\mathrm{A}$ (the ``theory of \emph{reduction modulo} $\mathfrak{p}$'' ($\mathfrak{p}$ a prime ideal of $\mathrm{A}$); cf.\ \sref[I]{I.3.7}).
The results of \textsection\textsection\,8 and~9 thus provide, among other things, some of the foundations of the language of ``reduction modulo $\mathfrak{p}$'' in arithmetic.

We note that in the example considered here, the morphisms $\mathrm{S}_\mu\to\mathrm{S}_\lambda$ (for $\lambda\leqslant\mu$) are the canonical \emph{open immersions} $\mathrm{U}_\mu\hookrightarrow\mathrm{U}_\lambda$, and \emph{a fortiori} are \emph{flat} morphisms (but not faithfully flat in general), which explains the interest of the statements that appeal to such a restriction.

  \item
  Suppose that the $\mathrm{A}_\lambda$ are \emph{fields}, so that $\mathrm{A}=\varinjlim\mathrm{A}_\lambda$ is also a field.
This case arises generally when one starts from geometric data over an arbitrary field $\mathrm{K}$, which one considers as an extension of a field $k$ (for example the prime subfield of $\mathrm{K}$).
There is then interest in considering $\mathrm{K}$ as an inductive limit of its sub-extensions that are \emph{of finite type over $k$}, which allows us to reduce many questions to the case where $\mathrm{K}$ is a finite type extension of $k$.
Using also the method sketched in \emph{a)}, one can then reduce generally to the case of a base ring $\mathrm{A}$ which is an \emph{integral algebra of finite type over $k$}.

We note that in this example, the morphisms $\mathrm{S}_\mu\to\mathrm{S}_\lambda$ are \emph{faithfully flat}.

  \item
  Suppose that we are interested in properties, local on $\mathrm{Y}$, of preschemes of finite presentation over an arbitrary prescheme $\mathrm{Y}$, which we can from now on suppose to be an affine scheme of ring $\mathrm{A}$.
There is then interest in considering $\mathrm{A}$ as an inductive limit of its subrings that are \emph{$\mathbf{Z}$-algebras of finite type}, which allows us to reduce many questions to the case where $\mathrm{Y}$ is the spectrum of such an algebra.
This is the explication of the ``Kroneckerian point of view'' according to which Algebraic Geometry reduces to the
\oldpage[IV-3]{7}
Algebraic Geometry of preschemes of finite type over $\mathbf{Z}$ (which one sometimes qualifies as ``Absolute Algebraic Geometry'').
This example shows us in particular that in most questions ``relative'' to a prescheme of base $\mathrm{Y}$, one can reduce to the case where $\mathrm{Y}$ is \emph{Noetherian}.

We note that in this example, in contrast to the previous ones, the morphisms $\mathrm{S}_\mu\to\mathrm{S}_\lambda$ have in general no particular regularity property.
\end{enumerate}

In what follows, when we apply the results that follow to any one of the three particular situations that have just been described, we will not go into detail about the procedure that allows us to make these applications, being content to refer to what precedes.
\end{env}

\begin{env}[8.1.3]
\label{IV.8.1.3}
In example \emph{a)} of \sref{IV.8.1.2}, we saw that if $\mathrm{Y}$ is an integral prescheme with generic point $y$, $f:\mathrm{X}\to\mathrm{Y}$ a morphism of finite presentation, and if the generic fibre $f^{-1}(y)$ is proper over $\boldsymbol{k}(y)$, then there exists an open neighbourhood $\mathrm{U}$ of $y$ in $\mathrm{Y}$ such that $f^{-1}(\mathrm{U})$ is proper over $\mathrm{U}$; it follows \emph{a fortiori} that for every $s\in\mathrm{U}$, $f^{-1}(s)$ is proper over $\boldsymbol{k}(s)$.
It turns out that one needs a converse, ensuring that if $f^{-1}(s)$ is proper over $\boldsymbol{k}(s)$, then $f^{-1}(y)$ is proper over $\boldsymbol{k}(y)$.
For example, suppose that $\mathrm{X}$ and $\mathrm{Y}$ are algebraic preschemes over an algebraically closed field $k$ (one can take for $k$ the field $\mathbf{C}$ of complex numbers, to fix ideas); one sometimes needs to know that, for every $s\in\mathrm{Y}$, \emph{rational over $k$}, the fibre $f^{-1}(s)$ is proper over $\boldsymbol{k}(s)$, then $f^{-1}(y)$ is proper over $\boldsymbol{k}(y)$, and consequently $f^{-1}(\mathrm{U})$ is proper over $\mathrm{U}$ for a neighbourhood $\mathrm{U}$ of $y$\footnote{We note that such a statement is ultimately purely geometric in nature, in the sense that it only appeals to rational points over $k$, and not to generic points; for example, when $k=\mathbf{C}$, this statement has an obvious topological significance for the analyst, interpreting ``proper'' in the topological sense of the term, for the underlying analytic spaces formed by the points of $\mathrm{X}$ and $\mathrm{Y}$ rational over $\mathbf{C}$.}.
Now this statement will follow easily from the following: the set $\mathrm{E}$ of points $s\in\mathrm{Y}$ such that $f^{-1}(s)$ is proper over $\boldsymbol{k}(s)$ is \emph{constructible} (and therefore identical to all of $\mathrm{Y}$ if it contains the closed points of $\mathrm{Y}$, by the theorem of zeros of Hilbert \sref[I]{I.10.4.8}); this amounts again to saying that if $f^{-1}(y)$ \emph{is not proper} over $\boldsymbol{k}(y)$, then there exists an open neighbourhood $\mathrm{U}$ of $y$ such that $f^{-1}(s)$ is not proper over $\boldsymbol{k}(s)$ for every $s\in\mathrm{U}$ (cf.\ \sref{IV.9.6.1},~(iv))).
This example illustrates the interest of systematically developing \emph{criteria of constructibility}: this is what will be done in \textsection\,9.
\end{env}

\subsection{Projective limits of preschemes}
\label{subsection:IV.8.2}

\begin{env}[8.2.1]
\label{IV.8.2.1}
Let $\mathrm{S}_0$ be a ringed space, $\mathrm{L}$ a filtered preordered increasing set, $(\sh{A}_\lambda,\varphi_{\mu\lambda})$ an inductive system of $\sh{O}_{\mathrm{S}_0}$-Algebras (not necessarily commutative), having $\mathrm{L}$ for index set.
We know that, considered as an inductive system of $\sh{O}_{\mathrm{S}_0}$-Modules, $(\sh{A}_\lambda,\varphi_{\mu\lambda})$ admits an inductive limit $\sh{A}$; let us denote by $\varphi_\lambda:\sh{A}_\lambda\to\sh{A}$ the canonical homomorphism of $\sh{O}_{\mathrm{S}_0}$-Modules.
Let $m_\lambda:\sh{A}_\lambda\otimes\sh{A}_\lambda\to\sh{A}_\lambda$ be the homomorphism of $\sh{O}_{\mathrm{S}_0}$-Modules that defines multiplication in the $\sh{O}_{\mathrm{S}_0}$-Algebra $\sh{A}_\lambda$; the hypothesis on the $\varphi_{\mu\lambda}$ implies that the $m_\lambda$ form an inductive system of homomorphisms,
\oldpage[IV-3]{8}
and since the functor $\varinjlim$ commutes with tensor products, $m=\varinjlim m_\lambda$ is a homomorphism $\sh{A}\otimes\sh{A}\to\sh{A}$ of $\sh{O}_{\mathrm{S}_0}$-Modules; by passing to the limit on the commutative diagrams expressing the associativity of $m_\lambda$ and the existence of the identity section in $\sh{A}_\lambda$, one sees that $m$ defines on $\sh{A}$ a structure of $\sh{O}_{\mathrm{S}_0}$-Algebra and that $\varphi_\lambda$ is a homomorphism of $\sh{O}_{\mathrm{S}_0}$-Algebras for every $\lambda\in\mathrm{L}$.
Furthermore, $\sh{A}$ is the inductive limit of the system $(\sh{A}_\lambda,\varphi_{\mu\lambda})$ \emph{in the category of $\sh{O}_{\mathrm{S}_0}$-Algebras}, in other words, for every $\sh{O}_{\mathrm{S}_0}$-Algebra $\sh{B}$, the canonical map
\begin{equation}
  \operatorname{Hom}_{\sh{O}_{\mathrm{S}_0}\text{-}\mathrm{Alg.}}(\sh{A},\sh{B})\to\varprojlim\operatorname{Hom}_{\sh{O}_{\mathrm{S}_0}\text{-}\mathrm{Alg.}}(\sh{A}_\lambda,\sh{B})
  \tag{8.2.1.1}
\label{IV.8.2.1.1}
\end{equation}
which to every homomorphism $f:\sh{A}\to\sh{B}$ of $\sh{O}_{\mathrm{S}_0}$-Algebras associates the family $(f\circ\varphi_\lambda)$, is \emph{bijective}.
Indeed, we already know that it is injective and identifies $\operatorname{Hom}_{\sh{O}_{\mathrm{S}_0}\text{-}\mathrm{Alg.}}(\sh{A},\sh{B})$ with a subset of $\varprojlim\operatorname{Hom}_{\sh{O}_{\mathrm{S}_0}\text{-}\mathrm{Mod.}}(\sh{A}_\lambda,\sh{B})$; it all comes down to seeing that if $(f_\lambda)$ is an inductive system of $\sh{O}_{\mathrm{S}_0}$-Algebra homomorphisms, $f_\lambda:\sh{A}_\lambda\to\sh{B}$, its inductive limit $f:\sh{A}\to\sh{B}$, which is by definition a homomorphism of $\sh{O}_{\mathrm{S}_0}$-\emph{Modules}, is also a homomorphism of $\sh{O}_{\mathrm{S}_0}$-\emph{Algebras}; but this results from passage to the inductive limit in the commutative diagram of $\sh{O}_{\mathrm{S}_0}$-Modules
\[
\xymatrix{
  \sh{A}_\lambda\otimes\sh{A}_\lambda \ar[r]^{f_\lambda\otimes f_\lambda} \ar[d]_{m_\lambda} & \sh{B}\otimes\sh{B} \ar[d] \\
  \sh{A}_\lambda \ar[r]_{f_\lambda} & \sh{B}
}
\]
and from the fact that the functor $\varinjlim$ commutes with tensor products.
We note finally that if the $\sh{A}_\lambda$ are commutative $\sh{O}_{\mathrm{S}_0}$-Algebras, the same holds for $\sh{A}$.
\end{env}

\begin{env}[8.2.2]
\label{IV.8.2.2}
Suppose now that $\mathrm{S}_0$ is a \emph{prescheme}, and that the $\sh{A}_\lambda$ are (commutative) \emph{quasi-coherent} $\sh{O}_{\mathrm{S}_0}$-Algebras; we know then that $\sh{A}=\varinjlim\sh{A}_\lambda$ is a quasi-coherent $\sh{O}_{\mathrm{S}_0}$-Algebra \sref[I]{I.4.1.1}.
Let us denote by $\mathrm{S}_\lambda$ (resp.\ $\mathrm{S}$) the spectrum of the $\sh{O}_{\mathrm{S}_0}$-Algebra $\sh{A}_\lambda$ (resp.\ $\sh{A}$) \sref[II]{II.1.3.1}, and let $u_{\lambda\mu}:\mathrm{S}_\mu\to\mathrm{S}_\lambda$ (for $\lambda\leqslant\mu$) and $u_\lambda:\mathrm{S}\to\mathrm{S}_\lambda$ be the $\mathrm{S}_0$-morphisms corresponding to the homomorphisms $\varphi_{\mu\lambda}$ and $\varphi_\lambda$ respectively \sref[II]{II.1.2.7}; it is clear that $(\mathrm{S}_\lambda,u_{\lambda\mu})$ is a \emph{projective system} in the category of $\mathrm{S}_0$-preschemes.
We note that the $u_{\lambda\mu}$ and $u_\lambda$ are \emph{affine} morphisms \sref[II]{II.1.6.2}, hence \emph{quasi-compact} and \emph{separated}.
\end{env}

\begin{proposition}[8.2.3]
\label{IV.8.2.3}
With the notations of \sref{IV.8.2.2}, the morphisms $u_\lambda:\mathrm{S}\to\mathrm{S}_\lambda$ make $\mathrm{S}$ a projective limit of the projective system $(\mathrm{S}_\lambda,u_{\lambda\mu})$ in the category of preschemes.
Furthermore, if $h:\mathrm{S}_0\to\mathrm{T}$ is a morphism, making every $\mathrm{S}_0$-prescheme a $\mathrm{T}$-prescheme, $\mathrm{S}$ is also a projective limit of the system $(\mathrm{S}_\lambda,u_{\lambda\mu})$ in the category of $\mathrm{T}$-preschemes.
\end{proposition}

\begin{proof}
Let us first prove the second assertion of the statement in the case $\mathrm{T}=\mathrm{S}_0$.
\oldpage[IV-3]{9}
It all comes down to proving that if $\mathrm{X}$ is any $\mathrm{S}_0$-prescheme, the canonical map
\begin{equation}
  \operatorname{Hom}_{\mathrm{S}_0}(\mathrm{X},\mathrm{S})\to\varprojlim\operatorname{Hom}_{\mathrm{S}_0}(\mathrm{X},\mathrm{S}_\lambda)
  \tag{8.2.3.1}
\label{IV.8.2.3.1}
\end{equation}
which to every $\mathrm{S}_0$-morphism $v:\mathrm{X}\to\mathrm{S}$ associates the family $(u_\lambda\circ v)$, is \emph{bijective}.
Now, if $g:\mathrm{X}\to\mathrm{S}_0$ is the structural morphism and if we set $\sh{B}=g_*(\sh{O}_\mathrm{X})$, which is an $\sh{O}_{\mathrm{S}_0}$-Algebra, the map \sref{IV.8.2.3.1} identifies canonically with \sref{IV.8.2.1.1} \sref[II]{II.1.2.7}, and the conclusion follows from what we have seen in \sref{IV.8.2.1}.

The other assertions of \sref{IV.8.2.3} are consequences of the following general lemma:
\end{proof}

\begin{lemma}[8.2.4]
\label{IV.8.2.4}
Let $\mathbf{C}$ be a category, $\mathrm{T}$ an object of $\mathbf{C}$, $\mathbf{C}_\mathrm{T}$ the subcategory of objects of $\mathbf{C}$ over $\mathrm{T}$.
Let $(\mathrm{S}_\lambda,u_{\lambda\mu})$ be a projective system in $\mathbf{C}_\mathrm{T}$; then every projective limit of this system in $\mathbf{C}_\mathrm{T}$ is also a projective limit in $\mathbf{C}$, and conversely.
\end{lemma}

\begin{proof}
Let $f_\lambda:\mathrm{S}_\lambda\to\mathrm{T}$ be the structural morphism.
Suppose that $\mathrm{S}$ is a projective limit of $(\mathrm{S}_\lambda,u_{\lambda\mu})$ in $\mathbf{C}$, and let us denote by $u_\lambda:\mathrm{S}\to\mathrm{S}_\lambda$ the canonical corresponding morphisms.
Consider a projective system of $\mathrm{T}$-morphisms (in $\mathbf{C}$) $w_\lambda:\mathrm{Y}\to\mathrm{S}_\lambda$, where $\mathrm{Y}\in\mathbf{C}_\mathrm{T}$.
There exists by hypothesis a unique morphism (in $\mathbf{C}$) $w:\mathrm{Y}\to\mathrm{S}$ such that $w_\lambda=u_\lambda\circ w$ for every $\lambda$.
The hypothesis that the $u_\lambda$ are $\mathrm{T}$-morphisms implies that the morphisms $f_\lambda\circ u_\lambda:\mathrm{S}\to\mathrm{T}$ are all the same, and this composite $f$ makes $\mathrm{S}$ a $\mathrm{T}$-object.
If $g:\mathrm{Y}\to\mathrm{T}$ is the structural morphism of $\mathrm{Y}$, one has $f\circ w=f_\lambda\circ u_\lambda\circ w=f_\lambda\circ w_\lambda=g$ for every $\lambda$, which proves that $w$ is a $\mathrm{T}$-morphism.
Conversely, suppose (with the same notations) that $\mathrm{S}$ is a projective limit of $(\mathrm{S}_\lambda,u_{\lambda\mu})$ in $\mathbf{C}_\mathrm{T}$, and consider now a projective system of morphisms (in $\mathbf{C}$) $w_\lambda:\mathrm{Y}\to\mathrm{S}_\lambda$.
The composites $f_\lambda\circ w_\lambda:\mathrm{Y}\to\mathrm{T}$ are then all the same: indeed, for any two indices $\lambda$, $\mu$, there is an index $\nu$ such that $\lambda\leqslant\nu$ and $\mu\leqslant\nu$, whence $f_\nu=f_\lambda\circ u_{\lambda\nu}=f_\mu\circ u_{\mu\nu}$; since $w_\lambda=u_{\lambda\nu}\circ w_\nu$ and $w_\mu=u_{\mu\nu}\circ w_\nu$, one has $f_\lambda\circ w_\lambda=f_\lambda\circ u_{\lambda\nu}\circ w_\nu=f_\nu\circ w_\nu$ and $f_\mu\circ w_\mu=f_\mu\circ u_{\mu\nu}\circ w_\nu=f_\nu\circ w_\nu$.
If $g:\mathrm{Y}\to\mathrm{T}$ is the unique morphism thus defined, $g$ makes $\mathrm{Y}$ a $\mathrm{T}$-object, and the $w_\lambda$ are $\mathrm{T}$-morphisms; they hence have a projective limit $w:\mathrm{Y}\to\mathrm{S}$ which is a $\mathrm{T}$-morphism, and \emph{a fortiori} a morphism of $\mathbf{C}$; moreover, the first part of the argument shows that any projective limit $w'$ (in $\mathbf{C}$) of the projective system $(w_\lambda)$ is also a $\mathrm{T}$-morphism, and as $w':\mathrm{Y}\to\mathrm{S}'$ in $\mathbf{C}_\mathrm{T}$ gives $p\circ w'=w$ and $p_\lambda\circ w_\lambda=p_\lambda\circ w'_\lambda$, the definition of $\mathrm{S}'_\lambda$ as a fibre product $\mathrm{S}_\lambda\times_\mathrm{T}\mathrm{T}'$ gives $u'_\lambda\circ w'=w'_\lambda$, and it is immediate that $w'$ is the unique $\mathrm{T}'$-morphism satisfying these relations, whence the lemma.
\end{proof}

\begin{proposition}[8.2.5]
\label{IV.8.2.5}
Under the conditions of \sref{IV.8.2.2}, let $\alpha$ be an element of $\mathrm{L}$, $\mathrm{X}_\alpha$ an $\mathrm{S}_\alpha$-prescheme.
For every $\lambda\geqslant\alpha$, set $\mathrm{X}_\lambda=\mathrm{X}_\alpha\times_{\mathrm{S}_\alpha}\mathrm{S}_\lambda$, and for $\alpha\leqslant\lambda\leqslant\mu$, set $v_{\lambda\mu}=1_{\mathrm{X}_\alpha}\times u_{\lambda\mu}$, so that $(\mathrm{X}_\lambda,v_{\lambda\mu})$ is a projective system of $\mathrm{X}_\alpha$-preschemes, whose index set is formed of the $\lambda\geqslant\alpha$ in $\mathrm{L}$.
Set likewise $\mathrm{X}=\mathrm{X}_\alpha\times_{\mathrm{S}_\alpha}\mathrm{S}$ and $v_\lambda=1_{\mathrm{X}_\alpha}\times u_\lambda$.
Then the $\mathrm{X}_\alpha$-morphisms $v_\lambda:\mathrm{X}\to\mathrm{X}_\lambda$ make $\mathrm{X}$ a projective limit of the projective system $(\mathrm{X}_\lambda,v_{\lambda\mu})$ in the category of $\mathrm{X}_\alpha$-preschemes, or in the category of all preschemes.
\end{proposition}

\begin{proof}
This will also follow from the following general lemma:
\end{proof}

\begin{lemma}[8.2.6]
\label{IV.8.2.6}
Let $\mathbf{C}$ be a category in which fibre products exist, $q:\mathrm{T}'\to\mathrm{T}$ a morphism of $\mathbf{C}$, $\mathbf{C}_\mathrm{T}$ (resp.\ $\mathbf{C}_{\mathrm{T}'}$) the category of objects of $\mathbf{C}$ over
\oldpage[IV-3]{10}
$\mathrm{T}$ (resp.\ $\mathrm{T}'$).
Let $(\mathrm{S}_\lambda,u_{\lambda\mu})$ be a projective system (not necessarily filtered) in $\mathbf{C}_\mathrm{T}$, and set $\mathrm{S}'_\lambda=\mathrm{S}_\lambda\times_\mathrm{T}\mathrm{T}'$, $u'_{\lambda\mu}=u_{\lambda\mu}\times 1_{\mathrm{T}'}$, so that $(\mathrm{S}'_\lambda,u'_{\lambda\mu})$ is a projective system in $\mathbf{C}_{\mathrm{T}'}$.
Then, if $(\mathrm{S},u_\lambda)$ is a projective limit of $(\mathrm{S}_\lambda,u_{\lambda\mu})$ in $\mathbf{C}_\mathrm{T}$, $(\mathrm{S}\times_\mathrm{T}\mathrm{T}',u_\lambda\times 1_{\mathrm{T}'})$ is a projective limit of $(\mathrm{S}'_\lambda,u'_{\lambda\mu})$ in $\mathbf{C}_{\mathrm{T}'}$.
\end{lemma}

\begin{proof}
By hypothesis, for every $\lambda$, one has a commutative diagram
\[
\xymatrix{
  \mathrm{S}' \ar[r]^{u'_\lambda} \ar[d]_p & \mathrm{S}'_\lambda \ar[r]^{h'_\lambda} & \mathrm{T}' \ar[d]^q \\
  \mathrm{S} \ar[r]_{u_\lambda} & \mathrm{S}_\lambda \ar[r]_{h_\lambda} & \mathrm{T}
}
\]
where we have set $\mathrm{S}'=\mathrm{S}\times_\mathrm{T}\mathrm{T}'$, $u'_\lambda=u_\lambda\times 1_{\mathrm{T}'}$, $h'_\lambda=h_\lambda\times 1_{\mathrm{T}'}$.
Let $\mathrm{Y}$ be a $\mathrm{T}'$-object, $g':\mathrm{Y}\to\mathrm{T}'$, and consider a projective system of $\mathrm{T}'$-morphisms $w'_\lambda:\mathrm{Y}\to\mathrm{S}'_\lambda$.
Then $\mathrm{Y}$ is a $\mathrm{T}$-object for the morphism $g=q\circ g'$, and the $w_\lambda=p_\lambda\circ w'_\lambda$ are $\mathrm{T}$-morphisms since $h_\lambda\circ w_\lambda=h_\lambda\circ p_\lambda\circ w'_\lambda=q\circ h'_\lambda\circ w'_\lambda=q\circ g'$ by hypothesis.
Furthermore, one verifies immediately that $(w_\lambda)$ is a projective system.
There thus exists by hypothesis a unique $\mathrm{T}$-morphism $w:\mathrm{Y}\to\mathrm{S}$ such that $u_\lambda\circ w=w_\lambda$ for every $\lambda$.
By the definition of the fibre product, there is a unique $\mathrm{T}'$-morphism $w':\mathrm{Y}\to\mathrm{S}'$ such that $p\circ w'=w$.
Then $u'_\lambda\circ w'=u_\lambda\circ w'=w_\lambda=p_\lambda\circ w'_\lambda$, and since $h'_\lambda\circ u'_\lambda\circ w'=g'=h'_\lambda\circ w'_\lambda$, the definition of $\mathrm{S}'_\lambda$ as the fibre product $\mathrm{S}_\lambda\times_\mathrm{T}\mathrm{T}'$ gives $u'_\lambda\circ w'=w'_\lambda$, and it is immediate that $w'$ is the unique $\mathrm{T}'$-morphism satisfying these relations, whence the lemma.
\end{proof}

\begin{remark}[8.2.7]
\label{IV.8.2.7}
Given a ringed space $\mathrm{S}$, the inductive limits with respect to a preordered set $\mathrm{L}$ (not necessarily filtered) exist in the category of $\sh{O}_\mathrm{S}$-Algebras, commutative or not, since the inductive limit exists by \sref{IV.8.2.1} and, on the other hand, for any two homomorphisms of $\sh{O}_\mathrm{S}$-Algebras $\sh{A}\to\sh{B}$, $\sh{A}\to\sh{C}$, the tensor product $\sh{B}\otimes_{\sh{A}}\sh{C}$ is the corresponding ``amalgamated sum'' in this category.

When $\mathrm{S}$ is a prescheme, we say that the tensor product $\sh{B}\otimes_{\sh{A}}\sh{C}$ is a quasi-coherent $\sh{O}_\mathrm{S}$-Algebra when this is so for $\sh{A}$, $\sh{B}$, and $\sh{C}$ \sref[I]{I.1.3.13}; from this one concludes that, in the category of quasi-coherent $\sh{O}_\mathrm{S}$-Algebras, inductive limits for any preordered index set always exist.
This allows us to generalise the definition of a projective limit of preschemes and propositions~\sref{IV.8.2.3} and \sref{IV.8.2.5} to the case where the preordered set $\mathrm{L}$ is not necessarily filtered.
\end{remark}

\begin{env}[8.2.8]
\label{IV.8.2.8}
The notations being those of \sref{IV.8.2.2}, set $u_{\lambda\mu}=(\psi_{\lambda\mu},\theta_{\lambda\mu})$ and $u_\lambda=(\psi_\lambda,\theta_\lambda)$; $(\mathrm{S}_\lambda,\psi_{\lambda\mu})$ is thus a projective system of topological spaces and $(\psi_\lambda)$ a projective system of continuous maps $\mathrm{S}\to\mathrm{S}_\lambda$.
\end{env}

\begin{proposition}[8.2.9]
\label{IV.8.2.9}
With the notations of \sref{IV.8.2.8}, the projective limit of the projective system $(\psi_\lambda)$ of continuous maps is a homeomorphism of the underlying space of $\mathrm{S}$ onto the projective limit of the projective system $(\mathrm{S}_\lambda,\psi_{\lambda\mu})$ of topological spaces.
\end{proposition}

\begin{proof}
\oldpage[IV-3]{11}
Let $\mathrm{T}$ be the topological limit space of the system $(\mathrm{S}_\lambda,\psi_{\lambda\mu})$ and set $\psi=\varprojlim\psi_\lambda:\mathrm{S}\to\mathrm{T}$.
We can restrict to the case where $\mathrm{S}_0=\mathrm{S}_\alpha$ for some $\alpha\in\mathrm{L}$, and $\lambda\geqslant\alpha$.

\medskip
\noindent
(i) Let us show first that the topology of $\mathrm{S}$ is the inverse image by $\psi$ of the topology of $\mathrm{T}$; in other words, if $\pi_\lambda:\mathrm{T}\to\mathrm{S}_\lambda$ is the canonical map, it amounts to showing that every open subset of $\mathrm{S}$ is a union of open subsets of the form $\psi^{-1}(\pi_\lambda^{-1}(\mathrm{U}_\lambda))$, where $\mathrm{U}_\lambda$ is open in $\mathrm{S}_\lambda$.
Now, every open subset of $\mathrm{S}$ is by definition a union of open subsets obtained in the following way: one considers an open affine $\mathrm{U}_0$ of $\mathrm{S}_0$, of ring $\mathrm{A}_0$, so that $\psi_\alpha^{-1}(\mathrm{U}_0)$ is the open affine subset of $\mathrm{S}$ of ring $\mathrm{A}=\Gamma(\mathrm{U}_0,\sh{A})$, then one takes an element $f\in\mathrm{A}$ and considers $\psi_\alpha^{-1}(\mathrm{U}_0)$, i.e., the open set $\mathrm{D}(f)$.
These are the open subsets $\mathrm{D}(f)$ which form a base of the topology of $\mathrm{S}$ \sref[II]{II.1.3.1}.
Now, if one sets $\mathrm{A}_\lambda=\Gamma(\mathrm{U}_0,\sh{A}_\lambda)$, one has $\mathrm{A}=\varinjlim\mathrm{A}_\lambda$ \sref[I]{I.1.3.9}, hence there exists an index $\lambda$ such that $f$ is the canonical image of an element $f_\lambda\in\mathrm{A}_\lambda$; one has then $\mathrm{D}(f)=\psi_\lambda^{-1}(\mathrm{D}(f_\lambda))$ \sref[I]{I.1.2.2}, and since $\psi_\lambda=\pi_\lambda\circ\psi$, our assertion is proved.

\medskip
\noindent
(ii) Let us now show that $\psi$ is \emph{bijective}, which will complete the proof.
Since $\mathrm{S}$ is a Kolmogoroff space, it already follows from (i) that $\psi$ is injective, and it remains to show that $\psi$ is surjective.
We can evidently replace $\mathrm{T}$ by an open subset $\pi_\alpha^{-1}(\mathrm{U}_0)$, where $\mathrm{U}_0$ is an open affine in $\mathrm{S}_\alpha=\mathrm{S}_0$, hence we can restrict to the case where the $\mathrm{S}_\lambda$ and $\mathrm{S}$ are affine, in other words $\sh{A}_\lambda$ is the sheaf of algebras associated to an $\mathrm{A}_0$-algebra $\mathrm{A}_\lambda$, and $\sh{A}$ the sheaf of algebras associated to $\mathrm{A}=\varinjlim\mathrm{A}_\lambda$; we will still denote $\varphi_{\mu\lambda}:\mathrm{A}_\lambda\to\mathrm{A}_\mu$ and $\varphi_\lambda:\mathrm{A}_\lambda\to\mathrm{A}$ the canonical homomorphisms.
By definition, an element of $\mathrm{T}$ is a family $(\mathfrak{p}_\lambda)_{\lambda\in\mathrm{L}}$, where $\mathfrak{p}_\lambda$ is a prime ideal of $\mathrm{A}_\lambda$ and where $\mathfrak{p}_\lambda=\varphi_{\mu\lambda}^{-1}(\mathfrak{p}_\mu)$ for $\lambda\leqslant\mu$.
We know (\sref[0]{0.5.13.3} and \sref{IV.5.13.1}) that there exists a prime ideal $\mathfrak{p}$ of $\mathrm{A}$ such that $\mathfrak{p}_\lambda=\varphi_\lambda^{-1}(\mathfrak{p})$ for every $\lambda\in\mathrm{L}$, which completes the proof.
\end{proof}

In particular, we have thus proved the

\begin{corollary}[8.2.10]
\label{IV.8.2.10}
Let $(\mathrm{A}_\lambda)_{\lambda\in\mathrm{L}}$ be a filtered inductive system of rings, and let $\varphi_\lambda:\mathrm{A}_\lambda\to\mathrm{A}$ be the canonical homomorphisms.
The canonical map $\mathfrak{p}\mapsto(\varphi_\lambda^{-1}(\mathfrak{p}))$ is a homeomorphism of $\operatorname{Spec}(\mathrm{A})$ onto the topological space $\varprojlim\operatorname{Spec}(\mathrm{A}_\lambda)$.
\end{corollary}

\begin{corollary}[8.2.11]
\label{IV.8.2.11}
With the notations of \sref{IV.8.2.8}, for every quasi-compact open subset $\mathrm{U}$ of $\mathrm{S}$, there exist an index $\lambda$ and a quasi-compact open subset $\mathrm{U}_\lambda$ of $\mathrm{S}_\lambda$ such that $\mathrm{U}=\psi_\lambda^{-1}(\mathrm{U}_\lambda)$.
\end{corollary}

\begin{proof}
This follows from the fact that, by definition of the projective limit topology, the $\psi_\lambda^{-1}(\mathrm{U}_\lambda)$ ($\mathrm{U}_\lambda$ a quasi-compact open subset of $\mathrm{S}_\lambda$) form a base of the topology of $\mathrm{S}$, and from the fact that the index set $\mathrm{L}$ is filtered.
\end{proof}

\begin{corollary}[8.2.12]
\label{IV.8.2.12}
With the notations of \sref{IV.8.2.8}, the inductive limit of the inductive system of homomorphisms $\theta_\lambda^*:\psi_\lambda^*(\sh{O}_{\mathrm{S}_\lambda})\to\sh{O}_\mathrm{S}$ of sheaves of rings on $\mathrm{S}$ is an isomorphism
\begin{equation}
  \varinjlim_\lambda\psi_\lambda^*(\sh{O}_{\mathrm{S}_\lambda})\xrightarrow{\sim}\sh{O}_\mathrm{S}.
  \label{IV.8.2.12.1}
\tag{8.2.12.1}
\end{equation}
\end{corollary}

\begin{proof}
We can evidently suppose the $\mathrm{S}_\lambda$ affine; with the notations of the proof of \sref{IV.8.2.9}, it all comes down to seeing that the inductive limit of the inductive system of canonical maps $(\mathrm{A}_\lambda)_{\mathfrak{p}_\lambda}\to\mathrm{A}_\mathfrak{p}$ is an isomorphism, which is nothing other than \sref[0]{0.5.13.3},~(ii).
\end{proof}

\begin{proposition}[8.2.13]
\label{IV.8.2.13}
Suppose that the morphisms $u_{\lambda\mu}:\mathrm{S}_\mu\to\mathrm{S}_\lambda$ are open immersions, so that $\mathrm{S}_\mu$ identifies with an open subset of $\mathrm{S}_\lambda$ for $\lambda\leqslant\mu$.
Then, for every $\alpha\in\mathrm{L}$, the underlying space of the prescheme $\mathrm{S}$ identifies with the intersection of the $\mathrm{S}_\lambda$ for $\lambda\geqslant\alpha$, and the structural sheaf $\sh{O}_\mathrm{S}$ of the sheaf induced \textnormal{(}$\mathrm{G}$, II, $1.5$\textnormal{)} by $\sh{O}_{\mathrm{S}_\alpha}$ on this intersection; more generally, for every $\sh{O}_{\mathrm{S}_\alpha}$-Module $\sh{F}_\alpha$, $u_\alpha^*(\sh{F}_\alpha)$ identifies with the $\sh{O}_\mathrm{S}$-Module induced by $\sh{F}_\alpha$ on $\mathrm{S}$.
\end{proposition}

\begin{proof}
\oldpage[IV-3]{12}
The first assertion follows from \sref{IV.8.2.9}, in view of the definition of a projective limit of topological spaces; furthermore, all the $\psi_\lambda^*(\sh{O}_{\mathrm{S}_\lambda})$ are equal to the sheaf induced by $\sh{O}_{\mathrm{S}_\alpha}$ on $\mathrm{S}$ by definition \sref[0]{0.1.3.7.1}, and with the notations of the proof of \sref{IV.8.2.9}, the homomorphisms $(\mathrm{A}_\lambda)_{\mathfrak{p}_\lambda}\to(\mathrm{A}_\mu)_{\mathfrak{p}_\mu}$ are bijective for a system $(\mathfrak{p}_\lambda)$ of ideals of primes corresponding to a single point of $\mathrm{S}$; the assertion relative to $\sh{O}_\mathrm{S}$ is then deduced from \sref{IV.8.2.12}.
The last assertion then follows from the definition of $u_\alpha^*(\sh{F}_\alpha)$ \sref[0]{0.1.4.3.1}.
\end{proof}

\begin{remark}[8.2.14]
\label{IV.8.2.14}
The results of \sref{IV.8.2.9} and \sref{IV.8.2.12} show that $\mathrm{S}$ is a projective limit of the projective system $(\mathrm{S}_\lambda,u_{\lambda\mu})$ in the category of \emph{all} ringed spaces (or of all ringed spaces with a local ringed space structure).
Indeed, let $\mathrm{Y}$ be a ringed space, and consider a projective system of morphisms of ringed spaces $w_\lambda:\mathrm{Y}\to\mathrm{S}_\lambda$.
If one writes $w_\lambda=(\rho_\lambda,\omega_\lambda)$, the $\rho_\lambda$ form a projective system of continuous maps and, by \sref{IV.8.2.9}, their projective limit identifies with a continuous map $\rho:\mathrm{Y}\to\mathrm{S}$ such that $\rho_\lambda=\psi_\lambda\circ\rho$. On the other hand, the homomorphisms
\[
 \omega_\lambda^\sharp:\rho_\lambda^*(\sh{O}_{\mathrm{S}_\lambda})\longrightarrow\sh{O}_\mathrm{Y}
\]
form an inductive system of homomorphisms of sheaves of rings. Since
\[
 \rho_\lambda^*(\sh{O}_{\mathrm{S}_\lambda})
 =\rho^*\bigl(\psi_\lambda^*(\sh{O}_{\mathrm{S}_\lambda})\bigr)
\]
and the functor $\rho^*$ is exact, the inductive limit of the $\rho_\lambda^*(\sh{O}_{\mathrm{S}_\lambda})$ is $\rho^*(\sh{O}_\mathrm{S})$ by \sref{IV.8.2.12}. There is therefore a unique homomorphism
\[
 \omega^\sharp:\rho^*(\sh{O}_\mathrm{S})\longrightarrow\sh{O}_\mathrm{Y}
\]
such that $\omega_\lambda^\sharp=\omega^\sharp\circ\rho^*(\theta_\lambda^*)$ for every $\lambda$, which proves the assertion.
\end{remark}

\subsection{Constructible subsets of a projective limit of preschemes}
\label{subsection:IV.8.3}

\begin{env}[8.3.1]
\label{IV.8.3.1}
In what follows throughout this paragraph, we suppose the conditions of \sref{IV.8.2.2} are satisfied, and we keep the notations.
\end{env}

\begin{theorem}[8.3.2]
\label{IV.8.3.2}
For every $\lambda$, let $\mathrm{E}_\lambda$, $\mathrm{F}_\lambda$ be two subsets of $\mathrm{S}_\lambda$.
Set
\begin{equation}
  \mathrm{E}=\bigcap_\lambda u_\lambda^{-1}(\mathrm{E}_\lambda),\qquad \mathrm{F}=\bigcup_\lambda u_\lambda^{-1}(\mathrm{F}_\lambda).
  \tag{8.3.2.1}
\label{IV.8.3.2.1}
\end{equation}
Suppose the following conditions are satisfied:
\begin{enumerate}[label=(\roman*)]
  \item For every $\lambda$, $\mathrm{E}_\lambda$ is pro-constructible and $\mathrm{F}_\lambda$ is ind-constructible \sref{IV.1.9.4}.
  \item For $\lambda\leqslant\mu$, we have
  \[
    \mathrm{E}_\mu\subset u_{\lambda\mu}^{-1}(\mathrm{E}_\lambda),\qquad \mathrm{F}_\mu\supset u_{\lambda\mu}^{-1}(\mathrm{F}_\lambda).
  \]
  \item There exists $\alpha\in\mathrm{L}$ such that $\mathrm{S}_\alpha$ is quasi-compact (which implies that $\mathrm{S}_\lambda$ is quasi-compact for $\lambda\geqslant\alpha$).
\end{enumerate}
Then the following properties are equivalent:
\begin{enumerate}[label=\emph{\alph*)}]
  \item $\mathrm{E}\subset\mathrm{F}$.
  \item There exists $\lambda\geqslant\alpha$ such that $u_\lambda^{-1}(\mathrm{E}_\lambda)\subset u_\lambda^{-1}(\mathrm{F}_\lambda)$ (and then we also have $u_\mu^{-1}(\mathrm{E}_\mu)\subset u_\mu^{-1}(\mathrm{F}_\mu)$ for
\oldpage[IV-3]{13}
$\mu\geqslant\lambda$).
  \item There exists $\lambda\geqslant\alpha$ such that $\mathrm{E}_\lambda\subset\mathrm{F}_\lambda$ (and then we also have $\mathrm{E}_\mu\subset\mathrm{F}_\mu$ for $\mu\geqslant\lambda$).
\end{enumerate}
The parenthetical remarks in \emph{b)} and \emph{c)} follow from (ii).
Set
\[
  \mathrm{G}_\lambda=\mathrm{E}_\lambda\cap(\mathrm{S}_\lambda-\mathrm{F}_\lambda),\qquad \mathrm{G}=\mathrm{E}\cap(\mathrm{S}-\mathrm{F}).
\]
Then $\mathrm{G}_\lambda$ is a pro-constructible subset of $\mathrm{S}_\lambda$ \sref{IV.1.9.5},~(i), and by virtue of \sref{IV.8.3.2.1} we have
\[
  \mathrm{G}_\mu\subset u_{\lambda\mu}^{-1}(\mathrm{G}_\lambda)\qquad\text{for }\lambda\leqslant\mu
\]
\[
  \mathrm{G}=\bigcap_\lambda u_\lambda^{-1}(\mathrm{G}_\lambda).
\]
We are thus reduced to proving the particular case of \sref{IV.8.3.2} corresponding to $\mathrm{F}_\lambda=\varnothing$ for every $\lambda$:
\end{theorem}

\begin{corollary}[8.3.3]
\label{IV.8.3.3}
For every $\lambda$, let $\mathrm{E}_\lambda$ be a pro-constructible subset of $\mathrm{S}_\lambda$ such that $\mathrm{E}_\mu\subset u_{\lambda\mu}^{-1}(\mathrm{E}_\lambda)$.
Suppose that there exists $\alpha\in\mathrm{L}$ such that $\mathrm{S}_\alpha$ is quasi-compact.
Then the following conditions are equivalent:
\begin{enumerate}[label=\emph{\alph*)}]
  \item $\mathrm{E}=\bigcap_\lambda u_\lambda^{-1}(\mathrm{E}_\lambda)=\varnothing$.
  \item There exists $\lambda$ such that $u_\lambda^{-1}(\mathrm{E}_\lambda)=\varnothing$ (and then $u_\mu^{-1}(\mathrm{E}_\mu)=\varnothing$ for $\mu\geqslant\lambda$).
  \item There exists $\lambda$ such that $\mathrm{E}_\lambda=\varnothing$ (and then $\mathrm{E}_\mu=\varnothing$ for $\mu\geqslant\lambda$).
\end{enumerate}
\end{corollary}

\begin{proof}
It is clear that \emph{c)} implies \emph{a)}.
Let us prove that \emph{a)} implies \emph{b)}: since $\mathrm{S}_\lambda$ is quasi-compact, $\mathrm{E}_\lambda$, being pro-constructible, is also quasi-compact, and it is the same for $u_\lambda^{-1}(\mathrm{E}_\lambda)$ \sref{IV.1.9.5},~(vi); then the filtered decreasing family of pro-constructible subsets $u_\lambda^{-1}(\mathrm{E}_\lambda)$ has empty intersection, hence \sref{IV.1.9.9} one of them is empty.

Finally, let us show that \emph{b)} implies \emph{c)}.
Since $\mathrm{S}_\alpha$ is quasi-compact and $\mathrm{L}$ is filtered, we can replace $\mathrm{S}_\alpha$ by an open affine, hence suppose \sref{IV.8.2.2} that $\mathrm{S}$ and the $\mathrm{S}_\lambda$ for $\lambda\geqslant\alpha$ are affine; one has then \sref{IV.1.9.2.1}, for $\lambda\geqslant\alpha$,
\[
  u_\lambda(\mathrm{S})=\bigcap_{\mu\geqslant\lambda}u_{\lambda\mu}(\mathrm{S}_\mu)
\]
whence
\[
  \mathrm{E}_\lambda\cap u_\lambda(\mathrm{S})=\bigcap_{\mu\geqslant\lambda}(\mathrm{E}_\lambda\cap u_{\lambda\mu}(\mathrm{S}_\mu)).
\]
Now, since $u_\lambda$ and the $u_{\lambda\mu}$ are quasi-compact, $u_\lambda(\mathrm{S})$ and $u_{\lambda\mu}(\mathrm{S}_\mu)$ are pro-constructible in $\mathrm{S}_\lambda$ \sref{IV.1.9.5},~(vii), hence the sets $\mathrm{E}_\lambda\cap u_{\lambda\mu}(\mathrm{S}_\mu)$ for $\mu\geqslant\lambda$ form a filtered decreasing family of pro-constructible subsets of $\mathrm{S}_\lambda$.
Since $\mathrm{S}_\lambda$ is quasi-compact, the hypothesis \emph{b)} implies that the intersection of this family is empty, hence \sref{IV.1.9.9} one of the sets $\mathrm{E}_\lambda\cap u_{\lambda\mu}(\mathrm{S}_\mu)$ is empty, hence $\mathrm{E}_\mu\subset u_{\lambda\mu}^{-1}(\mathrm{E}_\lambda)$ is empty.
\end{proof}

\begin{corollary}[8.3.4]
\label{IV.8.3.4}
For every $\lambda$, let $\mathrm{F}_\lambda$ be an ind-constructible subset of $\mathrm{S}_\lambda$ such that for $\lambda\leqslant\mu$ one has $u_{\lambda\mu}^{-1}(\mathrm{F}_\lambda)\subset\mathrm{F}_\mu$.
Suppose that there exists $\alpha\in\mathrm{L}$ such that $\mathrm{S}_\alpha$ is quasi-compact.
Then the following conditions are equivalent:
\begin{enumerate}[label=\emph{\alph*)}]
  \item The set $\mathrm{F}=\bigcup_\lambda u_\lambda^{-1}(\mathrm{F}_\lambda)$ is equal to $\mathrm{S}$.
  \item There exists $\lambda$ such that $u_\lambda^{-1}(\mathrm{F}_\lambda)=\mathrm{S}$ (and then $u_\mu^{-1}(\mathrm{F}_\mu)=\mathrm{S}$ for $\mu\geqslant\lambda$).
\oldpage[IV-3]{14}
  \item There exists $\lambda$ such that $\mathrm{F}_\lambda=\mathrm{S}_\lambda$ (and then $\mathrm{F}_\mu=\mathrm{S}_\mu$ for $\mu\geqslant\lambda$).
\end{enumerate}
This follows immediately from \sref{IV.8.3.3} by passing to complements.
\end{corollary}

\begin{corollary}[8.3.5]
\label{IV.8.3.5}
For every $\lambda$, let $\mathrm{E}_\lambda$, $\mathrm{F}_\lambda$ be two constructible subsets of $\mathrm{S}_\lambda$ such that, for $\mu\geqslant\lambda$, one has $\mathrm{E}_\mu\subset u_{\lambda\mu}^{-1}(\mathrm{E}_\lambda)$ and $\mathrm{F}_\mu\supset u_{\lambda\mu}^{-1}(\mathrm{F}_\lambda)$.
Suppose that there exists $\alpha$ such that $\mathrm{S}_\alpha$ is quasi-compact.
Then, for $\mathrm{E}\subset\mathrm{F}$ \textnormal{(}resp.\ $\mathrm{E}=\mathrm{F}$\textnormal{)}, it is necessary and sufficient that there exist $\lambda$ such that $\mathrm{E}_\lambda\subset\mathrm{F}_\lambda$ \textnormal{(}resp.\ $\mathrm{E}_\lambda=\mathrm{F}_\lambda$\textnormal{)}, in which case one also has $\mathrm{E}_\mu\subset\mathrm{F}_\mu$ \textnormal{(}resp.\ $\mathrm{E}_\mu=\mathrm{F}_\mu$\textnormal{)} for $\mu\geqslant\lambda$.

This is a particular case of \sref{IV.8.3.2}.
\end{corollary}

\begin{corollary}[8.3.6]
\label{IV.8.3.6}
Suppose that there exists $\alpha$ such that $\mathrm{S}_\alpha$ is quasi-compact.
For $\mathrm{S}=\varnothing$, it is necessary and sufficient that there exist $\lambda$ such that $\mathrm{S}_\lambda=\varnothing$.
\end{corollary}

\begin{corollary}[8.3.7]
\label{IV.8.3.7}
For every $\lambda$, we have
\begin{equation}
  u_\lambda(\mathrm{S})=\bigcap_{\mu\geqslant\lambda}u_{\lambda\mu}(\mathrm{S}_\mu).
  \tag{8.3.7.1}
\label{IV.8.3.7.1}
\end{equation}
\end{corollary}

\begin{proof}
It is clear that the first member of \sref{IV.8.3.7.1} is contained in the second.
Let $s$ be a point of $\mathrm{S}_\lambda$ and set $\mathrm{X}_\lambda=\operatorname{Spec}(\boldsymbol{k}(s))$; consider the projective system $(\mathrm{X}_\mu,v_{\mu\nu})$ where $\mathrm{X}_\mu=\mathrm{X}_\lambda\times_{\mathrm{S}_\lambda}\mathrm{S}_\mu$ and $v_{\mu\nu}=1\times u_{\mu\nu}$ for $\lambda\leqslant\mu\leqslant\nu$; its projective limit is $\mathrm{X}=\mathrm{X}_\lambda\times_{\mathrm{S}_\lambda}\mathrm{S}$ \sref{IV.8.2.5}.
If $s\in\bigcap_{\mu\geqslant\lambda}u_{\lambda\mu}(\mathrm{S}_\mu)$, this implies that $\mathrm{X}_\mu\neq\varnothing$ for every $\mu\geqslant\lambda$ \sref[I]{I.3.4.8}; it follows then from \sref{IV.8.3.6} that $\mathrm{X}\neq\varnothing$, hence $s\in u_\lambda(\mathrm{S})$ by \sref[I]{I.3.4.8}.
\end{proof}

\begin{proposition}[8.3.8]
\label{IV.8.3.8}
\begin{enumerate}[label=(\roman*)]
  \item For the morphism $u_\lambda:\mathrm{S}\to\mathrm{S}_\lambda$ to be dominant (resp.\ surjective), it is necessary and sufficient that for $\mu\geqslant\lambda$ the morphism $u_{\lambda\mu}:\mathrm{S}_\mu\to\mathrm{S}_\lambda$ be dominant (resp.\ surjective).
  \item If, for some index $\lambda$, the morphisms $u_{\lambda\mu}:\mathrm{S}_\mu\to\mathrm{S}_\lambda$ are flat (resp.\ faithfully flat) for every $\mu\geqslant\lambda$, then the morphism $u_\lambda:\mathrm{S}\to\mathrm{S}_\lambda$ is flat (resp.\ faithfully flat).
  \item Suppose that the morphisms $u_\mu:\mathrm{S}\to\mathrm{S}_\mu$ are surjective for $\mu\geqslant\lambda$.
For $u_\lambda$ to be an open morphism (resp.\ universally open), it is necessary and sufficient that, for every $\mu\geqslant\lambda$, $u_{\lambda\mu}$ be an open morphism (resp.\ universally open).
\end{enumerate}
\end{proposition}

\begin{proof}
\begin{enumerate}[label=(\roman*)]
  \item Since $u_\lambda(\mathrm{S})\subset u_{\lambda\mu}(\mathrm{S}_\mu)$ for every $\mu\geqslant\lambda$, the necessity of the conditions is trivial, and it follows immediately from \sref{IV.8.3.7.1} that if the $u_{\lambda\mu}$ are surjective, then so is $u_\lambda$.
Suppose now that the $u_{\lambda\mu}$ are dominant for $\mu\geqslant\lambda$, and consider in $\mathrm{S}_\lambda$ an open quasi-compact non-empty subset $\mathrm{U}_\lambda$; then the $\mathrm{U}_\mu=u_{\lambda\mu}^{-1}(\mathrm{U}_\lambda)$ for $\mu\geqslant\lambda$ form a projective system whose projective limit is $\mathrm{U}=u_\lambda^{-1}(\mathrm{U}_\lambda)$ \sref{IV.8.2.5}.
By hypothesis the $\mathrm{U}_\mu$ are all non-empty, hence it is the same for $\mathrm{U}$ by \sref{IV.8.3.6}, which proves that $u_\lambda$ is dominant.

  \item By virtue of (i), it suffices to consider the case where the $u_{\lambda\mu}$ are flat; one can then restrict to the case where $\mathrm{S}_\lambda$ is affine, hence also the $\mathrm{S}_\mu$ and $\mathrm{S}$, and the assertion follows from \sref{IV.2.1.2} and \sref[0]{0.6.2.3}.

  \item By virtue of \sref{IV.8.2.5} and of \sref[I]{I.3.5.2}, it suffices to treat the case of open morphisms.
Since $u_\lambda=u_{\lambda\mu}\circ u_\mu$ and $u_\mu$ is surjective, we know that if $u_\lambda$ is open then so is $u_{\lambda\mu}$ (Bourbaki, \emph{Top.\ g\'{e}n.}, chap.~I\textsuperscript{er}, 3\textsuperscript{e}~ed., \textsection\,5, n\textsuperscript{o}\,1, prop.~1).
\oldpage[IV-3]{15}
Conversely, to show that $u_\lambda$ is open when all the $u_{\lambda\mu}$ are open for $\mu\geqslant\lambda$, it suffices to see that for every quasi-compact open subset $\mathrm{V}$ of $\mathrm{S}$, $u_\lambda(\mathrm{V})$ is open in $\mathrm{S}_\lambda$; but there exists then $\mu\geqslant\lambda$ such that $\mathrm{V}=u_\mu^{-1}(\mathrm{V}_\mu)$, where $\mathrm{V}_\mu$ is open in $\mathrm{S}_\mu$ \sref{IV.8.2.11}; since $u_\mu$ is surjective, one has $\mathrm{V}_\mu=u_\mu(\mathrm{V})$ and $u_\lambda(\mathrm{V})=u_{\lambda\mu}(u_\mu(\mathrm{V}))$ is therefore open by hypothesis.
\end{enumerate}

One notes that it may happen that all the $u_{\lambda\mu}$ are open without $u_\lambda$ being open, when the $u_\mu$ are not surjective.
Indeed, one has an example by considering an integral ring $\mathrm{A}$ which is not a field, and its field of fractions $\mathrm{K}$, which is the inductive limit of the $\mathrm{A}_f$, where $f$ ranges over $\mathrm{A}-\{0\}$; setting $\mathrm{S}_1=\operatorname{Spec}(\mathrm{A})$, $\mathrm{S}_f=\operatorname{Spec}(\mathrm{A}_f)$, and $\mathrm{S}=\varprojlim\mathrm{S}_f=\operatorname{Spec}(\mathrm{K})$, the morphism $\mathrm{S}\to\mathrm{S}_1$ is not open, even though the morphisms $\mathrm{S}_f\to\mathrm{S}_1$ are open.
\end{proof}

\begin{env}[8.3.9]
\label{IV.8.3.9}
For every prescheme $\mathrm{X}$, we will denote as usual by $\mathfrak{P}(\mathrm{X})$ the set of all subsets of the underlying set of $\mathrm{X}$, by $\mathfrak{C}(\mathrm{X})$ (resp.\ $\mathfrak{Dc}(\mathrm{X})$, $\mathfrak{Fc}(\mathrm{X})$, $\mathfrak{LFc}(\mathrm{X})$) the set of constructible (resp.\ constructible and open, constructible and closed, and locally closed) subsets of $\mathrm{X}$.
It is clear that $(\mathfrak{P}(\mathrm{S}_\lambda),u_{\lambda\mu}^{-1})$ is an \emph{inductive system of sets} and that the maps $u_\lambda^{-1}:\mathfrak{P}(\mathrm{S}_\lambda)\to\mathfrak{P}(\mathrm{S})$ form an inductive system of maps, whence, by passing to the inductive limit, a canonical map
\begin{equation}
  u_\mathfrak{P}:\varinjlim\mathfrak{P}(\mathrm{S}_\lambda)\to\mathfrak{P}(\mathrm{S}).
  \tag{8.3.9.1}
\label{IV.8.3.9.1}
\end{equation}

Furthermore, it follows from \sref{IV.1.8.2} that $u_{\lambda\mu}^{-1}$ sends $\mathfrak{C}(\mathrm{S}_\lambda)$ (resp.\ $\mathfrak{Dc}(\mathrm{S}_\lambda)$, $\mathfrak{Fc}(\mathrm{S}_\lambda)$, $\mathfrak{LFc}(\mathrm{S}_\lambda)$) into $\mathfrak{C}(\mathrm{S}_\mu)$ (resp.\ $\mathfrak{Dc}(\mathrm{S}_\mu)$, $\mathfrak{Fc}(\mathrm{S}_\mu)$, $\mathfrak{LFc}(\mathrm{S}_\mu)$) and that $u_\lambda^{-1}$ sends $\mathfrak{C}(\mathrm{S}_\lambda)$ (resp.\ $\mathfrak{Dc}(\mathrm{S}_\lambda)$, $\mathfrak{Fc}(\mathrm{S}_\lambda)$, $\mathfrak{LFc}(\mathrm{S}_\lambda)$) into $\mathfrak{C}(\mathrm{S})$ (resp.\ $\mathfrak{Dc}(\mathrm{S})$, $\mathfrak{Fc}(\mathrm{S})$, $\mathfrak{LFc}(\mathrm{S})$).
We thus have, by restriction of \sref{IV.8.3.9.1}, canonical maps
\begin{align}
  &\varinjlim\mathfrak{C}(\mathrm{S}_\lambda)\to\mathfrak{C}(\mathrm{S}) \tag{8.3.9.2}\label{IV.8.3.9.2}\\
  &\varinjlim\mathfrak{Dc}(\mathrm{S}_\lambda)\to\mathfrak{Dc}(\mathrm{S}) \tag{8.3.9.3}\label{IV.8.3.9.3}\\
  &\varinjlim\mathfrak{Fc}(\mathrm{S}_\lambda)\to\mathfrak{Fc}(\mathrm{S}) \tag{8.3.9.4}\label{IV.8.3.9.4}\\
  &\varinjlim\mathfrak{LFc}(\mathrm{S}_\lambda)\to\mathfrak{LFc}(\mathrm{S}). \tag{8.3.9.5}
\label{IV.8.3.9.5}
\end{align}
\end{env}

\begin{env}[8.3.10]
\label{IV.8.3.10}
Let $g_\alpha:\mathrm{X}_\alpha\to\mathrm{S}_\alpha$ be a morphism; with the notations of \sref{IV.8.2.5} one has, on the one hand, a canonical map $v_\mathfrak{P}:\varinjlim\mathfrak{P}(\mathrm{S}_\lambda)\to\mathfrak{P}(\mathrm{X})$; on the other hand, from the morphisms of projection $g_\lambda:\mathrm{X}_\lambda\to\mathrm{S}_\lambda$ for every $\lambda\geqslant\alpha$ and a morphism of projection $g:\mathrm{X}\to\mathrm{S}$.
It is clear that the $g_\lambda^{-1}:\mathfrak{P}(\mathrm{S}_\lambda)\to\mathfrak{P}(\mathrm{X}_\lambda)$ form an inductive system of maps, and that the diagrams
\[
\xymatrix{
  \mathfrak{P}(\mathrm{S}_\lambda) \ar[r]^{u_\lambda^{-1}} \ar[d]_{g_\lambda^{-1}} & \mathfrak{P}(\mathrm{S}) \ar[d]^{g^{-1}} \\
  \mathfrak{P}(\mathrm{X}_\lambda) \ar[r]_{v_\lambda^{-1}} & \mathfrak{P}(\mathrm{X})
}
\]
\oldpage[IV-3]{16}
are commutative; one deduces from this by passing to the inductive limit a commutative diagram
\[
\xymatrix{
  \varinjlim\mathfrak{P}(\mathrm{S}_\lambda) \ar[r]^{u_\mathfrak{P}} \ar[d] & \mathfrak{P}(\mathrm{S}) \ar[d]^{g^{-1}} \\
  \varinjlim\mathfrak{P}(\mathrm{X}_\lambda) \ar[r]_{v_\mathfrak{P}} & \mathfrak{P}(\mathrm{X})
}
\]
and it follows from \sref{IV.1.8.2} that one has analogous diagrams replacing $\mathfrak{P}$ by $\mathfrak{C}$, $\mathfrak{Dc}$, $\mathfrak{Fc}$, or $\mathfrak{LFc}$.

\begin{equation}
  \label{IV.8.3.10.1}
\tag{8.3.10.1}
\end{equation}
\end{env}

It follows from \sref{IV.8.3.5} that under the hypothesis that for some $\alpha\in\mathrm{L}$, $\mathrm{S}_\alpha$ is \emph{quasi-compact}, the canonical map \sref{IV.8.3.9.2} is \emph{injective}.
Furthermore:

\begin{theorem}[8.3.11]
\label{IV.8.3.11}
Suppose that there exists $\alpha$ such that $\mathrm{S}_\alpha$ is quasi-compact and quasi-separated.
Then the canonical maps \sref{IV.8.3.9.2}, \sref{IV.8.3.9.3}, \sref{IV.8.3.9.4}, and \sref{IV.8.3.9.5} are \emph{bijective}.
\end{theorem}

\begin{proof}
By virtue of the preceding remark, it remains to prove that these maps are \emph{surjective}; as every constructible subset of $\mathrm{S}$ is a finite union of sets of the form $\mathrm{U}\cap\complement\mathrm{V}$, where $\mathrm{U}$ and $\mathrm{V}$ are open and constructible, it will suffice to prove that \sref{IV.8.3.9.3} is surjective for it to be the same for \sref{IV.8.3.9.2} (and also for \sref{IV.8.3.9.4}, by passing to complements).
Now, since the morphisms $\mathrm{S}_\lambda\to\mathrm{S}_\alpha$ are affine and separated, the $\mathrm{S}_\lambda$ for $\lambda\geqslant\alpha$ and $\mathrm{S}$ are quasi-compact and quasi-separated, and one knows that the constructible open subsets of such a prescheme are none other than the quasi-compact open subsets \sref{IV.1.8.1}.
The conclusion then follows from \sref{IV.8.2.11} except for the map \sref{IV.8.3.9.5}.
To prove that this last map is surjective, consider a locally closed and constructible subset $\mathrm{Z}$ of $\mathrm{S}$; $\mathrm{Z}$ is therefore quasi-compact ($\mathbf{0}_{\mathrm{III}}$,~9.1.4).
As every point $x\in\mathrm{Z}$ admits by hypothesis an open quasi-compact neighbourhood $\mathrm{V}_x$ in $\mathrm{S}$ such that $\mathrm{Z}\cap\mathrm{V}_x$ is closed in $\mathrm{V}_x$; we can cover $\mathrm{Z}$ by a finite number of the $\mathrm{V}_x$, hence there is a quasi-compact open subset $\mathrm{U}$ containing $\mathrm{Z}$ such that $\mathrm{Z}$ is closed in $\mathrm{U}$; in other words, $\mathrm{Z}$ is also constructible in $\mathrm{U}$ ($\mathbf{0}_{\mathrm{III}}$,~9.1.8).
We know \sref{IV.8.2.11} that there exist $\lambda$ and a quasi-compact open subset $\mathrm{U}_\lambda$ of $\mathrm{S}_\lambda$ such that $\mathrm{U}=u_\lambda^{-1}(\mathrm{U}_\lambda)$.
Applying to $\mathrm{U}$ (which is a projective limit of $\mathrm{U}_\lambda=u_{\lambda\mu}^{-1}(\mathrm{U}_\lambda)$) the map \sref{IV.8.3.9.4}, one sees that there exist $\mu\geqslant\lambda$ and a closed constructible subset $\mathrm{Z}_\mu$ in $\mathrm{U}_\mu$ such that $\mathrm{Z}=u_\mu^{-1}(\mathrm{Z}_\mu)$.
But since the canonical immersion $\mathrm{U}_\mu\to\mathrm{S}_\mu$ is quasi-compact by hypothesis \sref{IV.1.2.7}, she is of finite presentation \sref{IV.1.6.2},~(i), and $\mathrm{Z}_\mu$ is also a constructible subset \emph{of} $\mathrm{S}_\mu$ by virtue of \sref{IV.1.8.4} and \sref{IV.1.8.1}; since $\mathrm{Z}_\mu$ is evidently locally closed in $\mathrm{S}_\mu$, this completes the proof.
\end{proof}

\begin{corollary}[8.3.12]
\label{IV.8.3.12}
Suppose that there exists $\alpha$ such that $\mathrm{S}_\alpha$ is quasi-compact, and let, for every $\lambda$, $\mathrm{Z}_\lambda$ be a constructible subset of $\mathrm{S}_\lambda$ such that $\mathrm{Z}_\mu=u_{\lambda\mu}^{-1}(\mathrm{Z}_\lambda)$ for $\mu\geqslant\lambda$.
If $\mathrm{Z}=u_\lambda^{-1}(\mathrm{Z}_\lambda)$,
\oldpage[IV-3]{17}
then, for $\mathrm{Z}$ to be open (resp.\ closed, resp.\ locally closed) in $\mathrm{S}$, it is necessary and sufficient that there exist $\lambda\geqslant\alpha$ such that $\mathrm{Z}_\lambda$ be so in $\mathrm{S}_\lambda$.
\end{corollary}

\begin{proof}
Cover $\mathrm{S}_\alpha$ by a finite number of open affine subsets $\mathrm{U}_\alpha^{(i)}$; then the $\mathrm{U}^{(i)}=u_\alpha^{-1}(\mathrm{U}_\alpha^{(i)})$ form an open covering of $\mathrm{S}$, and for $\mathrm{Z}$ to be open (resp.\ closed, resp.\ locally closed) in $\mathrm{S}$, it is necessary and sufficient that each of the $\mathrm{Z}\cap\mathrm{U}^{(i)}$ be so in $\mathrm{U}^{(i)}$.
Since $\mathrm{L}$ is filtered and each $\mathrm{Z}\cap\mathrm{U}^{(i)}$ is constructible in $\mathrm{U}^{(i)}$ ($\mathbf{0}_{\mathrm{III}}$,~9.1.8), one can restrict to proving the corollary when $\mathrm{S}_\alpha$ is affine, hence quasi-compact and quasi-separated; but then it follows from \sref{IV.8.3.11}.
\end{proof}

\begin{proposition}[8.3.13]
\label{IV.8.3.13}
Suppose that the morphisms $u_{\lambda\mu}:\mathrm{S}_\mu\to\mathrm{S}_\lambda$ are flat and that there exists $\alpha$ such that $\mathrm{S}_\alpha$ is quasi-compact.
For every $\lambda$, let $\mathrm{Z}'_\lambda$, $\mathrm{Z}''_\lambda$ be two pro-constructible subsets of $\mathrm{S}_\lambda$, such that $\mathrm{Z}'_\mu=u_{\lambda\mu}^{-1}(\mathrm{Z}'_\lambda)$ and $\mathrm{Z}''_\mu=u_{\lambda\mu}^{-1}(\mathrm{Z}''_\lambda)$ for $\mu\geqslant\lambda$; suppose further that $\overline{\mathrm{Z}'_\alpha}$ is constructible in $\mathrm{S}_\alpha$.
Set $\mathrm{Z}'=u_\lambda^{-1}(\mathrm{Z}'_\lambda)$, $\mathrm{Z}''=u_\lambda^{-1}(\mathrm{Z}''_\lambda)$; for $\mathrm{Z}''\subset\overline{\mathrm{Z}'}$, it is necessary and sufficient that there exist $\lambda\geqslant\alpha$ such that $\mathrm{Z}''_\lambda\subset\overline{\mathrm{Z}'_\lambda}$.
\end{proposition}

\begin{proof}
Indeed, one knows that $u_\lambda:\mathrm{S}\to\mathrm{S}_\lambda$ is also a flat morphism for every $\lambda$ \sref{IV.8.3.8},~(i); as $\mathrm{Z}'_\lambda$ is pro-constructible, the closure of $\mathrm{Z}'_\mu$ for $\mu\geqslant\lambda$ (resp.\ of $\mathrm{Z}'$) in $\mathrm{S}_\mu$ (resp.\ $\mathrm{S}$) is equal to $u_{\lambda\mu}^{-1}(\overline{\mathrm{Z}'_\lambda})$ (resp.\ $u_\lambda^{-1}(\overline{\mathrm{Z}'_\lambda})$) \sref{IV.2.3.10}.
Since the $u_{\lambda\mu}^{-1}(\overline{\mathrm{Z}'_\lambda})$ and $u_\lambda^{-1}(\overline{\mathrm{Z}'_\lambda})$ are constructible \sref{IV.1.8.2}, the conclusion follows from \sref{IV.8.3.2}.
\end{proof}

\subsection{Criteria of irreducibility and connectedness for projective limits of preschemes}
\label{subsection:IV.8.4}

\begin{proposition}[8.4.1]
\label{IV.8.4.1}
Suppose that there exists an index $\alpha$ such that $\mathrm{S}_\alpha$ is quasi-compact.
\begin{enumerate}[label=(\roman*)]
  \item If $\mathrm{S}$ is not irreducible and if in addition the underlying space of $\mathrm{S}$ is Noetherian and $\mathrm{S}_\alpha$ is quasi-separated, there exists $\lambda\geqslant\alpha$ such that, for $\mu\geqslant\lambda$, $\mathrm{S}_\mu$ is not irreducible.
  \item If $\mathrm{S}$ is not connected, there exists $\lambda\geqslant\alpha$ such that, for $\mu\geqslant\lambda$, $\mathrm{S}_\mu$ is not connected.
\end{enumerate}
\end{proposition}

\begin{proof}
Suppose that $\mathrm{S}$ is a union of two distinct closed subsets $\mathrm{S}'$, $\mathrm{S}''$ of $\mathrm{S}$ (resp.\ disjoint non-empty closed subsets).
In case (i), $\mathrm{S}'$ and $\mathrm{S}''$ are constructible since the space $\mathrm{S}$ is Noetherian.
By virtue of \sref{IV.8.3.11}, there exist $\lambda\geqslant\alpha$ and two constructible closed subsets $\mathrm{S}'_\lambda$, $\mathrm{S}''_\lambda$ of $\mathrm{S}_\lambda$ such that $\mathrm{S}'=u_\lambda^{-1}(\mathrm{S}'_\lambda)$, $\mathrm{S}''=u_\lambda^{-1}(\mathrm{S}''_\lambda)$; since $\mathrm{S}=\mathrm{S}'\cup\mathrm{S}''$, it follows from \sref{IV.8.3.11} that one can suppose $\mathrm{S}_\lambda=\mathrm{S}'_\lambda\cup\mathrm{S}''_\lambda$; since $\mathrm{S}'_\lambda$ and $\mathrm{S}''_\lambda$ are distinct from $\mathrm{S}_\lambda$, this proves that $\mathrm{S}_\lambda$ is not irreducible.

In case (ii), $\mathrm{S}'$ and $\mathrm{S}''$ are open and quasi-compact, hence, by virtue of \sref{IV.8.2.11}, there exist $\lambda\geqslant\alpha$ and two open quasi-compact subsets $\mathrm{S}'_\lambda$, $\mathrm{S}''_\lambda$ of $\mathrm{S}_\lambda$ such that $\mathrm{S}'=u_\lambda^{-1}(\mathrm{S}'_\lambda)$, $\mathrm{S}''=u_\lambda^{-1}(\mathrm{S}''_\lambda)$.
Moreover, since $\mathrm{S}'$ and $\mathrm{S}''$ are open and closed in $\mathrm{S}$, they are both pro-constructible and ind-constructible \sref{IV.1.9.6}, hence constructible \sref{IV.1.9.5}, and it follows from \sref{IV.8.3.5} that one can suppose $\lambda$ taken such that $\mathrm{S}_\lambda=\mathrm{S}'_\lambda\cup\mathrm{S}''_\lambda$ and $\mathrm{S}'_\lambda\cap\mathrm{S}''_\lambda=\varnothing$, which shows that $\mathrm{S}_\lambda$ is not connected.
\end{proof}

\begin{proposition}[8.4.2]
\label{IV.8.4.2}
Suppose that the underlying space of $\mathrm{S}$ is Noetherian and that one of the two following conditions is satisfied:
\begin{enumerate}[label=\emph{\alph*)}]
  \item For $\lambda\leqslant\mu$, $u_{\lambda\mu}:\mathrm{S}_\mu\to\mathrm{S}_\lambda$ is dominant, and there exists $\alpha$ such that $\mathrm{S}_\alpha$ is quasi-compact.
\oldpage[IV-3]{18}
  \item There exists $\alpha$ such that the underlying space of $\mathrm{S}_\alpha$ is Noetherian, and for $\mu\geqslant\lambda$, $u_{\lambda\mu}$ is a homeomorphism of $\mathrm{S}_\mu$ onto a subspace of $\mathrm{S}_\lambda$.
\end{enumerate}
Under these conditions, there exists $\lambda$ such that, for every $\mu\geqslant\lambda$:
\begin{enumerate}[label=(\roman*)]
  \item For every irreducible component $\mathrm{Y}_i$ of $\mathrm{S}$ ($1\leqslant i\leqslant m$), $\overline{u_\mu(\mathrm{Y}_i)}$ is an irreducible component of $\mathrm{S}_\mu$, and the map $\mathrm{Y}_i\mapsto\overline{u_\mu(\mathrm{Y}_i)}$ is a bijection of the set of irreducible components of $\mathrm{S}$ onto the set of irreducible components of $\mathrm{S}_\mu$.
  \item For every connected component $\mathrm{C}_j$ of $\mathrm{S}$ ($1\leqslant j\leqslant n$), $\overline{u_\mu(\mathrm{C}_j)}$ is a connected component of $\mathrm{S}_\mu$, and the map $\mathrm{C}_j\mapsto\overline{u_\mu(\mathrm{C}_j)}$ is a bijection of the set of connected components of $\mathrm{S}$ onto the set of connected components of $\mathrm{S}_\mu$.
\end{enumerate}
\end{proposition}

\begin{proof}
Let us first establish the

\begin{lemma}[8.4.2.1]
\label{IV.8.4.2.1}
Under condition \emph{a)} or \emph{b)} of \sref{IV.8.4.2}, there exists $\lambda$ such that, for $\mu\geqslant\lambda$, $u_\mu:\mathrm{S}\to\mathrm{S}_\mu$ is dominant.
\end{lemma}

In case \emph{a)}, this has already been proved without supposing the space $\mathrm{S}$ Noetherian \sref{IV.8.3.8},~(i).
In case \emph{b)}, set $\mathrm{Z}_\alpha=\overline{u_\alpha(\mathrm{S})}$; as a closed subset of the Noetherian space $\mathrm{S}_\alpha$, $\mathrm{Z}_\alpha$ is constructible, and since $u_\alpha^{-1}(\mathrm{Z}_\alpha)=\mathrm{S}$, it follows from \sref{IV.8.3.5} that there exists $\lambda\geqslant\alpha$ such that $u_{\lambda\mu}^{-1}(\mathrm{Z}_\lambda)=\mathrm{S}_\mu$.
But since $u_{\alpha\mu}$ is a homeomorphism and $u_\alpha(\mathrm{S})$ is dense in $\mathrm{S}_\alpha$, it follows from $\mathrm{S}\to\mathrm{Z}_\mu=\mathrm{S}_\mu$.

This lemma being proved, one can suppose that for every $\lambda$, $u_\lambda$ is a dominant morphism.
\begin{enumerate}[label=(\roman*)]
  \item Each of the $\mathrm{S}_\lambda$ is a union of the $\overline{u_\lambda(\mathrm{Y}_i)}$, which are irreducible.
On the other hand, if $\mathrm{U}_i$ is the open subset of $\mathrm{S}$ complementary to the union of the $\mathrm{Y}_j$ of index $j\neq i$ ($1\leqslant i\leqslant m$), the $\mathrm{U}_i$ are pairwise disjoint and $\mathrm{Y}_i=\overline{\mathrm{U}_i}$ \sref[0]{0.2.1.6}.
Since the underlying space of $\mathrm{S}$ is Noetherian, the $\mathrm{U}_i$ are quasi-compact, hence there exist an index $\lambda$ and open subsets $\mathrm{U}_{i\lambda}$ of $\mathrm{S}_\lambda$ such that $\mathrm{U}_i=u_\lambda^{-1}(\mathrm{U}_{i\lambda})$ for $1\leqslant i\leqslant m$ \sref{IV.8.2.11}.
From this one concludes that if one sets $\mathrm{U}_{i\mu}=u_{\lambda\mu}^{-1}(\mathrm{U}_{i\lambda})$ for $\lambda\leqslant\mu$, the $\mathrm{U}_{i\mu}$ are pairwise disjoint, the $\mathrm{U}_{i\mu}$ are open and $u_\mu(\mathrm{U}_i)$ is dense in $\mathrm{U}_{i\mu}$ since $u_\mu$ is dominant.
By this, none of the closures $\overline{\mathrm{U}_{i\mu}}$ is contained in another and $u_\mu(\mathrm{U}_i)$ is dense in $\mathrm{U}_{i\mu}$ since $u_\mu$ is dominant; hence $\overline{\mathrm{U}_{i\mu}}=\overline{u_\mu(\mathrm{Y}_i)}$, which proves that the $\overline{u_\mu(\mathrm{Y}_i)}$ are the irreducible components of $\mathrm{S}_\mu$ \sref[0]{0.2.1.7} and completes the proof.

  \item Since the underlying space $\mathrm{S}$ is Noetherian, the $\mathrm{C}_j$ are open and closed in $\mathrm{S}$ and quasi-compact; the same reasoning as in (i) shows that there exist $\lambda$ and open subsets $\mathrm{V}_{j\lambda}$ of $\mathrm{S}_\lambda$ such that $\mathrm{C}_j=u_\lambda^{-1}(\mathrm{V}_{j\lambda})$ for $1\leqslant j\leqslant n$.
One sees also, as in (i), that if one sets $\mathrm{V}_{j\mu}=u_{\lambda\mu}^{-1}(\mathrm{V}_{j\lambda})$ for $\mu\geqslant\lambda$, the $\mathrm{V}_{j\mu}$ are pairwise disjoint, and $u_\mu(\mathrm{C}_j)$ is dense in $\mathrm{V}_{j\mu}$.
Furthermore, it follows from \sref{IV.8.3.4} that for $\mu$ sufficiently large, the union of the $\mathrm{V}_{j\mu}$ is $\mathrm{S}_\mu$, since every open subset of a Noetherian prescheme is ind-constructible \sref{IV.1.9.6}.
The $\mathrm{V}_{j\mu}$ are thus the connected components of $\mathrm{S}_\mu$, this completes the proof.
\end{enumerate}
\end{proof}

We note that if the morphisms $u_{\lambda\mu}$ are \emph{immersions}, they satisfy in particular condition \emph{b)} of \sref{IV.8.4.2}.

\begin{corollary}[8.4.3]
\label{IV.8.4.3}
Suppose one of the conditions \emph{a)}, \emph{b)} of \sref{IV.8.4.2} is satisfied, the underlying space of $\mathrm{S}$ being Noetherian; then, for $\mathrm{S}$ to be irreducible, it is necessary and sufficient that there exist $\lambda$ such that $\mathrm{S}_\mu$ be so for every $\mu\geqslant\lambda$.
\end{corollary}

\begin{proposition}[8.4.4]
\label{IV.8.4.4}
\oldpage[IV-3]{19}
Suppose that there exists $\alpha$ such that $\mathrm{S}_\alpha$ is quasi-compact and that, for $\lambda\leqslant\mu$, $u_{\lambda\mu}:\mathrm{S}_\mu\to\mathrm{S}_\lambda$ is dominant.
Then, for $\mathrm{S}$ to be connected, it is necessary and sufficient that there exist $\lambda$ such that $\mathrm{S}_\mu$ be so for every $\mu\geqslant\lambda$.
\end{proposition}

\begin{proof}
The condition is sufficient by virtue of \sref{IV.8.4.1},~(i); on the other hand, we have seen \sref{IV.8.3.8},~(i) that $u_\lambda:\mathrm{S}\to\mathrm{S}_\lambda$ is dominant for $\lambda$ sufficiently large, hence, if $\mathrm{S}$ is connected, it is the same for $\mathrm{S}_\lambda$, since $u_\lambda(\mathrm{S})$ is dense in $\mathrm{S}_\lambda$ and is connected.
\end{proof}

\begin{corollary}[8.4.5]
\label{IV.8.4.5}
Let $k$ be a field, $\mathrm{X}$ a quasi-compact $k$-prescheme.
For $\mathrm{X}$ to be geometrically connected \sref{IV.4.5.2}, it is necessary and sufficient that, for every finite and separable extension $\mathrm{K}$ of $k$, $\mathrm{X}\otimes_k\mathrm{K}$ be connected.
\end{corollary}

\begin{proof}
The condition is trivially necessary.
To see that it is sufficient, we must prove that if $\Omega$ is an algebraic closure of $k$, $\mathrm{X}\otimes_k\Omega$ is connected \sref{IV.4.5.1}.
Now, $\Omega$ is a filtered inductive limit of the finite sub-extensions of $k$, and for $k'\subset k''\subset\Omega$, the morphism $\mathrm{X}\otimes_k k''\to\mathrm{X}\otimes_k k'$ is surjective.
We are thus reduced, by virtue of \sref{IV.8.4.4}, to proving that $\mathrm{X}\otimes_k k'$ is connected for every finite extension $k'$ of $k$.
But if $\mathrm{K}$ is the largest separable extension contained in $k'$, the morphism $\mathrm{X}\otimes_k k'\to\mathrm{X}\otimes_k\mathrm{K}$ is finite, surjective, and radicial, hence \sref{IV.2.4.5} a homeomorphism, and since $\mathrm{X}\otimes_k\mathrm{K}$ is connected by hypothesis, so is $\mathrm{X}\otimes_k k'$.
\end{proof}

\begin{remark}[8.4.6]
\label{IV.8.4.6}
\begin{enumerate}[label=(\roman*)]
  \item The proof of \sref{IV.8.4.2} shows that the conclusion of this proposition is valid if one supposes that the underlying space of $\mathrm{S}$ is Noetherian, that there exists $\alpha$ such that $\mathrm{S}_\alpha$ is quasi-compact, and finally that the $u_\lambda:\mathrm{S}\to\mathrm{S}_\lambda$ are dominant.
  \item On the other hand, the conclusion of \sref{IV.8.4.2} can be at fault when the $u_\lambda$ are not dominant for $\lambda$ sufficiently large, even when the $\mathrm{S}_\lambda$ and $\mathrm{S}$ are Noetherian, as shown by the following example.
Take $\mathbf{N}$ for index set, all the $\mathrm{S}_n$ equal to $\operatorname{Spec}(\mathrm{A}\times\mathrm{K})=\operatorname{Spec}(\mathrm{A})\amalg\operatorname{Spec}(\mathrm{K})$, where $\mathrm{K}$ is a field, $\mathrm{A}$ any $\mathrm{K}$-algebra, and all the morphisms $u_{n,n+1}$ equal to the same morphism corresponding to the homomorphism $(x,y)\mapsto(j(y),y)$ of $\mathrm{A}\times\mathrm{K}$ into itself, where $j:\mathrm{K}\to\mathrm{A}$ is the canonical homomorphism.
One easily verifies that the inductive limit of this system of rings is $\mathrm{K}$, the canonical homomorphism $\mathrm{A}\times\mathrm{K}\to\mathrm{K}$ being the second projection.
One thus sees that $\mathrm{S}=\operatorname{Spec}(\mathrm{K})$ is irreducible even though none of the $\mathrm{S}_n$ is connected.
\end{enumerate}
\end{remark}

\subsection{Modules of finite presentation on a projective limit of preschemes}
\label{subsection:IV.8.5}

\begin{env}[8.5.1]
\label{IV.8.5.1}
We always keep the notations of \sref{IV.8.2.2}; we will further restrict ourselves to the case where $\mathrm{S}_0$ is one of the $\mathrm{S}_\lambda$, to which one can always reduce.

\emph{Whenever, in this section, we consider a family $(\sh{F}_\lambda)$, where, for every $\lambda\in\mathrm{L}$, $\sh{F}_\lambda$ is an $\sh{O}_{\mathrm{S}_\lambda}$-Module, it will be understood that this family satisfies the condition}
\begin{equation}
  \sh{F}_\mu=u_{\lambda\mu}^*(\sh{F}_\lambda) \qquad\text{for }\lambda\leqslant\mu.
  \tag{8.5.1.1}
\label{IV.8.5.1.1}
\end{equation}
\oldpage[IV-3]{20}
We will then set
\begin{equation}
  \sh{F}=u_\lambda^*(\sh{F}_\lambda)
  \label{IV.8.5.1.2}
\tag{8.5.1.2}
\end{equation}
which is an $\sh{O}_\mathrm{S}$-Module not depending on the index $\lambda\in\mathrm{L}$, by virtue of the hypothesis \sref{IV.8.5.1.1}.

Let $(\sh{F}_\lambda)$, $(\sh{G}_\lambda)$ be two such families of $\sh{O}_{\mathrm{S}_\lambda}$-Modules.
It is clear that the maps $f_\lambda\mapsto u_{\lambda\mu}^*(f_\lambda)$ from $\operatorname{Hom}_{\sh{O}_{\mathrm{S}_\lambda}}(\sh{F}_\lambda,\sh{G}_\lambda)$ to $\operatorname{Hom}_{\sh{O}_{\mathrm{S}_\mu}}(\sh{F}_\mu,\sh{G}_\mu)$ define an \emph{inductive system} of abelian groups $(\operatorname{Hom}_{\sh{O}_{\mathrm{S}_\lambda}}(\sh{F}_\lambda,\sh{G}_\lambda),u_{\lambda\mu}^*)$, and that the maps $f_\lambda\mapsto u_\lambda^*(f_\lambda)$ form an inductive system of homomorphisms of abelian groups
\[
  u_\lambda^*:\operatorname{Hom}_{\sh{O}_{\mathrm{S}_\lambda}}(\sh{F}_\lambda,\sh{G}_\lambda)\to\operatorname{Hom}_\mathrm{S}(\sh{F},\sh{G});
\]
whence, by passing to the inductive limit, a canonical homomorphism of abelian groups
\begin{equation}
  u_{\sh{F},\sh{G}}^*:\varinjlim\operatorname{Hom}_{\sh{O}_{\mathrm{S}_\lambda}}(\sh{F}_\lambda,\sh{G}_\lambda)\to\operatorname{Hom}_\mathrm{S}(\sh{F},\sh{G}).
  \tag{8.5.1.3}
\label{IV.8.5.1.3}
\end{equation}

Note that when $\sh{F}_\lambda=\sh{O}_{\mathrm{S}_\lambda}$, the condition \sref{IV.8.5.1.1} is satisfied, and one has $\sh{F}=\sh{O}_\mathrm{S}$; the homomorphism \sref{IV.8.5.1.3} thus gives \sref[0]{0.5.1.1} a canonical homomorphism of abelian groups
\begin{equation}
  u_{\sh{G}}:\varinjlim\Gamma(\mathrm{S}_\lambda,\sh{G}_\lambda)\to\Gamma(\mathrm{S},\sh{G}).
  \tag{8.5.1.4}
\label{IV.8.5.1.4}
\end{equation}
\end{env}

\begin{theorem}[8.5.2]
\label{IV.8.5.2}
\begin{enumerate}[label=(\roman*)]
  \item Suppose $\mathrm{S}_0$ quasi-compact (resp.\ quasi-compact and quasi-separated) and that, for some $\lambda\in\mathrm{L}$, $\sh{F}_\lambda$ is quasi-coherent and of finite type (resp.\ of finite presentation) and $\sh{G}_\lambda$ quasi-coherent.
Then the homomorphism $u_{\sh{F},\sh{G}}^*$ is injective (resp.\ bijective).
  \item Suppose $\mathrm{S}_0$ quasi-compact and quasi-separated.
For every quasi-coherent $\sh{O}_\mathrm{S}$-Module $\sh{F}$ of finite presentation, there exist $\lambda\in\mathrm{L}$ and an $\sh{O}_{\mathrm{S}_\lambda}$-Module quasi-coherent of finite presentation $\sh{F}_\lambda$ such that $\sh{F}$ is isomorphic to $u_\lambda^*(\sh{F}_\lambda)$.
\end{enumerate}
\end{theorem}

\begin{proof}
\begin{enumerate}[label=(\roman*)]
  \item We can evidently restrict to the case where $\mathrm{S}_0=\mathrm{S}_\lambda$ since the morphisms $u_{0\lambda}:\mathrm{S}_\lambda\to\mathrm{S}_0$ are affine, hence quasi-compact and separated.
Consider first the case where $\mathrm{S}_0=\operatorname{Spec}(\mathrm{A}_0)$ is \emph{affine}.
Then assertion (i) is equivalent to the

\begin{lemma}[8.5.2.1]
\label{IV.8.5.2.1}
Let $\mathrm{A}_0$ be a ring, $(\mathrm{A}_\lambda)_{\lambda\in\mathrm{L}}$ an inductive system of $\mathrm{A}_0$-algebras, $\mathrm{A}=\varinjlim\mathrm{A}_\lambda$; let $\mathrm{M}_0$, $\mathrm{N}_0$ be two $\mathrm{A}_0$-modules, and set $\mathrm{M}_\lambda=\mathrm{M}_0\otimes_{\mathrm{A}_0}\mathrm{A}_\lambda$, $\mathrm{N}_\lambda=\mathrm{N}_0\otimes_{\mathrm{A}_0}\mathrm{A}_\lambda$, $\mathrm{M}=\mathrm{M}_0\otimes_{\mathrm{A}_0}\mathrm{A}=\varinjlim\mathrm{M}_\lambda$, $\mathrm{N}=\mathrm{N}_0\otimes_{\mathrm{A}_0}\mathrm{A}=\varinjlim\mathrm{N}_\lambda$.
If $\mathrm{M}_0$ is of finite type (resp.\ of finite presentation), the canonical homomorphism
\begin{equation}
  \varinjlim\operatorname{Hom}_{\mathrm{A}_\lambda}(\mathrm{M}_\lambda,\mathrm{N}_\lambda)\to\operatorname{Hom}_\mathrm{A}(\mathrm{M},\mathrm{N})
  \tag{8.5.2.2}
\label{IV.8.5.2.2}
\end{equation}
is injective (resp.\ bijective).
\end{lemma}

We know indeed (Bourbaki, \emph{Alg.}, chap.~II, 3\textsuperscript{e}~ed., \textsection\,5, n\textsuperscript{o}\,1) that one has canonical functorial isomorphisms
\[
  \operatorname{Hom}_{\mathrm{A}_\lambda}(\mathrm{M}_\lambda,\mathrm{N}_\lambda)\xrightarrow{\sim}\operatorname{Hom}_{\mathrm{A}_0}(\mathrm{M}_0,\mathrm{N}_\lambda),\qquad\operatorname{Hom}_\mathrm{A}(\mathrm{M},\mathrm{N})\xrightarrow{\sim}\operatorname{Hom}_{\mathrm{A}_0}(\mathrm{M}_0,\mathrm{N})
\]
so that the homomorphism \sref{IV.8.5.2.2} is none other, up to canonical isomorphisms, than the canonical homomorphism
\begin{equation}
  \varinjlim\operatorname{Hom}_{\mathrm{A}_0}(\mathrm{M}_0,\mathrm{N}_\lambda)\to\operatorname{Hom}_{\mathrm{A}_0}(\mathrm{M}_0,\varinjlim\mathrm{N}_\lambda),
  \tag{8.5.2.3}
\label{IV.8.5.2.3}
\end{equation}
\oldpage[IV-3]{21}
which, to every inductive system of homomorphisms of $\mathrm{A}_0$-modules $\theta_\lambda:\mathrm{M}_0\to\mathrm{N}_\lambda$, associates its inductive limit.

Now, if $\mathrm{M}_0$ is of finite type (resp.\ of finite presentation), one has an exact sequence $\mathrm{A}_0^m\to\mathrm{M}_0\to 0$ (resp.\ $\mathrm{A}_0^n\to\mathrm{A}_0^m\to\mathrm{M}_0\to 0$); since it is clear that \sref{IV.8.5.2.3} is bijective when $\mathrm{M}_0$ is of the form $\mathrm{A}_0^q$, it suffices to use the left-exactness of the functor $\mathrm{M}_0\mapsto\operatorname{Hom}_{\mathrm{A}_0}(\mathrm{M}_0,\mathrm{P})$ and the exactness of the functor $\varinjlim$ (in the category of abelian groups) to conclude.

\medskip
Let us pass to the case where $\mathrm{S}_0$ is quasi-compact, and let $(\mathrm{U}_i)$ be a finite open affine covering of $\mathrm{S}_0$; for every $\lambda$, the $\mathrm{U}_{i\lambda}=u_{0\lambda}^{-1}(\mathrm{U}_i)$ form a covering of the open affine subsets of $\mathrm{S}_\lambda$, and the $\mathrm{V}_i=u_0^{-1}(\mathrm{U}_i)$ a covering of the open affine subsets of $\mathrm{S}$.
To show that $u_{\sh{F},\sh{G}}^*$ is injective, it suffices to show that if $f_\lambda:\sh{F}_\lambda\to\sh{G}_\lambda$ is such that $f=u_\lambda^*(f_\lambda)=0$, then there exists $\mu\geqslant\lambda$ such that $u_{\lambda\mu}^*(f_\lambda)=0$.
By virtue of the lemma \sref{IV.8.5.2.1}, for each $i$ there exists $\mu_i\geqslant\lambda$ such that $f_\mu|_{\mathrm{U}_{i\mu_i}}=0$.
It suffices therefore to take $\mu$ greater than all the $\lambda_i$.

Suppose further that $\mathrm{S}_0$ is quasi-separated and $\sh{F}_0$ is a quasi-coherent $\sh{O}_{\mathrm{S}_0}$-Module of finite presentation.
Note now that $\mathrm{S}_\lambda$ is quasi-separated and $\sh{F}_\lambda$ is a quasi-coherent $\sh{O}_{\mathrm{S}_\lambda}$-Module of finite presentation \sref[0]{0.5.2.5}; since, for every pair of indices $i$, $j$, $\mathrm{U}_{ij\lambda}=\mathrm{U}_{i\lambda}\cap\mathrm{U}_{j\lambda}$ is quasi-compact and quasi-separated \sref{IV.1.2.7}, it follows from \sref{IV.8.5.2.5} that for every pair of indices $i$, $j$, $u_\lambda^*(f_\lambda^{(i)})=u_\lambda^*(f_\lambda^{(j)})$ on $\mathrm{U}_{ij\lambda}$ by definition, it follows from what we have already seen above that there exists $\lambda_{ij}$ such that $u_{\lambda\mu}^*(f_\lambda^{(i)})|_{\mathrm{U}_{ij\lambda}}=u_{\lambda\mu}^*(f_\lambda^{(j)})|_{\mathrm{U}_{ij\lambda}}$ for $\mu\geqslant\lambda_{ij}$; taking $\mu$ greater than all the $\lambda_{ij}$, one sees that $u_{\lambda\mu}^*(f_\lambda^{(i)})$ and $u_{\lambda\mu}^*(f_\lambda^{(j)})$ coincide in $\mathrm{U}_{i\mu}\cap\mathrm{U}_{j\mu}$ for every couple $(i,j)$, and therefore define an homomorphism $f_\mu:\sh{F}_\mu\to\sh{G}_\mu$ such that $f=u_\mu^*(f_\mu)$.

  \item Before passing to the proof of (ii), let us note the following corollaries of (i):
\end{enumerate}
\end{proof}

\begin{corollary}[8.5.2.4]
\label{IV.8.5.2.4}
Suppose $\mathrm{S}_0$ quasi-compact, $\sh{F}_\lambda$ quasi-coherent of finite type, $\sh{G}_\lambda$ quasi-coherent.
Let $f_\lambda:\sh{F}_\lambda\to\sh{G}_\lambda$ be a homomorphism.
For $f=u_\lambda^*(f_\lambda):\sh{F}\to\sh{G}$ to be an isomorphism, it is necessary and sufficient that there exist $\mu\geqslant\lambda$ such that $f_\mu=u_{\lambda\mu}^*(f_\lambda):\sh{F}_\mu\to\sh{G}_\mu$ be so.
\end{corollary}

\begin{proof}
We can always suppose $\mathrm{S}_\lambda=\mathrm{S}_0$, and we can further reduce to the case where $\mathrm{S}_0$ is affine, hence also quasi-separated.
Then the assertion is equivalent to the lemma \sref{IV.8.5.2.1}, for there is by hypothesis an $\sh{O}_\mathrm{S}$-homomorphism $g:\sh{G}\to\sh{F}$ such that $g\circ f=1_\sh{F}$ and $f\circ g=1_\sh{G}$.
Since $\sh{G}$ is of finite presentation, there exist by virtue of \sref{IV.8.5.2},~(i), an index $\nu\geqslant\lambda$ and a homomorphism $g_\nu:\sh{G}_\nu\to\sh{F}_\nu$ such that $g=u_\nu^*(g_\nu)$ by virtue of \sref{IV.8.5.2},~(i); one has therefore $u_\nu^*(g_\nu\circ f_\nu)=1_\sh{F}$ and $u_\nu^*(f_\nu\circ g_\nu)=1_\sh{G}$, whence the corollary.
\end{proof}

\begin{corollary}[8.5.2.5]
\label{IV.8.5.2.5}
Suppose $\mathrm{S}_0$ quasi-compact and quasi-separated.
Suppose that $\sh{F}_\lambda$, $\sh{G}_\lambda$ are $\sh{O}_{\mathrm{S}_\lambda}$-Modules quasi-coherent of finite presentation.
For $\sh{F}$ and $\sh{G}$ to be isomorphic, it is necessary and sufficient that there exist $\mu\geqslant\lambda$ such that $\sh{F}_\mu$ and $\sh{G}_\mu$ are isomorphic.
Moreover, for every isomorphism $f:\sh{F}\xrightarrow{\sim}\sh{G}$, there exist $\nu\geqslant\mu$ and an isomorphism $f_\nu:\sh{F}_\nu\xrightarrow{\sim}\sh{G}_\nu$ such that $f=u_\nu^*(f_\nu)$.
\end{corollary}

\begin{proof}
\oldpage[IV-3]{22}
This follows from \sref{IV.8.5.2.4} and from \sref{IV.8.5.2},~(i) since every homomorphism $f:\sh{F}\to\sh{G}$ is of the form $u_\mu^*(f_\mu)$ for some $\mu\geqslant\lambda$ and some homomorphism $f_\mu:\sh{F}_\mu\to\sh{G}_\mu$.
\end{proof}

\begin{proof}[Proof of \sref{IV.8.5.2},~(ii)]
Consider again first the case where $\mathrm{S}_0=\operatorname{Spec}(\mathrm{A}_0)$ is affine.
Then the assertion is equivalent to the lemma \sref{IV.5.13.7.1}.

In the general case, we start from a finite open affine covering $(\mathrm{U}_i)$ of $\mathrm{S}_0$, and for every $\lambda$, the $\mathrm{U}_{i\lambda}=u_{0\lambda}^{-1}(\mathrm{U}_i)$ form an open affine covering of $\mathrm{S}_\lambda$, and the $\mathrm{V}_i=u_0^{-1}(\mathrm{U}_i)$ an open affine covering of $\mathrm{S}$.
By the affine case, for every $i$, there exist an index $\lambda(i)$ and an $\sh{O}_{\mathrm{U}_{i,\lambda(i)}}$-Module quasi-coherent of finite presentation $\sh{F}_i$ such that $u_{\lambda(i)}^*(\sh{F}_i)=\sh{F}|_{\mathrm{V}_i}$ (with the notations of (i)).
Furthermore, since $\mathrm{L}$ is filtered, we can suppose that all the $\lambda_i$ are equal to a single $\lambda$.
Note now that $\mathrm{S}_\lambda$ is quasi-separated and $\sh{F}_\lambda$ is a quasi-coherent $\sh{O}_{\mathrm{S}_\lambda}$-Module of finite presentation \sref[0]{0.5.2.5}; since, for every pair of indices $i$, $j$, $\mathrm{U}_{ij\lambda}=\mathrm{U}_{i\lambda}\cap\mathrm{U}_{j\lambda}$ is quasi-compact and quasi-separated \sref{IV.1.2.7}, it follows from \sref{IV.8.5.2.5} that for every pair $(i,j)$, there exist an index $\lambda_{ij}\geqslant\lambda$ and an isomorphism $\theta_{ij}:\sh{F}_i|_{\mathrm{U}_{ij\lambda}}\xrightarrow{\sim}u_{\lambda,\lambda_{ij}}^*(\sh{F}_j)|_{\mathrm{U}_{ij\lambda}}$ such that $u_{\lambda_{ij}}^*(\theta_{ij})$ is the identity automorphism of $\sh{F}|_{(\mathrm{V}_i\cap\mathrm{V}_j)}$; one can further suppose all the $\lambda_{ij}$ equal to $\lambda$.
Changing the notations, one can therefore suppose that it holds for every pair $(i,j)$ an isomorphism $\theta_{ij}:\sh{F}_i|_{\mathrm{U}_{ij\lambda}}\xrightarrow{\sim}\sh{F}_j|_{\mathrm{U}_{ij\lambda}}$ such that $u_\lambda^*(\theta_{ij})$ is the identity automorphism of $\sh{F}|_{(\mathrm{V}_i\cap\mathrm{V}_j)}$.
Finally, for any three indices $i$, $j$, $k$, if one sets $\mathrm{U}_{ijk,\lambda}=\mathrm{U}_{i\lambda}\cap\mathrm{U}_{j\lambda}\cap\mathrm{U}_{k\lambda}$, $\mathrm{U}_{ijk,\lambda}$ is quasi-compact, and if $\theta'_{ij}$, $\theta'_{jk}$, and $\theta'_{ik}$ denote the restrictions of $\theta_{ij}$, $\theta_{jk}$, and $\theta_{ik}$ to $\mathrm{U}_{ijk,\lambda}$, one has $u_\lambda^*(\theta'_{ij}\circ\theta'_{jk})=u_\lambda^*(\theta'_{ik})$.
There is, hence, by virtue of (i), an index $\mu\geqslant\lambda$ such that $u_{\lambda\mu}^*(\theta'_{ik})=u_{\lambda\mu}^*(\theta'_{ij}\circ\theta'_{jk})$, hence the isomorphisms $u_{\lambda\mu}^*(\theta'_{ij}):u_{\lambda\mu}^*(\sh{F}_i)|_{\mathrm{U}_{ij\mu}}\xrightarrow{\sim}u_{\lambda\mu}^*(\sh{F}_j)|_{\mathrm{U}_{ij\mu}}$ satisfy the gluing condition, and define on $\mathrm{S}_\mu$ a quasi-coherent $\sh{O}_{\mathrm{S}_\mu}$-Module $\sh{F}_\mu$ of finite presentation \sref[0]{0.3.3.1} such that $\sh{F}$ is isomorphic to $u_\mu^*(\sh{F}_\mu)$.
\end{proof}

\begin{env}[8.5.3]
\label{IV.8.5.3}
The result of \sref{IV.8.5.2} can be expressed again by saying that if $\mathrm{S}_0$ is quasi-compact and quasi-separated, the category of $\sh{O}_\mathrm{S}$-Modules quasi-coherent of finite presentation is determined up to equivalence near the datum of the categories of $\sh{O}_{\mathrm{S}_\lambda}$-Modules quasi-coherent of finite presentation, of the functors $u_{\lambda\mu}^*$ between these categories, and of the transition isomorphisms $u_{\mu\nu}^*\circ u_{\lambda\mu}^*\xrightarrow{\sim}u_{\lambda\nu}^*$.
In a figurative way, one can say that the datum of a quasi-coherent $\sh{O}_\mathrm{S}$-Module of finite presentation $\sh{F}$ is ``functorially'' equivalent to the datum of a quasi-coherent $\sh{O}_{\mathrm{S}_\lambda}$-Module of finite presentation $\sh{F}'_\lambda$ for some sufficiently large $\lambda$ and that, if a quasi-coherent $\sh{O}_{\mathrm{S}_\mu}$-Module of finite presentation $\sh{F}''_\mu$ also has $\sh{F}$ as inverse image, then $\sh{F}'_\lambda$ and $\sh{F}''_\mu$ have the same inverse image in some suitable $\mathrm{S}_\nu$ ($\nu\geqslant\lambda$, $\nu\geqslant\mu$).

We will now interpret from this point of view various notions related to quasi-coherent $\sh{O}_\mathrm{S}$-Modules.
\end{env}

\begin{corollary}[8.5.4]
\label{IV.8.5.4}
Suppose $\mathrm{S}_0$ quasi-compact and quasi-separated; then, for every quasi-coherent $\sh{O}_{\mathrm{S}_\lambda}$-Module $\sh{G}_\lambda$, the canonical homomorphism \sref{IV.8.5.1.4} is bijective.
\end{corollary}

\begin{proof}
Indeed, it suffices to apply \sref{IV.8.5.2},~(i) to the case where $\sh{F}_\lambda=\sh{O}_{\mathrm{S}_\lambda}$, which is of finite presentation.
\end{proof}

\begin{proposition}[8.5.5]
\label{IV.8.5.5}
\oldpage[IV-3]{23}
Suppose $\mathrm{S}_0$ quasi-compact, and suppose that $\sh{F}_\lambda$ is an $\sh{O}_{\mathrm{S}_\lambda}$-Module quasi-coherent of finite presentation.
For $\sh{F}$ to be locally free (resp.\ locally free of rank $n$), it is necessary and sufficient that there exist $\mu\geqslant\lambda$ such that $\sh{F}_\mu$ be so.
\end{proposition}

\begin{proof}
The condition being trivially sufficient, let us prove that it is necessary.
If $\sh{F}$ is locally free (resp.\ locally free of rank $n$), there exists a finite open affine covering $(\mathrm{V}_i)$ of $\mathrm{S}$ such that $\sh{F}|_{\mathrm{V}_i}$ is isomorphic to $\sh{O}_\mathrm{S}^{n_i}|_{\mathrm{V}_i}$ (resp.\ $\sh{O}_\mathrm{S}^n|_{\mathrm{V}_i}$) for every $i$.
By virtue of \sref{IV.8.2.11}, there exist $\nu\geqslant\lambda$ and for each $i$ a quasi-compact open subset $\mathrm{U}_{i\nu}$ of $\mathrm{S}_\nu$ such that $\mathrm{V}_i=u_\nu^{-1}(\mathrm{U}_{i\nu})$.
Since $\mathrm{S}_\nu$ is quasi-compact, each $\mathrm{U}_{i\nu}$ is a finite union of open affines; we are thus reduced to the case where $\mathrm{S}_0$ is affine and $\sh{F}=\sh{O}_\mathrm{S}^n$; we then know that there exist $\mu\geqslant\lambda$ such that $\sh{F}_\mu$ is isomorphic to $\sh{O}_{\mathrm{S}_\mu}^n$ \sref{IV.8.5.2.5}.
\end{proof}

\begin{proposition}[8.5.6]
\label{IV.8.5.6}
Suppose $\mathrm{S}_0$ quasi-compact, and consider a sequence
\[
  \sh{F}_\lambda\to\sh{G}_\lambda\to\sh{H}_\lambda\to 0
\]
of homomorphisms of $\sh{O}_{\mathrm{S}_\lambda}$-Modules quasi-coherent, where $\sh{F}_\lambda$ and $\sh{G}_\lambda$ are of finite type and $\sh{H}_\lambda$ of finite presentation.
For the corresponding sequence $\sh{F}\to\sh{G}\to\sh{H}\to 0$ to be exact, it is necessary and sufficient that there exist $\mu\geqslant\lambda$ such that the sequence $\sh{F}_\mu\to\sh{G}_\mu\to\sh{H}_\mu\to 0$ be so (in which case it is the same for $\nu\geqslant\mu$).
\end{proposition}

\begin{proof}
The fact that the condition is sufficient and the last assertion result from the fact that the functor $u_\lambda^*$ (resp.\ $u_{\mu\nu}^*$) is right-exact.
To prove that the condition is necessary, note that it follows from the hypothesis and from \sref{IV.8.5.2},~(i) that there exists $\nu\geqslant\lambda$ such that the composite $\sh{F}_\nu\to\sh{G}_\nu\to\sh{H}_\nu$ is zero.
If one sets $\sh{H}'_\nu=\operatorname{Coker}(\sh{F}_\nu\to\sh{G}_\nu)$, one has therefore a homomorphism $f_\nu:\sh{H}'_\nu\to\sh{H}_\nu$; by hypothesis, $u_\nu^*(f_\nu)$ is an isomorphism, and it follows from \sref{IV.8.5.2.4} that there exists $\mu\geqslant\nu$ such that $u_{\nu\mu}^*(f_\nu)$ is an isomorphism, which completes the proof.
\end{proof}

\begin{corollary}[8.5.7]
\label{IV.8.5.7}
Suppose $\mathrm{S}_0$ quasi-compact, $\sh{F}_\lambda$ quasi-coherent of finite type, and let $f_\lambda:\sh{F}_\lambda\to\sh{G}_\lambda$ be a homomorphism.
For $f=u_\lambda^*(f_\lambda):\sh{F}\to\sh{G}$ to be surjective, it is necessary and sufficient that there exist $\mu\geqslant\lambda$ such that $f_\mu=u_{\lambda\mu}^*(f_\lambda):\sh{F}_\mu\to\sh{G}_\mu$ be so.
\end{corollary}

\begin{proof}
This is the particular case of \sref{IV.8.5.6} applied to the sequence $0\to\sh{H}_\lambda\to 0\to 0$, where $\sh{H}_\lambda=\operatorname{Coker}(f_\lambda)$, which is quasi-coherent and of finite type (taking into account the fact that one has $\sh{H}=\operatorname{Coker}(f)$ and $\sh{H}_\mu=\operatorname{Coker}(f_\mu)$, by virtue of the right-exactness of the functors $u_\lambda^*$ and $u_{\lambda\mu}^*$).
\end{proof}

\begin{corollary}[8.5.8]
\label{IV.8.5.8}
Suppose $\mathrm{S}_0$ quasi-compact and the morphisms $u_{\lambda\mu}:\mathrm{S}_\mu\to\mathrm{S}_\lambda$ flat.
Then:
\begin{enumerate}[label=(\roman*)]
  \item Let $\sh{F}_\lambda\xrightarrow{f_\lambda}\sh{G}_\lambda\xrightarrow{g_\lambda}\sh{H}_\lambda$ be a sequence of homomorphisms of $\sh{O}_{\mathrm{S}_\lambda}$-Modules quasi-coherent, such that $\operatorname{Im} f_\lambda$ and $\operatorname{Ker} g_\lambda$ are of finite type.
For the corresponding sequence $\sh{F}\xrightarrow{f}\sh{G}\xrightarrow{g}\sh{H}$ to be
exact, it is necessary and sufficient that there exist $\mu\geqslant\lambda$ such that the sequence $\sh{F}_\mu\xrightarrow{f_\mu}\sh{G}_\mu\xrightarrow{g_\mu}\sh{H}_\mu$ be so (in which case the same holds for $\nu\geqslant\mu$).
  \item Let $f_\lambda:\sh{F}_\lambda\to\sh{G}_\lambda$ be a homomorphism of quasi-coherent $\sh{O}_{\mathrm{S}_\lambda}$-Modules, and suppose that $\operatorname{Ker} f_\lambda$ is of finite type.
For $f=u_\lambda^*(f_\lambda):\sh{F}\to\sh{G}$ to be injective, it is necessary and sufficient that there exist $\mu\geqslant\lambda$ such that $f_\mu=u_{\lambda\mu}^*(f_\lambda):\sh{F}_\mu\to\sh{G}_\mu$ be so.
\end{enumerate}
\end{corollary}

\begin{proof}
(i) Taking account of \sref{IV.8.3.8},~(ii), observe that, by flatness, $\operatorname{Im} f$ and $\operatorname{Ker} g$ (resp. $\operatorname{Im} f_\mu$ and $\operatorname{Ker} g_\mu$ for $\mu\geqslant\lambda$) are the inverse images of $\operatorname{Im} f_\lambda$ and $\operatorname{Ker} g_\lambda$ \sref[0\textsubscript{I}]{0.6.7.2}. Suppose that the sequence $\sh{F}\xrightarrow{f}\sh{G}\xrightarrow{g}\sh{H}$ is exact. Since $\operatorname{Im} f_\lambda$ is of finite type, it follows from \sref{IV.8.5.2},~(i) that there exists $\mu\geqslant\lambda$ such that the composite $\sh{F}_\mu\xrightarrow{f_\mu}\sh{G}_\mu\xrightarrow{g_\mu}\sh{H}_\mu$ is zero. Changing notation, we may therefore suppose already that $g_\lambda\circ f_\lambda=0$. Since $\operatorname{Ker} g_\lambda$ is of finite type, it then follows from \sref{IV.8.5.7} that there exists $\mu\geqslant\lambda$ such that the homomorphism $\sh{F}_\mu\to\operatorname{Ker} g_\mu$ is surjective, which proves (i).

(ii) This is the special case of (i) applied to the sequence $0\to\sh{F}_\lambda\to\sh{G}_\lambda$.
\end{proof}

\begin{lemma}[8.5.9]
\label{IV.8.5.9}
Suppose that $\mathrm{S}_0$ is quasi-compact and quasi-separated, and that $\sh{F}_\lambda$ is quasi-coherent and of finite type. Let $\sh{G}'_\lambda$ and $\sh{G}''_\lambda$ be two quasi-coherent quotients of $\sh{F}_\lambda$, with $\sh{G}'_\lambda$ moreover supposed to be of finite presentation. For $\sh{G}''$ to be a quotient of $\sh{G}'$, it is necessary and sufficient that there exist $\mu\geqslant\lambda$ such that $\sh{G}''_\mu$ is a quotient of $\sh{G}'_\mu$.
\end{lemma}

\begin{proof}
By hypothesis, there are two surjective homomorphisms $p'_\lambda:\sh{F}_\lambda\to\sh{G}'_\lambda$ and $p''_\lambda:\sh{F}_\lambda\to\sh{G}''_\lambda$. By the right exactness of $u_\lambda^*$ and $u_{\lambda\mu}^*$, the homomorphisms $p'=u_\lambda^*(p'_\lambda)$, $p''=u_\lambda^*(p''_\lambda)$, $p'_\mu=u_{\lambda\mu}^*(p'_\lambda)$, and $p''_\mu=u_{\lambda\mu}^*(p''_\lambda)$ are also surjective. Moreover, if there exists a homomorphism $f:\sh{G}'\to\sh{G}''$ (resp. $f_\mu:\sh{G}'_\mu\to\sh{G}''_\mu$) such that $p''=f\circ p'$ (resp. $p''_\mu=f_\mu\circ p'_\mu$), that homomorphism is necessarily unique; this shows that the question is local on $\mathrm{S}_\lambda$, and consequently, since $\mathrm{S}_0$ is quasi-compact and $\mathbf{L}$ is filtered, that we may suppose $\mathrm{S}_\lambda$ affine and therefore quasi-separated. The condition in the statement is plainly sufficient. Conversely, since $\sh{G}'_\lambda$ is of finite presentation and $\mathrm{S}_\lambda$ is quasi-compact and quasi-separated, it follows from \sref{IV.8.5.2},~(i) that if there exists $f:\sh{G}'\to\sh{G}''$ such that $p''=f\circ p'$, then there exist $\mu\geqslant\lambda$ and $f_\mu:\sh{G}'_\mu\to\sh{G}''_\mu$ such that $f=u_\mu^*(f_\mu)$ and $p''_\mu=f_\mu\circ p'_\mu$, proving the lemma.
\end{proof}

\begin{env}[8.5.10]
\label{IV.8.5.10}
To finish this section, for every quasi-coherent Module $\sh{F}$ on a prescheme, denote by $\mathfrak{Q}(\sh{F})$ the set of quotient Modules of $\sh{F}$ that are of finite presentation. If $\sh{F}_\lambda$ is quasi-coherent and $\sh{G}_\lambda\in\mathfrak{Q}(\sh{F}_\lambda)$, it follows from the right exactness of $u_{\lambda\mu}^*$ and $u_\lambda^*$ that $\sh{G}_\mu\in\mathfrak{Q}(\sh{F}_\mu)$ for $\mu\geqslant\lambda$ and $\sh{G}\in\mathfrak{Q}(\sh{F})$. It is clear that $(\mathfrak{Q}(\sh{F}_\lambda),u_{\lambda\mu}^*)$ is an inductive system of sets and that the $u_\lambda^*:\mathfrak{Q}(\sh{F}_\lambda)\to\mathfrak{Q}(\sh{F})$ form an inductive system of maps; hence, by passing to the inductive limit, there is a canonical map
\[
u_{\mathfrak{Q}}:\varinjlim\mathfrak{Q}(\sh{F}_\lambda)\longrightarrow\mathfrak{Q}(\sh{F}).
\tag{8.5.10.1}\label{IV.8.5.10.1}
\]
Furthermore, if $(\sh{F}'_\lambda)$ is a second family of quasi-coherent $\sh{O}_{\mathrm{S}_\lambda}$-Modules and, for every $\lambda$, $\sh{F}'_\lambda$ is a quotient of $\sh{F}_\lambda$, then $\sh{F}'$ is a quotient of $\sh{F}$, and one has a commutative diagram

\oldpage[IV-3]{25}

\[
\begin{tikzcd}
\varinjlim\mathfrak{Q}(\sh{F}_\lambda) \ar[r,"u_{\mathfrak{Q}}"] & \mathfrak{Q}(\sh{F}) \\
\varinjlim\mathfrak{Q}(\sh{F}'_\lambda) \ar[r,"u_{\mathfrak{Q}}"'] \ar[u] & \mathfrak{Q}(\sh{F}') \ar[u]
\end{tikzcd}
\tag{8.5.10.2}\label{IV.8.5.10.2}
\]
\end{env}

\begin{proposition}[8.5.11]
\label{IV.8.5.11}
Suppose that $\mathrm{S}_0$ is quasi-compact (resp. quasi-compact and quasi-separated). Suppose that $\sh{F}_\lambda$ is quasi-coherent and of finite type (resp. of finite presentation) for every $\lambda$; then the canonical map \eqref{IV.8.5.10.1} is injective (resp. bijective).
\end{proposition}

\begin{proof}
The first assertion follows more precisely from Lemma \sref{IV.8.5.9}. For the second, consider a quotient $\sh{O}_{\mathrm{S}}$-Module $\sh{G}$ of $\sh{F}$ that is of finite presentation. It follows from \sref{IV.8.5.2},~(ii) that there exist $\lambda'\in\mathbf{L}$ and a quasi-coherent $\sh{O}_{\mathrm{S}_{\lambda'}}$-Module $\sh{G}_{\lambda'}$ of finite presentation such that $\sh{G}=u_{\lambda'}^*(\sh{G}_{\lambda'})$; since $\mathbf{L}$ is filtered, we may replace $\lambda$ and $\lambda'$ by an index majorizing both and suppose that $\lambda'=\lambda$. Consider then the canonical homomorphism $p:\sh{F}\to\sh{G}$. It follows from \sref{IV.8.5.2},~(i) that there exist $\mu\geqslant\lambda$ and a homomorphism $p_\mu:\sh{F}_\mu\to\sh{G}_\mu$ such that $p=u_\mu^*(p_\mu)$. By \sref{IV.8.5.7}, we may moreover suppose that $\mu$ is large enough for $p_\mu$ to be surjective, which completes the proof.
\end{proof}

\subsection{Sub-preschemes of finite presentation of a projective limit of preschemes}
\label{subsection:IV.8.6}

\begin{env}[8.6.1]
\label{IV.8.6.1}
Given a prescheme $Y$, we denote in this section by $\mathfrak{Spr}(Y)$ the ordered set
\sref[I]{I.4.1.10} of sub-preschemes of $Y$ that are \emph{of finite presentation}
\sref[I]{I.1.6.1}, by $\mathfrak{Spro}(Y)$ (resp.\ $\mathfrak{Sprf}(Y)$) the subset of $\mathfrak{Spr}(Y)$ formed by the sub-preschemes induced on opens (resp.\ on closed sub-preschemes) of $Y$, of finite presentation over $Y$; it amounts to the same to say that a sub-prescheme $Z$ of $Y$ belongs to $\mathfrak{Spro}(Y)$ (resp.\ $\mathfrak{Sprf}(Y)$) or that it is induced on an open and that the underlying space is \emph{retrocompact} in $Y$ (resp.\ that it is closed and that the ideal $\sh{J}$ of $\sh{O}_Y$ that defines it is \emph{of finite type}, which also means that $j_*(\sh{O}_Z)\in\mathfrak{Q}(\sh{O}_Y)$ where $j:Z\to Y$ is the canonical injection) \sref[I]{I.1.6.1} and \sref[I]{I.1.4.5}.
\end{env}

\begin{env}[8.6.2]
\label{IV.8.6.2}
We always preserve the notations of \sref{IV.8.2.2} and suppose that $S_0$ is one of the $S_\lambda$.
Let $Y_\lambda$ be a sub-prescheme of $S_\lambda$; then $Y_\mu=u_{\mu\lambda}^{-1}(Y_\lambda)$ (resp.\ $Y=u_\lambda^{-1}(Y_\lambda)$) is a sub-prescheme of $S_\mu$ for $\mu\geqslant\lambda$ (resp.\ of $S$); it is induced on an open (resp.\ it is closed) and is of finite presentation on $S_\mu$ (resp.\ $S$) if $Y_\lambda$ is of finite presentation on $S_\lambda$ (resp.\ $S$) \sref[I]{I.4.3.2}. It follows that $(\mathfrak{Spr}(S_\lambda),u_{\lambda\mu}^{-1})$ (resp.\ $(\mathfrak{Spro}(S_\lambda),u_{\lambda\mu}^{-1})$, $(\mathfrak{Sprf}(S_\lambda),u_{\lambda\mu}^{-1})$) is an \emph{inductive system} of sets and the maps $u_\lambda^{-1}:\mathfrak{Spr}(S_\lambda)\to\mathfrak{Spr}(S)$ (resp.\ $\mathfrak{Spro}(S_\lambda)\to\mathfrak{Spro}(S)$, $\mathfrak{Sprf}(S_\lambda)\to\mathfrak{Sprf}(S)$) form an inductive system of maps; whence, by passing to the inductive limit, canonical maps
\end{env}

\begin{align}
  \varinjlim \mathfrak{Spr}(S_\lambda) &\to \mathfrak{Spr}(S) \tag{8.6.2.1}\label{IV.8.6.2.1}\\
  \varinjlim \mathfrak{Spro}(S_\lambda) &\to \mathfrak{Spro}(S) \tag{8.6.2.2}\label{IV.8.6.2.2}\\
  \varinjlim \mathfrak{Sprf}(S_\lambda) &\to \mathfrak{Sprf}(S). \tag{8.6.2.3}\label{IV.8.6.2.3}
\end{align}

\oldpage[IV-3]{26}

Recall \sref[I]{I.4.1.10} that $\mathfrak{Spr}(X)$, for every prescheme $X$, is an ordered set by the relation ``$Z$ is a sub-prescheme of the sub-prescheme $Y$'' which we write $Z\leqslant Y$.
The maps $u_{\lambda\mu}^{-1}$ and $u_\lambda^{-1}$ are \emph{increasing} for the corresponding order relations in $\mathfrak{Spr}(S_\lambda)$, $\mathfrak{Spr}(S_\mu)$, $\mathfrak{Spr}(S)$.
Furthermore, by writing that $\eta\leqslant\eta'$ for two elements of this set whenever there exist $\lambda$ and two elements $Y_\lambda$, $Y'_\lambda$ of $\mathfrak{Spr}(S_\lambda)$, of which $\eta$ and $\eta'$ are the canonical images, and which are such that $Y_\lambda\leqslant Y'_\lambda$; one easily verifies that this does not depend on the choice of representatives, and that one thus has indeed a \emph{relation of order}.
The fact that the $u_\lambda^{-1}$ are increasing immediately implies that the canonical map \eqref{IV.8.6.2.1} is \emph{increasing}; the same is evidently true of \eqref{IV.8.6.2.2} and \eqref{IV.8.6.2.3}.

\begin{proposition}[8.6.3]
\label{IV.8.6.3}
Suppose $S_0$ quasi-compact (resp.\ quasi-compact and quasi-separated).
Then the maps \eqref{IV.8.6.2.1}, \eqref{IV.8.6.2.2}, \eqref{IV.8.6.2.3} are injective (resp.\ bijective).
\end{proposition}

\begin{proof}
Taking into account the remarks of \sref{IV.8.6.1}, the assertions relative to \eqref{IV.8.6.2.3} follow from \sref{IV.8.5.11} applied to $\sh{F}_\lambda=\sh{O}_{S_\lambda}$; similarly, the assertions relative to \eqref{IV.8.6.2.2} are special cases of \sref{IV.8.3.5} and \sref{IV.8.3.11}, taking into account that the $S_\lambda$ and $S$ are quasi-compact.
It remains to consider the map \eqref{IV.8.6.2.1}.
Let us first prove that it is surjective when $S_0$ is quasi-compact and quasi-separated.
Let $Z$ be a sub-prescheme of $S$, of finite presentation over $S$; as $S$ is of the same nature as $Z$, there exists an open quasi-compact $U$ of $S$ such that $Z$ is a sub-prescheme of $U$, closed in $U$, of finite presentation over $U$.
There then exist an index $\lambda$ and an open quasi-compact $U_\lambda$ of $S_\lambda$ such that $U=u_\lambda^{-1}(U_\lambda)$ \sref{IV.8.2.11}; as $S_\lambda$ is quasi-separated, it is of the same nature as $U_\lambda$ \sref[I]{I.1.2.7}, and consequently one can restrict to the case where $U=S$; but in this case, one is reduced to the fact that \eqref{IV.8.6.2.3} is surjective.

Finally, to see that \eqref{IV.8.6.2.1} is injective when $S_0$ is quasi-compact, it will suffice to prove the following more precise result:
\end{proof}

\begin{corollary}[8.6.3.1]
\label{IV.8.6.3.1}
Suppose $S_0$ quasi-compact and let $Z'_\lambda$, $Z''_\lambda$ be two sub-preschemes of $S_\lambda$, of finite presentation.
For $Z'=u_\lambda^{-1}(Z'_\lambda)$ to be majorised by $Z''=u_\lambda^{-1}(Z''_\lambda)$ \sref[I]{I.4.1.10}, it is necessary and sufficient that there exist $\mu\geqslant\lambda$ such that $Z'_\mu=u_{\mu\lambda}^{-1}(Z'_\lambda)$ be majorised by $Z''_\mu=u_{\mu\lambda}^{-1}(Z''_\lambda)$.
\end{corollary}

\begin{proof}
It is trivial that the condition is sufficient.
To see that it is necessary, note first that the underlying sets $Z'_\lambda$ and $Z''_\lambda$ are locally constructible in $S_\lambda$ by hypothesis \sref[I]{I.1.8.4}, hence the hypothesis implies the existence of a $\nu\geqslant\lambda$ such that $Z'_\lambda\subset Z''_\lambda$ \sref{IV.8.3.5}; replacing $\lambda$ by $\nu$, we can therefore already suppose that $Z'_\lambda\subset Z''_\lambda$.
Furthermore, by hypothesis \sref[I]{I.1.6.1}, the sub-spaces $Z'_\lambda$ and $Z''_\lambda$ of $S_\lambda$ are quasi-compact; for every point $x\in Z'_\lambda$, there is an open quasi-compact neighbourhood $V(x)$ in $S_\lambda$ such that $V(x)\cap Z'_\lambda$ and $V(x)\cap Z''_\lambda$ are closed in $V(x)$.
By covering $S_\lambda$ by a finite number of such neighbourhoods, we see that there is an open quasi-compact $U_\lambda$ of $S_\lambda$ containing $Z'_\lambda$ and such that $Z'_\lambda$ and $Z''_\lambda\cap U_\lambda$ are closed in $U_\lambda$.
If we denote by $Y_\lambda$ the sub-prescheme of $S_\lambda$ induced by $Z''_\lambda$ on $U_\lambda\cap Z''_\lambda$, it is clear that with the usual notations, $Y_\mu$ (resp.\ $Y$) is majorised by $Z''_\mu$ (resp.\ $Z''$) for $\mu\geqslant\lambda$ \sref[I]{I.4.4.1}, and as it suffices to prove that $Z'_\mu$ is majorised by $Y_\mu$ for $\mu$
\oldpage[IV-3]{27}
large enough, we see finally that one is reduced (replacing $S_\lambda$ by $U_\lambda$) to the case where $Z'_\lambda$ and $Z''_\lambda$ are closed sub-preschemes of $S_\lambda$.
But then, this has already been proved since \eqref{IV.8.6.2.3} is increasing and injective.
\end{proof}

\begin{corollary}[8.6.4]
\label{IV.8.6.4}
Suppose $S_0$ quasi-compact, and let $Z_\lambda$ be a sub-prescheme of $S_\lambda$, of finite presentation over $S_\lambda$.
For $Z=u_\lambda^{-1}(Z_\lambda)$ to be the subprescheme induced on an open subset of $S$ (resp.\ a closed subprescheme of $S$), it is necessary and sufficient that there exist $\mu\geqslant\lambda$ such that $Z_\mu=u_{\mu\lambda}^{-1}(Z_\lambda)$ is the subprescheme induced on an open subset of $S_\mu$ (resp.\ a closed subprescheme of $S_\mu$).
\end{corollary}

\begin{proof}
Let $(U_\lambda^{(i)})$ be a finite affine open covering of $S_\lambda$, and set $U^{(i)}_\mu=u_{\mu\lambda}^{-1}(U_\lambda^{(i)})$ for $\mu\geqslant\lambda$ and $U^{(i)}=u_\lambda^{-1}(U_\lambda^{(i)})$.
If $Z$ is open (resp.\ closed) in $S$, $Z\cap U^{(i)}$ is open (resp.\ closed) in $U^{(i)}$, and conversely if each of the $Z_\mu\cap U^{(i)}_\mu$ is open (resp.\ closed) in $U^{(i)}_\mu$, $Z_\mu$ is open (resp.\ closed) in $S_\mu$.
Since $L$ is filtered, it suffices to prove the corollary when $S_\lambda$ is affine, hence quasi-separated, and the result follows from the fact that the maps \eqref{IV.8.6.2.1}, \eqref{IV.8.6.2.2}, and \eqref{IV.8.6.2.3} are bijective.
\end{proof}

\subsection{Criteria for a projective limit of preschemes to be a reduced (resp.\ integral) prescheme}
\label{subsection:IV.8.7}

We preserve the hypotheses and notations of \sref{IV.8.2.2} and suppose always that $S_0$ is one of the $S_\lambda$.

\begin{proposition}[8.7.1]
\label{IV.8.7.1}
Suppose that $S$ is not reduced.
Then there exists $\lambda_0$ such that for $\lambda\geqslant\lambda_0$, $S_\lambda$ is not reduced.
\end{proposition}

\begin{proof}
The question being local on $S_0$, we can suppose $S_0=\operatorname{Spec}(A_0)$ affine, whence $S=\operatorname{Spec}(A)$, where $A=\varinjlim A_\lambda$ is the inductive limit of an inductive system of $A_0$-algebras $(A_\lambda)$.
We know then \sref{IV.5.13.2} that the nilradical of $A$ is the inductive limit of those of the $A_\lambda$; if it is not zero, one of the $A_\lambda$ contains therefore a nilpotent element $a_\lambda\neq 0$, and the image of $a_\lambda$ in the $A_\mu$ for $\mu\geqslant\lambda$ is therefore a nilpotent element $\neq 0$, and the image of $a_\lambda$ in $A$ is not zero for $\mu\geqslant\lambda$ since the image in $A$ is not zero.
\end{proof}

\begin{proposition}[8.7.2]
\label{IV.8.7.2}
Suppose one of the following hypotheses satisfied:
\begin{enumerate}[label=\emph{\alph*)}]
\item $S_0$ is quasi-compact, the nilradical $\sh{N}_0$ of $\sh{O}_{S_0}$ is an ideal of finite type (which is satisfied for example when $S_0$ is Noetherian), and the morphisms $u_{\lambda\mu}:S_\mu\to S_\lambda$ are open immersions.
\item The morphisms $u_{\lambda\mu}:S_\mu\to S_\lambda$ are faithfully flat.
\end{enumerate}
Under these conditions, for $S$ to be reduced, it is necessary and sufficient that there exist $\lambda_0$ such that $S_\lambda$ be reduced for $\lambda\geqslant\lambda_0$.
\end{proposition}

\begin{proof}
Furthermore, in case b), if $S$ is reduced, all the $S_\lambda$ are reduced.

The last assertion follows from the fact that the morphism $S\to S_\lambda$ is then faithfully flat for every $\lambda$ \sref{IV.8.3.8} and from \sref{IV.2.1.13}.
On the other hand, \sref{IV.8.7.1} proves that the condition of the statement is sufficient (without hypothesis on $S_0$ or on the $u_{\lambda\mu}$).
It therefore remains to see that the condition is necessary in hypothesis $a)$; then \sref{IV.8.2.13}, the underlying space of $S$ is identified with $S$ (the $S_\lambda$ being identified with sub-preschemes induced on opens of $S_0$), and the structural sheaf $\sh{O}_S$ is identified with
\oldpage[IV-3]{28}
the sheaf induced on $S$ by all the $\sh{O}_{S_\lambda}$; in particular for every $s\in S$, the local ring $\sh{O}_s$ is the same for $S$ and for all the $S_\lambda$.
If $\sh{N}_\lambda$ is the nilradical of $\sh{O}_{S_\lambda}$, the nilradical $\sh{N}$ of $\sh{O}_S$ has at each point of $S$ the same fibre $\sh{N}_s$ (induced on $S$ by $\sh{N}_\lambda$).
The hypothesis that $S$ is reduced implies therefore $u_\lambda^*(\sh{N}_\lambda)=0$; as $\sh{N}_0$ is supposed of finite type, the same is true of the $\sh{N}_\lambda$, and the conclusion follows then from \sref{IV.8.5.7}.
\end{proof}

\begin{corollary}[8.7.3]
\label{IV.8.7.3}
Suppose one of the following hypotheses satisfied:
\begin{enumerate}[label=\emph{\alph*)}]
\item $S_0$ is a Noetherian prescheme and the morphisms $u_{\lambda\mu}:S_\mu\to S_\lambda$ are open immersions.
\item The morphisms $u_{\lambda\mu}:S_\mu\to S_\lambda$ are faithfully flat.
\end{enumerate}
Then, for $S$ to be integral, it is necessary and sufficient that there exist $\lambda_0$ such that $S_\lambda$ be integral for $\lambda\geqslant\lambda_0$.
\end{corollary}

\begin{proof}
To say that a prescheme is integral means that it is both reduced and irreducible; the corollary thus follows from \sref{IV.8.7.2} and from \sref{IV.8.4.3}.
\end{proof}

\begin{remark}[8.7.4]
\label{IV.8.7.4}
If one makes no hypothesis on the $u_{\lambda\mu}$, it can happen that $S$ is integral even though all the $S_\lambda$ are neither reduced nor connected, as the example \sref{IV.8.4.4} shows, where one takes the ring $A$ non reduced.
\end{remark}

\subsection{Preschemes of finite presentation over a projective limit of preschemes}
\label{subsection:IV.8.8}

\begin{env}[8.8.1]
\label{IV.8.8.1}
Always preserving the notations and hypotheses of \sref{IV.8.2.2}, we shall suppose given in this section two $S_\alpha$-preschemes $X_\alpha$, $Y_\alpha$, which defines \sref{IV.8.2.5} two projective systems of preschemes $(X_\lambda,v_{\lambda\mu})$ and $(Y_\lambda,w_{\lambda\mu})$ by setting $X_\lambda=X_\alpha\times_{S_\alpha}S_\lambda$, $v_{\lambda\mu}=1_{X_\alpha}\times u_{\lambda\mu}$, $w_{\lambda\mu}=1_{Y_\alpha}\times u_{\lambda\mu}$ (for $\alpha\leqslant\lambda\leqslant\mu$), whose projective limits are respectively $X=X_\alpha\times_{S_\alpha}S$, $Y=Y_\alpha\times_{S_\alpha}S$, with the canonical morphisms $v_\lambda:X\to X_\lambda$ and $w_\lambda:Y\to Y_\lambda$.
For $\alpha\leqslant\lambda\leqslant\mu$, there is a canonical map $\varepsilon_{\mu\lambda}:\operatorname{Hom}_{S_\lambda}(X_\lambda,Y_\lambda)\to\operatorname{Hom}_{S_\mu}(X_\mu,Y_\mu)$, which to every $S_\lambda$-morphism $f_\lambda:X_\lambda\to Y_\lambda$ associates $f_\mu=f_\lambda\times 1_{S_\mu}:X_\lambda\times_{S_\lambda}S_\mu\to Y_\lambda\times_{S_\lambda}S_\mu$, and it is clear that $(\operatorname{Hom}_{S_\lambda}(X_\lambda,Y_\lambda),\varepsilon_{\mu\lambda})$ is an \emph{inductive system} of sets.
Similarly, there is a canonical map $\varepsilon_\lambda:\operatorname{Hom}_{S_\lambda}(X_\lambda,Y_\lambda)\to\operatorname{Hom}_S(X,Y)$ which to $f_\lambda$ associates $f=f_\lambda\times 1_S:X_\lambda\times_{S_\lambda}S\to Y_\lambda\times_{S_\lambda}S$ and $(\varepsilon_\lambda)$ is an inductive system of maps; whence, by passing to the inductive limit, a canonical map, functorial in $S_\alpha$, $X_\alpha$, and $Y_\alpha$:
\[
  \varepsilon:\varinjlim\operatorname{Hom}_{S_\lambda}(X_\lambda,Y_\lambda)\to\operatorname{Hom}_S(X,Y).
  \tag{8.8.1.1}\label{IV.8.8.1.1}
\]
\end{env}

\begin{theorem}[8.8.2]
\label{IV.8.8.2}
\begin{enumerate}[label=(\roman*)]
\item Suppose $X_\alpha$ quasi-compact (resp.\ quasi-compact and quasi-separated), and $Y_\alpha$ locally of finite type (resp.\ locally of finite presentation) over $S_\alpha$.
Then the map \eqref{IV.8.8.1.1} is injective (resp.\ bijective).
\item Suppose $S_0$ quasi-compact and quasi-separated.
For every prescheme $X$ of finite presentation over $S$, there exist a $\lambda\in L$, a prescheme $X_\lambda$ of finite presentation over $S_\lambda$, and an $S$-isomorphism $X\xrightarrow{\sim}X_\lambda\times_{S_\lambda}S$.
\end{enumerate}
\end{theorem}

\begin{proof}
\oldpage[IV-3]{29}
(i) Consider first the case where $S_0=\operatorname{Spec}(A_0)$, $X_\alpha=\operatorname{Spec}(B_\alpha)$, and $Y_\alpha=\operatorname{Spec}(C_\alpha)$ are affine; then the $S_\lambda=\operatorname{Spec}(A_\lambda)$ and $S=\operatorname{Spec}(A)$ are also affine, with $A=\varinjlim A_\lambda$, and the assertions of (i) are equivalent to the

\begin{lemma}[8.8.2.1]
\label{IV.8.8.2.1}
Let $A_0$ be a ring, $(A_\lambda)_{\lambda\in L}$ an inductive system of $A_0$-algebras, $A=\varinjlim A_\lambda$; let $B_\alpha$ be an $A_\alpha$-algebra, $C_\alpha$ an $A_\alpha$-algebra of finite type (resp.\ of finite presentation).
Then the canonical homomorphism
\[
  \varinjlim_\lambda\operatorname{Hom}_{A_\lambda\text{-alg.}}(C_\alpha\otimes_{A_\alpha}A_\lambda,B_\alpha\otimes_{A_\alpha}A_\lambda)\to\operatorname{Hom}_{A\text{-alg.}}(C_\alpha\otimes_{A_\alpha}A,B_\alpha\otimes_{A_\alpha}A)
  \tag{8.8.2.2}\label{IV.8.8.2.2}
\]
is injective (resp.\ bijective).
\end{lemma}

\begin{proof}
We know that there are functorial canonical isomorphisms
\[
  \operatorname{Hom}_{A_\lambda\text{-alg.}}(C_\alpha\otimes_{A_\alpha}A_\lambda,B_\alpha\otimes_{A_\alpha}A_\lambda)\xrightarrow{\sim}\operatorname{Hom}_{A_\alpha\text{-alg.}}(C_\alpha,B_\alpha\otimes_{A_\alpha}A_\lambda)
\]
\[
  \operatorname{Hom}_{A\text{-alg.}}(C_\alpha\otimes_{A_\alpha}A,B_\alpha\otimes_{A_\alpha}A)\xrightarrow{\sim}\operatorname{Hom}_{A_\alpha\text{-alg.}}(C_\alpha,B_\alpha\otimes_{A_\alpha}A)
\]
by virtue of the universal property of the tensor product of two algebras.
It therefore suffices to prove the
\end{proof}

\begin{lemma}[8.8.2.3]
\label{IV.8.8.2.3}
Let $E$ be a ring, $G$ an $E$-algebra of finite type (resp.\ of finite presentation), $(F_\lambda)$ an inductive system of $E$-algebras.
Then the canonical homomorphism
\[
  \varinjlim\operatorname{Hom}_{E\text{-alg.}}(G,F_\lambda)\to\operatorname{Hom}_{E\text{-alg.}}(G,\varinjlim F_\lambda)
\]
which, to every inductive system of homomorphisms $\theta_\lambda:G\to F_\lambda$ of $E$-algebras, associates its inductive limit, is injective (resp.\ bijective).
\end{lemma}

\begin{proof}
Suppose first the $E$-algebra $G$ of finite type, and let $(t_i)_{1\leqslant i\leqslant n}$ be a system of generators of this $E$-algebra; let us show that if $(\theta'_\lambda)$, $(\theta''_\lambda)$ are two inductive systems of homomorphisms $G\to F_\lambda$, with $\varinjlim\theta'_\lambda=\varinjlim\theta''_\lambda$, there exists $\mu\geqslant\lambda$ such that $\theta'_\mu=\theta''_\mu$.
Indeed, if $\varphi_{\mu\lambda}:F_\lambda\to F_\mu$ and $\varphi_\lambda:F_\lambda\to F=\varinjlim F_\lambda$ are the homomorphisms of the inductive system $(F_\lambda)$, by hypothesis, for each index $i$, there exists an index $\lambda_i$ such that $\varphi_{\lambda_i}(\theta'_{\lambda_i}(t_i))=\varphi_{\lambda_i}(\theta''_{\lambda_i}(t_i))$, and one can suppose all the $\lambda_i$ equal to a single $\lambda$; it follows in the same way the existence of $\mu\geqslant\lambda$ such that $\varphi_{\mu\lambda}(\theta'_\lambda(t_i))=\varphi_{\mu\lambda}(\theta''_\lambda(t_i))$ for $1\leqslant i\leqslant n$, that is to say $\theta'_\mu(t_i)=\theta''_\mu(t_i)$ for $1\leqslant i\leqslant n$, whence $\theta'_\mu=\theta''_\mu$.

Suppose in the second place $G$ of finite presentation, so that one has
\[
  G=E[T_1,\ldots,T_n]/\mathfrak{J},
\]
where $\mathfrak{J}$ is an ideal of finite type, $t_i$ being the class of $T_i$ modulo $\mathfrak{J}$.
Let $(P_j)_{1\leqslant j\leqslant m}$ be a system of generators of $\mathfrak{J}$.
Suppose given a homomorphism of $E$-algebras $\theta:G\to F$; set $b_i=\theta(t_i)$; by definition, one has therefore $P_j(b_1,b_2,\ldots,b_n)=\theta(P_j(t_1,\ldots,t_n))=0$ for $1\leqslant j\leqslant m$.
Now, there exist $\lambda$ and elements $x_1,\ldots,x_n$ in $F_\lambda$ such that $b_i=\varphi_\lambda(x_i)$ for $1\leqslant i\leqslant n$; one has therefore $\varphi_\lambda(P_j(x_1,\ldots,x_n))=P_j(b_1,\ldots,b_n)=0$, and consequently there exists
\oldpage[IV-3]{30}
$\mu\geqslant\lambda$ such that $\varphi_{\mu\lambda}(P_j(x_1,\ldots,x_n))=P_j(\varphi_{\mu\lambda}(x_1),\ldots,\varphi_{\mu\lambda}(x_n))=0$ for $1\leqslant j\leqslant m$; one therefore deduces for every $\nu\geqslant\mu$ a homomorphism of $E$-algebras $\theta_\nu=\varphi_{\nu\mu}\circ\theta_\mu$ of $G$ into $F_\nu$, and it is clear that $\theta$ is the inductive limit of this inductive system of homomorphisms, which completes the proof of the lemma.
\end{proof}

Let us now pass to the case where $X_\alpha$ is quasi-compact and $Y_\alpha$ locally of finite type over $S_\alpha$.
Set $Z_\alpha=X_\alpha\times_{S_\alpha}Y_\alpha$ and introduce the projective system corresponding to $Z_\lambda=X_\lambda\times_{S_\lambda}Y_\lambda$ and its limit $Z=X\times_S Y$; the canonical bijections \sref[I]{I.3.3.14} give commutative diagrams
\[
\xymatrix{
  \operatorname{Hom}_{S_\lambda}(X_\lambda,Y_\lambda) \ar[r]^-{\varepsilon_{\mu\lambda}} \ar[d]^{\wr} & \operatorname{Hom}_{S_\mu}(X_\mu,Y_\mu) \ar[r]^-{\varepsilon_\mu} \ar[d]^{\wr} & \operatorname{Hom}_S(X,Y) \ar[d]^{\wr}\\
  \operatorname{Hom}_{X_\lambda}(X_\lambda,Z_\lambda) \ar[r]^-{\varepsilon_{\mu\lambda}} & \operatorname{Hom}_{X_\mu}(X_\mu,Z_\mu) \ar[r]^-{\varepsilon_\mu} & \operatorname{Hom}_X(X,Z)
}
\]
and consequently we are reduced to proving that \eqref{IV.8.8.1.1} is injective in the particular case where $S_\alpha=X_\alpha$ (taking into account \sref[I]{I.1.3.4}).
Moreover, as $X_\alpha$ is quasi-compact, hence a finite union of open affines, and as $X$ is a union of the opens $f^{-1}(V_i)$, where the $V_i$ are open affines of $Y_\alpha$, every point of $X$ has an open affine neighbourhood $W(x)$ in $S_\alpha$ such that $f'_\alpha$ and $f''_\alpha$ both map $W(x)$ into a single open affine $V_{i\alpha}$.
As $X_\alpha$ is quasi-compact, it is covered by a finite number of open affines $W(x_k)$; this implies that for each $x\in X_\alpha$, there is an open affine neighbourhood of $x$ such that the restrictions of $f'_\alpha$ and $f''_\alpha$ to $W(x)$ map it into the same open affine $V_{i\alpha}$.
As $X_\alpha$ is quasi-compact, it is covered by a finite number of such open affines $W(x_k)$; by virtue of the lemma \sref{IV.8.8.2.1} and of the fact that $L$ is filtered, there exists $\lambda\geqslant\alpha$ such that the restrictions of $f'_\lambda$ and $f''_\lambda$ to each of the opens $v_\alpha^{-1}(W(x_k))$ are equal, whence $f'_\lambda=f''_\lambda$.

Suppose now $X_\alpha$ quasi-compact and quasi-separated and $Y_\alpha$ locally of finite presentation over $S_\alpha$, and let us prove that \eqref{IV.8.8.1.1} is surjective.
Suppose therefore given an $X$-morphism $f:X\to Y$.
As $X$ is quasi-compact, it is of the same nature as $Y$, which contains $f(X)$; there exists therefore $\lambda\geqslant\alpha$ and an open quasi-compact $Y'_\lambda$ in $Y_\lambda$ such that $Y'=w_\lambda^{-1}(Y'_\lambda)$ \sref{IV.8.2.11}.
Replacing as needed $\alpha$ by $\lambda$ and $Y_\lambda$ by $Y'_\lambda$, we can therefore restrict to the case where $Y_\alpha$ is \emph{quasi-compact}, hence also $Y$ and the $Y_\lambda$.
As $Y$ is quasi-compact, it is a finite union of open affines $V_i$, and $X$ is a union of the opens $f^{-1}(V_i)$.
As $X$ is quasi-compact, we can cover $X$ by a finite number of such open affines.
\oldpage[IV-3]{31}
And since $X$ is quasi-compact, for each $i$, the open quasi-compact sets in $X_\alpha$ may, by replacing $\alpha$ by an index $\lambda\geqslant\alpha$, be supposed such that $U_i=v_\alpha^{-1}(U_{i\alpha})$ where the $U_{i\alpha}$ are open quasi-compacts in $X_\alpha$; moreover, using \sref{IV.8.3.4}, on can moreover suppose the $(U_{i\alpha})$ form a \emph{covering} of $X_\alpha$.

It will then suffice to show that there exist $\lambda\geqslant\alpha$ and for each $i$ a morphism $f_{i\lambda}:U_{i\lambda}\to V_{i\lambda}$ (where $U_{i\lambda}=v_{\alpha\lambda}^{-1}(U_{i\alpha})$ and $V_{i\lambda}=w_{\alpha\lambda}^{-1}(V_{i\alpha})$) such that the morphism $f_i=\varepsilon_\lambda(f_{i\lambda}):U_i\to V_i$ be equal to the \emph{restriction} of $f$ to $U_i$.
Indeed, if this is so, as $X_\alpha$ is quasi-separated \sref[I]{I.1.2.3}, the $U_{i\lambda}\cap U_{j\lambda}$ are \emph{quasi-compact} and the uniqueness result proved above (which applies since $Y_\lambda$ is locally of finite type over $S_\lambda$) proves that there exists $\mu\geqslant\lambda$ such that $f_{i\mu}$ and $f_{j\mu}$ coincide on $U_{i\mu}\cap U_{j\mu}$ for every couple $(i,j)$, and consequently define a $S_\mu$-morphism $f_\mu:X_\mu\to Y_\mu$ such that $\varepsilon_\mu(f_\mu)=f$.

We are thus reduced to the case where $Y_\alpha$ is \emph{affine}, and as in addition we can also suppose that $S_\alpha$ is affine, we can also suppose $S_\alpha=\operatorname{Spec}(A_\alpha)$, $Y_\alpha=\operatorname{Spec}(C_\alpha)$, $C_\alpha$ being an $A_\alpha$-algebra of finite presentation, $S=\operatorname{Spec}(A)$, $Y=\operatorname{Spec}(C)$, with $A=\varinjlim(C_\alpha\otimes_{A_\alpha}A_\lambda)=C_\alpha\otimes_{A_\alpha}A$.
One has then
\[
  \operatorname{Hom}_S(X,Y)=\operatorname{Hom}_{A\text{-alg.}}(C,\Gamma(X,\sh{O}_X))=\operatorname{Hom}_{A_\alpha\text{-alg.}}(C_\alpha,\Gamma(X,\sh{O}_X))
\]
\sref[I]{I.2.2.4} and similarly
\[
  \operatorname{Hom}_{S_\lambda}(X_\lambda,Y_\lambda)=\operatorname{Hom}_{A_\lambda\text{-alg.}}(C_\alpha\otimes_{A_\alpha}A_\lambda,\Gamma(X_\lambda,\sh{O}_{X_\lambda}))=\operatorname{Hom}_{A_\alpha\text{-alg.}}(C_\alpha,\Gamma(X_\lambda,\sh{O}_{X_\lambda})).
\]
But as $X_\alpha$ is quasi-compact and quasi-separated, one knows \sref{IV.8.5.4} that $\varinjlim\Gamma(X_\lambda,\sh{O}_{X_\lambda})=\Gamma(X,\sh{O}_X)$; since $C_\alpha$ is an $A_\alpha$-algebra of finite presentation, the fact that \eqref{IV.8.8.1.1} is bijective then follows from \sref{IV.8.8.2.3}.

Before passing to the proof of (ii), let us note the following corollaries of (i):
\end{proof}

\begin{corollary}[8.8.2.4]
\label{IV.8.8.2.4}
Suppose $S_0$ quasi-compact, $X_\alpha$ of finite presentation over $S_\alpha$, $Y_\alpha$ of finite type over $S_\alpha$ and quasi-separated over $S_\alpha$ (which is the case, for example, if $Y_\alpha$ is also
\oldpage[IV-3]{32}
of finite presentation over $S_\alpha$).
Let $f_\alpha:X_\alpha\to Y_\alpha$ be an $S_\alpha$-morphism.
For $f:X\to Y$ to be an isomorphism, it is necessary and sufficient that there exist $\lambda\geqslant\alpha$ such that $f_\lambda:X_\lambda\to Y_\lambda$ is an isomorphism.
\end{corollary}

\begin{proof}
The condition is evidently sufficient. To prove that it is necessary, first observe that, since the question is local on $S_0$ (because $\mathbf{L}$ is filtered), we may always suppose $S_0$ affine and hence quasi-separated. By hypothesis there is an $S$-morphism $g:Y\to X$ such that $g\circ f=1_X$ and $f\circ g=1_Y$. Since $X_\alpha$ is of finite presentation over $S_\alpha$, and $Y_\alpha$ is quasi-compact and quasi-separated (because $S_\alpha$ is quasi-compact and quasi-separated), it follows from \sref{IV.8.8.2},~(ii) that there exist $\mu\geqslant\alpha$ and an $S_\mu$-morphism $g_\mu:Y_\mu\to X_\mu$ such that $g=g_\mu\times 1_S$. It follows moreover from \sref{IV.8.8.2},~(i), together with the relations $g\circ f=1_X$ and $f\circ g=1_Y$, that there exists $\nu\geqslant\mu$ such that $g_\nu\circ f_\nu=1_{X_\nu}$ and $f_\nu\circ g_\nu=1_{Y_\nu}$, since $X_\alpha$ and $Y_\alpha$ are of finite type over $S_\alpha$ and quasi-compact. Thus $f_\nu:X_\nu\to Y_\nu$ is an isomorphism, proving the corollary.
\end{proof}

\begin{corollary}[8.8.2.5]
\label{IV.8.8.2.5}
Suppose $S_0$ quasi-compact and quasi-separated, $X_\alpha$ and $Y_\alpha$ of finite presentation over $S_\alpha$.
For $X$ and $Y$ to be $S$-isomorphic, it is necessary and sufficient that there exist $\lambda\geqslant\alpha$ such that $X_\lambda$ and $Y_\lambda$ are $S_\lambda$-isomorphic. Moreover, for every $S$-isomorphism $f:X\xrightarrow{\sim}Y$, there exist $\mu\geqslant\lambda$ and an $S_\mu$-isomorphism $f_\mu:X_\mu\xrightarrow{\sim}Y_\mu$ such that $f=f_\mu\times 1_S$.
\end{corollary}

\begin{proof}
The condition is evidently sufficient. Conversely, if there exists an $S$-isomorphism $f:X\to Y$, it follows from \sref{IV.8.8.2},~(ii) that $f$ is of the form $f_\mu\times 1_S$ for some $\mu\geqslant\alpha$ and some $S_\mu$-morphism $f_\mu:X_\mu\to Y_\mu$. Since $f$ is an isomorphism, it follows from \sref{IV.8.8.2.4} that there exists $\nu\geqslant\mu$ such that $f_\nu:X_\nu\to Y_\nu$ is an isomorphism.
\end{proof}

\begin{proof}[Proof of \sref{IV.8.8.2}, (ii)]
Consider first the case where $S_0=\operatorname{Spec}(A_0)$ and $X=\operatorname{Spec}(B)$ are affine; then assertion~(ii) is equivalent to Lemma~\sref[I]{I.8.4.2}.

To prove (ii) in the general case, note that $S$ is quasi-compact and quasi-separated. Since $X$ is of finite presentation over $S$, and $S$ is affine over $S_0$, there is a finite open cover $(U_i)$ of $S_0$ and, if $(W_i)$ denotes the affine open cover of $S$ formed by the inverse images of the $U_i$ under $S\to S_0$, a finite cover $(X_r)$ of $X$ by affine opens such that the image of each $X_r$ under $X\to S$ is contained in some $W_{i(r)}$. The ring $A(X_r)$ is then an algebra of finite presentation over $A(W_{i(r)})$ \sref[I]{I.1.4.6}.

By Lemma~\sref[I]{I.8.4.2} and the fact that $L$ is filtered, there exist an index $\lambda\in L$ and, for every $r$, an affine scheme $Z_{r\lambda}$ of finite presentation over the inverse image $W_{i(r),\lambda}$ of $U_{i(r)}$ in $S_\lambda$, together with an $S$-isomorphism
\[
 g_r:Z_{r\lambda}\times_{S_\lambda}S\xrightarrow{\sim}X_r.
\]
Let $Z_{rs}$ be the inverse image under $g_r$ of the sub-prescheme induced on the open $X_r\cap X_s$ of $X_r$, which is quasi-compact because $X$ is quasi-separated, and let
\[
 g'_{rs}:Z_{rs}\xrightarrow{\sim}X_r\cap X_s
\]
be the restriction of $g_r$. By \sref{IV.8.8.2.5}, there exist $\mu\geqslant\lambda$ and, for every pair $(r,s)$, a quasi-compact open $Z_{rs\mu}$ of $Z_{r\mu}=v_{\lambda\mu}^{-1}(Z_{r\lambda})$ such that $Z_{rs}$ is the inverse image of $Z_{rs\mu}$. Moreover, since $S_\mu$ is quasi-separated and $W_{i(r),\mu}$ is a quasi-compact open of $S_\mu$, each $Z_{rs\mu}$ is of finite presentation over $S_\mu$ \sref[I]{I.1.6.2}.

For every pair $(r,s)$, consider the isomorphism
\[
 h_{sr}=(g'_{sr})^{-1}\circ g'_{rs}:Z_{rs}\xrightarrow{\sim}Z_{sr}.
\]
By \sref{IV.8.8.2.4}, there exist $\nu\geqslant\mu$ and, for every pair $(r,s)$, an isomorphism
\[
 h_{sr\nu}:Z_{rs\nu}\xrightarrow{\sim}Z_{sr\nu}
\]
such that $h_{sr}=h_{sr\nu}\times 1_S$. Finally, for every triple $(r,s,t)$, let $h'_{sr}$ denote the restriction of $h_{sr}$ to

\oldpage[IV-3]{33}

$Z_{rs}\cap Z_{rt}$. It follows immediately from the definitions that $h'_{sr}$ is an isomorphism from $Z_{rs}\cap Z_{rt}$ onto $Z_{sr}\cap Z_{st}$ and that
\[
 h'_{ts}\circ h'_{sr}=h'_{tr}.
\]
By \sref{IV.8.8.2},~(i), there therefore exists $\rho\geqslant\nu$ such that, for every triple $(r,s,t)$,
\[
 h'_{ts\rho}\circ h'_{sr\rho}=h'_{tr\rho}.
\]
Consequently the isomorphisms $h_{sr\rho}$ satisfy the gluing condition \sref[0\textsubscript{I}]{0.4.1.7} and hence define a prescheme $X_\rho$ such that $X$ is isomorphic to $X_\rho\times_{S_\rho}S$. Moreover, the $Z_{r\rho}$ are of finite presentation over $S_\rho$, and, identifying them with open subsets of $X_\rho$, their intersections $Z_{r\rho}\cap Z_{s\rho}$, which are isomorphic to the $Z_{rs\rho}$, are quasi-compact. Thus $X_\rho$ is of finite presentation over $S_\rho$ \sref[I]{I.1.6.3}.
\end{proof}

\begin{env}[8.8.3]
\label{IV.8.8.3}
The essential content of \sref{IV.8.8.2} is expressed again by saying that if $S_0$ is quasi-compact and quasi-separated, the category of $S$-preschemes of finite presentation is determined up to equivalence by the data of the categories of $S_\lambda$-preschemes of finite presentation, the functors $\rho_{\mu\lambda}:X_\lambda\mapsto X_\lambda\times_{S_\lambda}S_\mu$ (for $\lambda\leqslant\mu$) between these categories and the transitivity isomorphisms $\rho_{\mu\lambda}\circ\rho_{\nu\mu}\xrightarrow{\sim}\rho_{\nu\lambda}$.
In the form thus stated, one can say that the data of an $S_\lambda$-prescheme of finite presentation $X$ is equivalent ``functorially'' to the data of an $S_\mu$-prescheme of finite presentation $X'_\mu$; if an $S_\mu$-prescheme $X''_\mu$ is such that $X$ is $S$-isomorphic to $X''_\mu\times_{S_\mu}S$, there exists $\nu$ such that $\lambda\leqslant\nu$ and $X'_\lambda\times_{S_\lambda}S_\nu$ and $X''_\mu\times_{S_\mu}S_\nu$ are $S_\nu$-isomorphic.
The fact that $L$ is filtered implies in addition that if one gives oneself a finite \emph{family} $(X^{(i)})_{i\in I}$ of $S$-preschemes of finite presentation, and a finite family $(f^{(ij)})_{(i,j)\in J}$ of $S$-morphisms between these preschemes ($f^{(ij)}$ being therefore of the form $X^{(\sigma(j))}\to X^{(\tau(j))}$ where $\sigma$ and $\tau$ are two maps from $J$ to $I$), then there is an index $\lambda\in L$, a family $(X^{(i)}_\lambda)_{i\in I}$ of $S_\lambda$-preschemes of finite presentation and a family $(f^{(ij)}_\lambda)_{(i,j)\in J}$ of $S_\lambda$-morphisms such that $X^{(i)}$ is identified with $X^{(i)}_\lambda\times_{S_\lambda}S$ and $f^{(ij)}$ with $f^{(ij)}_\lambda\times 1_S$; furthermore, if one has a relation $f^{(ij)}=f^{(ik)}$, one can suppose $\lambda$ chosen such that $f^{(ij)}_\lambda=f^{(ik)}_\lambda$.

Consider in particular three $S$-preschemes of finite presentation $X$, $Y$, $Z$ and two $S$-morphisms $f:X\to Y$, $g:Y\to Z$, so that the product $X\times_Z Y$ relative to these morphisms is also an $S$-prescheme of finite presentation \sref[I]{I.1.6.2}.
It then results from the preceding and from \sref[I]{I.3.3.11} that one can write $X\times_Z Y=(X_\lambda\times_{Z_\lambda}Y_\lambda)\times_{S_\lambda}S$ for a convenient $\lambda$, $X_\lambda$, $Y_\lambda$, $Z_\lambda$ being $S_\lambda$-preschemes of finite presentation; one can therefore say that the determination of $S$-preschemes of finite presentation by the data of $S_\lambda$-preschemes of finite presentation is ``\emph{compatible with fibre products}''.
We have seen on the other hand \sref{IV.8.6.3} that if $g$ is an immersion, we can suppose the same for $g_\lambda$; identifying $Y$ (resp.\ $Y_\lambda$) with a sub-prescheme of $Z$ (resp.\ $Z_\lambda$), we see that one can write, for a convenient $\lambda$, $f^{-1}(Y)=f_\lambda^{-1}(Y_\lambda)\times_{S_\lambda}S$ \sref[I]{I.4.4.1};
\oldpage[IV-3]{33}
there is therefore also ``\emph{compatibility with the formation of inverse images of sub-preschemes}''.
More particularly, if $f_1$, $f_2$ are two $S$-morphisms from $X$ to $Y$, one calls the \emph{kernel} of this pair of morphisms the inverse image $N$ of the diagonal sub-prescheme of $Y\times_S Y$ by the $S$-morphism $(f_1,f_2)_S$; one deduces from what precedes that one will have for a convenient $\lambda$, $N=N_\lambda\times_{S_\lambda}S$, where $N_\lambda$ is the kernel of the couple of $S_\lambda$-morphisms $(f_{1\lambda},f_{2\lambda})$ of $X_\lambda$ into $Y_\lambda$.
These remarks extend to finite products and to ``kernels'' of arbitrary finite systems of $S$-morphisms from an $S$-prescheme of finite presentation into
\oldpage[IV-3]{34}
another; one concludes that in a general fashion the formation of $S$-preschemes of finite presentation by the data of $S_\lambda$-preschemes of finite presentation is ``\emph{compatible with finite projective limits}'', such a limit being by definition the kernel of a finite system of morphisms from one $S$-prescheme into a product of $S$-preschemes \sref[T]{T.1.8}.

We shall allow ourselves in what follows to freely make the translations implied by the preceding properties (as well as by \sref{IV.8.3.11}, \sref{IV.8.5.2}, and \sref{IV.8.6.3}) without always striving to justify them as above.
For example, the data of a ``group prescheme'' $G$ of finite presentation over $S$ is equivalent to the data of a group prescheme on $S_\lambda$ for a large enough $\lambda$; to write indeed the associativity condition for the $S$-morphism ``composition law'' $G\times_S G\to G$ of the group prescheme amounts to writing that the kernel of the two $S$-morphisms from $G\times_S G\times_S G\to G$ is equal to $G\times_S G\times_S G$ \sref[II]{II.8.3.9}, and the other conditions intervening in the definition of a group prescheme are interpreted similarly.
\end{env}

\subsection{First applications to the elimination of Noetherian hypotheses}
\label{subsection:IV.8.9}

\begin{proposition}[8.9.1]
\label{IV.8.9.1}
Let $A$ be a ring, $X$ an $A$-prescheme.
\begin{enumerate}[label=(\roman*)]
\item The following conditions are equivalent:
\begin{enumerate}[label=\emph{\alph*)}]
\item $X$ is of finite presentation over $A$.
\item There exist a Noetherian ring $A_0$, a prescheme $X_0$ of finite type over $A_0$, a ring homomorphism $A_0\to A$, and an $A$-isomorphism $X_0\otimes_{A_0}A\xrightarrow{\sim}X$.
\item There exist a subring $A_0$ of $A$, which is a $\mathbf{Z}$-algebra of finite type, a prescheme $X_0$ of finite type over $A_0$, and an $A$-isomorphism $X_0\otimes_{A_0}A\xrightarrow{\sim}X$.
\end{enumerate}
\item If $\sh{F}$ is a quasi-coherent $\sh{O}_X$-Module of finite presentation, one can suppose the subring $A_0$ such that there exists a coherent $\sh{O}_{X_0}$-Module $\sh{F}_0$ such that $\sh{F}$ is isomorphic to $\sh{F}_0\otimes_{A_0}A$; $\operatorname{Supp}(\sh{F})$ is constructible and closed in $X$, and there is a closed sub-prescheme $Z$ of $X$, having $\operatorname{Supp}(\sh{F})$ as underlying space, such that the canonical immersion $Z\to X$ be of finite presentation.
\item If $Y$ is a second $A$-prescheme of finite presentation, and $f:X\to Y$ an $A$-morphism, one can suppose the subring $A_0$ of $A$ such that there exists a prescheme $Y_0$ of finite type over $A_0$, an $A$-isomorphism $Y_0\otimes_{A_0}A\xrightarrow{\sim}Y$, and an $A_0$-morphism $f_0:X_0\to Y_0$ (necessarily of finite type) such that $f$ is identified with $f_0\otimes 1_A$.
\end{enumerate}
\end{proposition}

\begin{proof}
\begin{enumerate}[label=(\roman*)]
\item As $A$ is the inductive limit of its subrings which are of finite type over $\mathbf{Z}$, \emph{a)} implies \emph{c)} by virtue of \sref{IV.8.8.2},~(ii); \emph{c)} implies \emph{b)} since a $\mathbf{Z}$-algebra of finite type is a Noetherian ring; finally, if $A_0$ is Noetherian, an $A_0$-prescheme of finite type is of finite presentation \sref[I]{I.1.6.1}, hence \emph{b)} implies \emph{a)} by virtue of \sref[I]{I.1.6.2},~(iii).
\end{enumerate}
The assertion (ii) follows immediately from \sref{IV.8.5.2},~(ii), \sref{IV.8.3.11}, and \sref[I]{I.1.8.2}, and the assertion (iii) from \sref{IV.8.8.2},~(i).
\end{proof}

\begin{env}[8.9.2]
\label{IV.8.9.2}
The proposition \sref{IV.8.9.1} and the results of \sref{IV.8.5.2} and \sref{IV.8.8.2} allow one to extend in numerous cases to morphisms of finite presentation $X\to Y$ the results proved by Noetherian techniques by supposing $Y$ \emph{locally Noetherian}.
\oldpage[IV-3]{35}
We shall see numerous examples in what follows and we shall limit ourselves here to giving a few results of this type.
\end{env}

\begin{proposition}[8.9.3]
\label{IV.8.9.3}
Let $A$ be a ring, $M$ an $A$-module of finite presentation.
Then every surjective $A$-endomorphism of $M$ is bijective.
\end{proposition}

\begin{proof}
Indeed, consider $A$ as an inductive limit of its sub-$\mathbf{Z}$-algebras of finite type.
It follows from \sref{IV.8.5.2},~(ii) that there exists one of these sub-algebras $A_0$ and an $A_0$-module $M_0$ of finite presentation such that $M$ is $A$-isomorphic to $M_0\otimes_{A_0}A$; furthermore, if $f:M\to M$ is a surjective $A$-endomorphism, one can suppose \sref{IV.8.5.2},~(i) that there exists an $A_0$-endomorphism $f_0:M_0\to M_0$ such that $f=f_0\otimes 1_A$; finally, by \sref{IV.8.5.7}, one can suppose that $f_0$ is \emph{surjective}.
But as $A_0$ is Noetherian and $M_0$ an $A_0$-module of finite type, $M_0$ is a Noetherian module, hence (Bourbaki, \emph{Alg.}, chap.~VIII, \textsection\,2, n\textsuperscript{o}\,2, lemma~3) $f_0$ is bijective, and it is the same for $f$ by base change.
\end{proof}

\begin{proposition}[8.9.4]
\label{IV.8.9.4}
(``generic flatness theorem''). --- Let $Y$ be an integral prescheme, $u:X\to Y$ a morphism of finite type and locally of finite presentation, $\sh{F}$ an $\sh{O}_X$-Module quasi-coherent of finite presentation.
Then there exists an open non-empty $U$ of $Y$ such that $\sh{F}|u^{-1}(U)$ be flat over $U$.
\end{proposition}

\begin{proof}
The reasoning from the beginning of \sref{IV.6.9.1} leads to proving the

\begin{lemma}[8.9.4.1]
\label{IV.8.9.4.1}
Let $A$ be an integral ring, $B$ an $A$-algebra of finite presentation, $M$ a $B$-module of finite presentation.
Then there exists $f\neq 0$ in $A$ such that $M_f$ is a free $A_f$-module.
\end{lemma}

\begin{proof}
Indeed, by \sref{IV.8.9.1} there is a finitely generated $\mathbf{Z}$-subalgebra $A_0$ of $A$, a finite-type $A_0$-algebra $B_0$, and a finite $B_0$-module $M_0$ such that $B$ is $A$-isomorphic to $B_0\otimes_{A_0}A$ and $M$ is $B$-isomorphic to $M_0\otimes_{A_0}A$; by \sref{IV.6.9.2}, there exists $f_0\neq 0$ in $A_0$ such that $(M_0)_{f_0}$ is a free $(A_0)_{f_0}$-module.
As $M_{f_0}=(M_0)_{f_0}\otimes_{A_0}A$ and $A_{f_0}=(A_0)_{f_0}\otimes_{A_0}A$, $M_{f_0}$ is a free $A_{f_0}$-module.
\end{proof}
\end{proof}

\begin{corollary}[8.9.5]
\label{IV.8.9.5}
Let $S$ be a quasi-compact prescheme and quasi-separated, $u:X\to S$ a morphism of finite presentation, $\sh{F}$ an $\sh{O}_X$-Module quasi-coherent of finite presentation.
Then there exists a partition $(S_i)_{1\leqslant i\leqslant n}$ of $S$ into finitely many locally closed subsets of $S$ such that, for $1\leqslant i\leqslant n$, there is a subprescheme $S'_i$ of $S$, with underlying space $S_i$, which is of finite presentation over $S$, and such that, on putting $X_i=X\times_S S'_i$, the $\sh{O}_{X_i}$-Module $\sh{F}_i=\sh{F}\otimes_{\sh{O}_X}\sh{O}_{X_i}$ is flat over $S'_i$.
\end{corollary}

\begin{proof}
Consider a finite covering $(U_j)_{1\leqslant j\leqslant m}$ of $S$ by affine open subsets and define, inductively, $T_1=U_1$ and
\[
T_j=U_j\setminus\bigcup_{h<j}(U_j\cap U_h)\qquad(j\geqslant 2).
\]
Each $T_j$ is closed in the affine open $U_j$, and the $T_j$ are pairwise disjoint.
Moreover, the $U_j\cap U_h$ are quasi-compact since $S$ is quasi-separated and hence, since $S$ is also quasi-compact, are constructible in $S$ \sref[I]{I.1.8.1}; consequently each $T_j$ is constructible as well.
It will therefore suffice to prove the corollary for a suitable subprescheme $T'_j$ of $S$ having $T_j$ as underlying space and of finite presentation over $S$, and for the morphism and the Module deduced
\oldpage[IV-3]{36}
respectively from $u$ and $\sh{F}$ by the base change $T'_j\to S$.
For this, note that if $U_j$ is given the prescheme structure induced by that of $S$, the open immersion $U_j\to S$ is quasi-compact because $S$ is quasi-separated \sref[I]{I.1.2.7}, and hence is of finite presentation \sref[I]{I.1.6.2},~(i).
It therefore suffices that $T'_j$ be of finite presentation over $U_j$; in other words, we may already reduce to the case where $U_j=S$ and $T_j=T$ is closed and constructible in $S$.

Write $S=\operatorname{Spec}(A)$, and regard $A$ as the inductive limit of its finitely generated $\mathbf{Z}$-subalgebras, so that $S=\varprojlim S_\lambda$, where the $S_\lambda$ are affine and Noetherian.
By \sref{IV.8.3.11}, there exist $\lambda$ and a closed (necessarily constructible) subset $T_\lambda$ of $S_\lambda$ such that $T=u_\lambda^{-1}(T_\lambda)$, where $u_\lambda:S\to S_\lambda$ is the canonical morphism.
Give $T_\lambda$ a subprescheme structure and take $T'=T_\lambda\times_{S_\lambda}S$.
Since the canonical immersion $T_\lambda\to S_\lambda$ is of finite presentation \sref[I]{I.1.6.1}, so is the immersion $T'\to S$.
As $T'$ is affine, we see finally that it is enough to treat the case where $T'=S$ is affine.

Then, with the notation of \sref{IV.8.9.1}, there exist $\lambda$, a morphism of finite type $u_\lambda:X_\lambda\to S_\lambda$, and a coherent $\sh{O}_{X_\lambda}$-Module $\sh{F}_\lambda$ such that $X$ is isomorphic to $X_\lambda\times_{S_\lambda}S$ and $\sh{F}$ to $\sh{F}_\lambda\otimes_{\sh{O}_{X_\lambda}}\sh{O}_X$.
Applying \sref{IV.6.9.3} to $S_\lambda$, $X_\lambda$, and $\sh{F}_\lambda$, we obtain subpreschemes $S'_{\lambda i}$ of $S_\lambda$ whose underlying sets form a partition of $S_\lambda$ and which are such that, on putting $X_{\lambda i}=X_\lambda\times_{S_\lambda}S'_{\lambda i}$, the Module
\[
\sh{F}_{\lambda i}=\sh{F}_\lambda\otimes_{\sh{O}_{X_\lambda}}\sh{O}_{X_{\lambda i}}
\]
is flat over $S'_{\lambda i}$.
The subpreschemes $S'_i=S'_{\lambda i}\times_{S_\lambda}S$ of $S$ then have the required properties by \sref{IV.2.1.4}.
\end{proof}

\subsection{Properties of permanence of morphisms by passage to the projective limit}
\label{subsection:IV.8.10}

In this section, we preserve the general hypotheses and the notations of \sref{IV.8.8.1}.

\begin{proposition}[8.10.1]
\label{IV.8.10.1}
If there exists $\lambda$ such that, for $\mu\geqslant\lambda$, $f_\mu$ is an open morphism (resp.\ universally open), then $f$ is an open morphism (resp.\ universally open).
\end{proposition}

\begin{proof}
As $X=X_\lambda\times_{Y_\lambda}Y$, the assertion relative to universally open morphisms is a particular case of \sref{IV.2.4.3},~(iii).
Let us therefore suppose $f_\mu$ open for $\mu\geqslant\lambda$; it suffices to see that for every open \emph{quasi-compact} $U$ of $X$, $f(U)$ is open in $Y$.
Now, there exists $\mu\geqslant\lambda$ such that $U=v_\mu^{-1}(U_\mu)$, where $U_\mu$ is an open quasi-compact in $X_\mu$ \sref{IV.8.2.11}; one has then $f(v_\mu^{-1}(U_\mu))=w_\mu^{-1}(f_\mu(U_\mu))$ \sref[I]{I.3.4.8}, hence $f(U)$ is open in $Y$.
\end{proof}

\begin{corollary}[8.10.2]
\label{IV.8.10.2}
Let $f:X\to Y$ be a morphism.
For $f$ to be universally open, it suffices that, for every integer $n>0$, if one sets $Y_n=Y\otimes_\mathbf{Z}\mathbf{Z}[T_1,\ldots,T_n](=\mathbf{V}_Y^n)$, and $X_n=X\times_Y Y_n$, $f_n=f\times 1_{Y_n}:X_n\to Y_n$, the projection $f_n$ be an open morphism.
\end{corollary}

\begin{proof}
To prove that $f$ is universally open, it suffices to prove that it remains so for every affine $U$ of $Y$; by hypothesis, if $U_n=U\otimes_\mathbf{Z}\mathbf{Z}[T_1,\ldots,T_n]$ is the inverse image of $U$ in $Y_n$, the morphism $f_n^{-1}(U_n)\to U_n$, restriction of $f_n$, is open, and one sees that one can restrict to the case where $Y=\operatorname{Spec}(A)$ is affine.
Furthermore, it suffices evidently to show that for every morphism $Y'\to Y$, where $Y'=\operatorname{Spec}(A')$ is also \emph{affine}, $f'=f_{|Y'}$ is open.
Let us suppose first that $A'$ is an $A$-algebra of finite type, hence a \emph{quotient} of a polynomial algebra $A[T_1,\ldots,T_n]$; then $Y'$ is a closed sub-prescheme of $Y_n$ and $f'$ the restriction
\oldpage[IV-3]{37}
of $f_n$ to $f_n^{-1}(Y')$; but for every open $V$ of $X_n$, one has $f_n(V\cap f_n^{-1}(Y'))=f_n(V)\cap Y'$, and as by hypothesis $f_n(V)$ is open in $Y_n$, this shows that $f'$ is also an open morphism.
When $A'$ is arbitrary, it can be considered as an inductive limit of its sub-$A$-algebras of finite type $A'_\lambda$, and the fact that $f'$ is an open morphism follows from the preceding and from \sref{IV.8.10.1}.
\end{proof}

\begin{proposition}[8.10.3]
\label{IV.8.10.3}
Suppose that there exists $\alpha$ such that: $1^\circ$ $S_\alpha$ is quasi-compact; $2^\circ$ the morphisms $X_\alpha\to S_\alpha$, $Y_\alpha\to S_\alpha$ are quasi-compact and the morphism $Y_\alpha\to S_\alpha$ quasi-separated; $3^\circ$ for $\alpha\leqslant\lambda\leqslant\mu$, the morphisms $u_{\lambda\mu}:S_\mu\to S_\lambda$ are flat; $4^\circ$ $f_\alpha(X_\alpha)$ is constructible in $Y_\alpha$.
Then, for $f$ to be dominant, it is necessary and sufficient that there exist $\lambda\geqslant\alpha$ such that $f_\lambda$ be dominant.
\end{proposition}

\begin{proof}
The hypotheses imply that $Y_\alpha$ is quasi-compact and that the morphism $f_\alpha$ is quasi-compact \sref[I]{I.1.2.4}; hence $f_\alpha(X_\alpha)=Z_\alpha$ is pro-constructible \sref[I]{I.1.9.5},~(vii) in $Y_\alpha$.
If one sets $Z_\lambda=w_{\alpha\lambda}^{-1}(Z_\alpha)$ for $\lambda\geqslant\alpha$ and $Z=w_\alpha^{-1}(Z_\alpha)$, we have $Z_\lambda=f_\lambda(X_\lambda)$ and $Z=f(X)$ \sref[I]{I.3.4.8}.
As $Z_\lambda$ is pro-constructible in $Y_\lambda$ \sref{IV.8.3.11}, the canonical map $\varinjlim\mathfrak{C}(S_\lambda)\to\mathfrak{C}(S)$ is injective, which proves the theorem in the case (vi).
It suffices then to apply \sref{IV.8.3.12} by replacing $S_\lambda$, $Z'_\lambda$, and $Z''_\lambda$ by $Y_\lambda$, $Y_\lambda$, and $Z_\lambda$ respectively.
\end{proof}

\begin{proposition}[8.10.4]
\label{IV.8.10.4}
Suppose that there exists $\alpha$ such that $Y_\alpha$ is quasi-compact and $f_\alpha$ is of finite type and quasi-separated.
For the morphism $f$ to be separated, it is necessary and sufficient that there exist $\lambda\geqslant\alpha$ such that $f_\lambda$ be separated.
\end{proposition}

\begin{proof}
The question being local on $Y_\alpha$ (since $Y_\alpha$ is quasi-compact and $L$ filtered), one can restrict to the case where $Y_\alpha$ is affine, hence quasi-separated, and the hypothesis implies that $X_\alpha$ (hence the $X_\lambda$ and $X$) are quasi-compact and quasi-separated.
Set $X'_\lambda=X_\lambda\times_{Y_\lambda}X_\lambda$ for $\lambda\geqslant\alpha$ and $X'=X\times_Y X$; the first projection $X'_\alpha\to X_\alpha$ is quasi-compact and quasi-separated by hypothesis \sref[I]{I.1.2.3},~(iii), hence $X'_\lambda$ is quasi-compact and quasi-separated.
Note now that if we denote by $\Delta_\lambda$ (resp.\ $\Delta$) the diagonal of $X_\lambda\times_{Y_\lambda}X_\lambda$ (resp.\ $X\times_Y X$), it results from \sref[I]{I.5.3.4} and \sref{IV.3.4.8} that $\Delta_\lambda$ is the inverse image of $\Delta_\alpha$ by the map $v'_{\lambda\alpha}:X'_\lambda\to X'_\alpha$ (resp.\ $v'_\lambda:X'\to X'_\alpha$).
On the other hand, $\Delta_\alpha$ is \emph{constructible} in $X'_\alpha$: indeed, since $f_\alpha$ is quasi-separated, the diagonal immersion $X_\alpha\to X'_\alpha$ is quasi-compact, and locally of finite presentation since $f_\alpha$ is of finite type \sref[I]{I.1.4.3} and \sref[I]{I.1.5.4},~(iii); it results then from \sref[I]{I.1.8.4.1} that $\Delta_\lambda$ is locally constructible, hence constructible since $X'_\lambda$ is quasi-compact.
One can now apply \sref{IV.8.3.12} by replacing $S_\lambda$ and $Z_\lambda$ by $X'_\lambda$ and $\Delta_\lambda$ respectively.
\end{proof}

\begin{theorem}[8.10.5]
\label{IV.8.10.5}
Suppose $S_0$ quasi-compact, $X_\alpha$ and $Y_\alpha$ of finite presentation over $S_\alpha$, and let $f_\alpha:X_\alpha\to Y_\alpha$ be an $S_\alpha$-morphism.
Consider, for a morphism, the property of being:
\begin{enumerate}[label=(\roman*)]
\item an isomorphism;
\item[\rm(i\,\emph{bis})] a monomorphism;
\item an immersion;
\item an open immersion;
\item a closed immersion;
\item separated;
\item surjective;
\oldpage[IV-3]{38}
\item radicial;
\item affine;
\item quasi-affine;
\item finite;
\item quasi-finite;
\item proper.
\end{enumerate}
Then, if $\mathbf{P}$ denotes one of the preceding properties, for $f$ to possess the property $\mathbf{P}$, it is necessary and sufficient that there exist $\lambda\geqslant\alpha$ such that $f_\lambda$ possess the property $\mathbf{P}$ (in which case $f_\mu$ also possesses it for $\mu\geqslant\lambda$).

If $S_0$ is furthermore supposed quasi-separated, the same conclusion holds when $\mathbf{P}$ is the property of being:
\begin{enumerate}[label=(\roman*)]
\setcounter{enumi}{12}
\item projective;
\item quasi-projective.
\end{enumerate}
\end{theorem}

\begin{proof}
The case where $\mathbf{P}$ is one of the properties (i) or (v) was not inserted in the statement for the record; in these cases, the theorem follows from what has been proved respectively in \sref{IV.8.8.2.4} and \sref{IV.8.10.4}.
Furthermore, taking into account \sref[I]{I.5.4.1} and \sref{I.5.3.4}, (v) also results from (iv).
The case (i\,\emph{bis}) follows immediately from (i), by using \sref[I]{I.5.3.8} and by noting (as we have already used in \sref{IV.8.10.4}) that the diagonal $\Delta_f$ is deduced from $\Delta_{f_\lambda}$ by the base change $S\to S_\lambda$.

We will note on the other hand that (vi), (vii), and (xi) are in fact conditions on the \emph{fibres} $f^{-1}(y)$ of the morphisms considered, taking into account the transitivity of the fibres by base change \sref[I]{I.3.6.4}: condition (vi) means indeed that all the fibres must be non-empty, condition (vii) that they must be radicial \sref[I]{I.3.5.8}, and condition (xi) that they must be finite \sref[II]{II.6.2.2} and \sref{IV.6.2.3} and \sref[I]{I.6.4.4}, taking into account that $f$ and the $f_\lambda$ are morphisms of finite type by \sref[I]{I.1.5.4},~(v).
The theorem in these three cases will be yet another consequence of a general result on this type of properties, which will be established in \sref{IV.9.3.3}; we therefore postpone until then the proof of the theorem in the case (xi) (of course, the reader will be able to verify that, except for n\textsuperscript{os}~{\hyperlink{ega4.c04.subsec.8-11}{8.11}} and {\hyperlink{ega4.c04.subsec.8-12}{8.12}}, we will make no use of the theorem in these cases until \sref{IV.9.3.3} and that \sref{IV.8.11} and \sref{IV.8.12} will not be used before \sref{IV.9.3.3}).

For the cases that remain to be proved, one can restrict to proving that the condition of the statement is necessary, all the properties $\mathbf{P}$ considered being invariant by base change (see chap.~I\textsuperscript{er} and II for the references concerning each of these properties).
One can furthermore suppose that $S_\alpha=S_0$ and that $Y_\alpha=S_\alpha$, hence $Y_\lambda=S_\lambda$ for every $\lambda\geqslant\alpha$.
Finally, the properties (i) through (xii) are local on $S_0$, hence as $S_0$ is a finite union of open affines and $L$ filtered, one can restrict to proving them for $S_0=\operatorname{Spec}(A_0)$ \emph{affine} (hence quasi-separated).
We shall denote by $A_\lambda$ (resp.\ $A$) the ring of $S_\lambda$ (resp.\ $S$).

\emph{Cases} (ii), (iii), (iv): Suppose that $f$ is an immersion (resp.\ an open immersion, a closed immersion), and let $X'$ be the sub-prescheme (resp.\ induced on an open, resp.\ closed) of $S$ associated to $f$, which is therefore an $S$-prescheme of finite presentation.
\oldpage[IV-3]{39}
By virtue of \sref{IV.8.6.3}, there then exists $\lambda\geqslant\alpha$ and a sub-prescheme $X'_\lambda$ of $S_\lambda$, of finite presentation over $S_\lambda$, such that $X'$ is isomorphic to $X'_\lambda\times_{S_\lambda}S$.
For every $\mu\geqslant\lambda$, $X'_\mu=X'_\lambda\times_{S_\lambda}S_\mu$ is therefore a sub-prescheme (resp.\ induced on an open, resp.\ closed) of $S_\mu$, and it results therefore from \sref{IV.8.8.2.4} and \sref{IV.8.8.2.5} that there exists $\mu\geqslant\lambda$ and an isomorphism $g_\mu:X_\mu\to X'_\mu$ such that $g=g_\mu\times 1_S$ is the isomorphism $X\to X'$ associated to $f$; whence the conclusion.

\emph{Cases} (vi) \emph{and} (vii): We know \sref[I]{I.1.8.1} that $Z_\alpha=f_\alpha(X_\alpha)$ is constructible in $S_\alpha$; if we set $Z_\lambda=u_{\alpha\lambda}^{-1}(Z_\alpha)$ for $\lambda\geqslant\alpha$ and $Z=u_\alpha^{-1}(Z_\alpha)$, we have $Z_\lambda=f_\lambda(X_\lambda)$ and $Z=f(X)$ \sref[I]{I.3.4.8}.
As $Z_\alpha$ is pro-constructible in $S_\alpha$, and by virtue of \sref{IV.8.3.11} the canonical map $\varinjlim\mathfrak{C}(S_\lambda)\to\mathfrak{C}(S)$ is injective, the relation $f(X)=S$ implies the existence of $\lambda\geqslant\alpha$ such that $f_\lambda(X_\lambda)=S_\lambda$, which proves the theorem in the case (vi).
To prove it in case~(vii), note that the structural morphism $X_\alpha\times_{S_\alpha}X_\alpha\to S_\alpha$ is of finite presentation, since the first projection $X_\alpha\times_{S_\alpha}X_\alpha\to X_\alpha$ is so \sref{IV.1.6.2}. It is therefore enough, by \sref[I]{I.1.8.7.1}, to apply case~(vi) to the diagonal morphism
\[
  \Delta_{f_\alpha}:X_\alpha\longrightarrow X_\alpha\times_{S_\alpha}X_\alpha,
\]
noting that $\Delta_{f_\lambda}=\Delta_{f_\alpha}\times1_{S_\lambda}$ and $\Delta_f=\Delta_{f_\alpha}\times1_S$ \sref[I]{I.5.3.4} and \sref[I]{I.3.3.11}.

\emph{Cases} (viii) \emph{and} (ix): As $S=\operatorname{Spec}(A)$ is affine, to say that $f:X\to S$ is affine (resp.\ quasi-affine) means that there exists an integer $r$ and a closed immersion (resp.\ an immersion) $j:X\to\mathbf{V}^r_S=\operatorname{Spec}(A[T_1,\ldots,T_r])$ of $S$-preschemes, since $f$ is of finite type and $S$ quasi-compact \sref[II]{II.5.1.9}.
As $\mathbf{V}^r_S=\mathbf{V}^r_{S_\lambda}\times_{S_\lambda}S$, and $\mathbf{V}^r_{S_\lambda}$ is an $S_0$-prescheme of finite presentation, it results from \sref{IV.8.8.2},~(i) applied to the $S$-morphism $j$, that there exists $\lambda$ and an $S_\lambda$-morphism $j_\lambda:X_\lambda\to\mathbf{V}^r_{S_\lambda}$ such that $j=j_\lambda\times 1_S$; applying (iv) (resp.\ (ii)) to $j$, one deduces that there exists $\mu\geqslant\lambda$ such that $j_\mu$ is a closed immersion (resp.\ an immersion); hence $f_\mu$ is affine (resp.\ quasi-affine).

\emph{Case} (x): By hypothesis, $X=\operatorname{Spec}(B)$, where $B$ is an $A$-algebra which is an $A$-module of finite presentation \sref[I]{I.1.4.7}; it results therefore from \sref{IV.8.5.2},~(ii) that there is $\lambda$ and an $A_\lambda$-module of finite presentation $B_\lambda$ such that the $A$-module $B$ is isomorphic to $B_\lambda\otimes_{A_\lambda}A$.
The structure of $A$-algebra of $B$ is defined by an $A$-homomorphism $m:B\otimes_A B\to B$; as $B\otimes_A B=(B_\lambda\otimes_{A_\lambda}B_\lambda)\otimes_A A$, there exists by \sref{IV.8.5.2},~(i) $\mu\geqslant\lambda$ and an $A_\mu$-homomorphism $m_\mu:B_\mu\otimes_{A_\mu}B_\mu\to B_\mu$ such that $m=m_\mu\otimes 1$.
Considering the usual diagrams expressing the associativity and the commutativity of $m$, one sees by applying again \sref{IV.8.5.2},~(i) that one can suppose $\nu$ taken large enough so that $m_\nu$ defines on $B_\nu$ an associative and commutative multiplication.
In the same way one can suppose $\nu$ large enough so that $B_\nu$ thus defined admits a unit element.
If $X'_\nu=\operatorname{Spec}(B_\nu)$, it is clear then that $X$ is $S$-isomorphic to $X'_\nu\times_{S_\nu}S$, hence, by virtue of (i), there exists $\rho\geqslant\nu$ such that $X_\rho$ and $X'_\rho$ are $S_\rho$-isomorphic, which completes the proof in this case.

To prove the theorem in the case (xii), we shall first prove the following proposition:
\end{proof}

\begin{proposition}[8.10.5.1]
\label{IV.8.10.5.1}
(Chow's lemma for morphisms of finite presentation). --- Let $\mathrm{A}$ be a ring, let $X$ and $Y$ be two $\mathrm{A}$-preschemes of finite presentation, and let $f:X\to Y$ be a separated $\mathrm{A}$-morphism. Then there exist two $\mathrm{A}$-preschemes $X'$ and $P$ of finite presentation, and $\mathrm{A}$-morphisms $p:P\to Y$, $j:X'\to P$, and $g:X'\to X$, such that the diagram
\oldpage[IV-3]{40}
\[
\begin{tikzcd}[column sep=large,row sep=large]
X' \arrow[r,"j"] \arrow[d,"g"'] & P \arrow[d,"p"] \\
X \arrow[r,"f"'] & Y
\end{tikzcd}
\]
is commutative, and such that: $1^\circ$ $p$ is projective; $2^\circ$ $g$ is projective and surjective; $3^\circ$ $j$ is an open immersion.
\end{proposition}

\begin{proof}
Choose $\mathrm{A}_0\subset\mathrm{A}$, $X_0$, $Y_0$, and $f_0$ as in \sref{IV.8.9.1}, with $Y_0$ Noetherian and $f_0$ of finite type; by \sref{IV.8.10.4}, one may moreover suppose that $f_0$ is separated. Chow's lemma \sref[II]{II.5.6.1} then gives morphisms of finite type $p_0:P_0\to Y_0$, $g_0:X'_0\to X_0$, and $j_0:X'_0\to P_0$ such that
\[
\begin{tikzcd}[column sep=large,row sep=large]
X'_0 \arrow[r,"j_0"] \arrow[d,"g_0"'] & P_0 \arrow[d,"p_0"] \\
X_0 \arrow[r,"f_0"'] & Y_0
\end{tikzcd}
\]
is commutative, $p_0$ is projective, $g_0$ is projective and surjective, and $j_0$ is an open immersion. The assertions follow from the invariance of these properties under base change \sref[II]{II.5.5.5},~(iii), \sref[I]{I.3.5.2}, and \sref{IV.4.3.2}.
\end{proof}

\begin{proof}[Continuation of the proof of \sref{IV.8.10.5}]
\emph{Case (xii).} Apply \sref{IV.8.10.5.1} to $f_0:X_0\to S_0$. We obtain a commutative diagram
\[
\begin{tikzcd}[column sep=large,row sep=large]
X'_0 \arrow[r,"j_0"] \arrow[d,"g_0"'] & P_0 \arrow[d,"p_0"] \\
X_0 \arrow[r,"f_0"'] & S_0
\end{tikzcd}
\]
where $p_0$ is projective, $g_0$ is projective and surjective, and $j_0$ is an open immersion. For each $\lambda$ one obtains the analogous diagram in which $p_\lambda=p_0\times1_{S_\lambda}$, $g_\lambda=g_0\times1_{S_\lambda}$, and $j_\lambda=j_0\times1_{S_\lambda}$ have the same respective properties, and similarly for the inverse-limit morphisms $p=p_0\times1_S$, $g=g_0\times1_S$, and $j=j_0\times1_S$. Since $g$ is proper \sref[II]{II.5.5.3}, so is $p\circ j=f\circ g$ \sref[II]{II.5.4.2}; and since $p$ is separated, $j$ is proper, hence a closed immersion. Apply case~(iv) to $j$, noting that $X'_0$ and $P_0$ are $S_0$-preschemes of finite presentation by \sref{IV.8.10.5.1} and \sref{IV.1.6.2}. There exists $\lambda$ such that $j_\lambda$ is a closed immersion, hence proper \sref[II]{II.5.4.2}. Thus
\[
  f_\lambda\circ g_\lambda=p_\lambda\circ j_\lambda
\]
is proper \sref[II]{II.5.5.3} and \sref[II]{II.5.4.2}. Since $g_\lambda$ is surjective and, by the hypothesis on $f$ and case~(v), $f_\lambda$ may be supposed separated, \sref[II]{II.5.4.3} shows that $f_\lambda$ is proper.

\emph{Cases (xiii) and (xiv).} By case~(xii) and \sref[II]{II.5.5.3}, which applies because the $S_\lambda$ are quasi-compact and quasi-separated by \sref[I]{I.1.7.19}, it is enough to consider the case where $f$ is quasi-projective. Suppose there is an invertible $\sh{O}_X$-Module $\sh{L}$ which is ample for $f$. Since $S_0$ is quasi-compact and quasi-separated, so is $X_0$ \sref[I]{I.1.2.3}; hence there are $\lambda$ and a quasi-coherent $\sh{O}_{X_\lambda}$-Module $\sh{L}_\lambda$ of finite presentation such that $\sh{L}=v_\lambda^*(\sh{L}_\lambda)$ \sref{IV.8.5.2},~(ii). By \sref{IV.8.5.5}, one may moreover suppose $\sh{L}_\lambda$ invertible. The assertion now follows from the more precise lemma below.

\begin{lemma}[8.10.5.2]
\label{IV.8.10.5.2}
\oldpage[IV-3]{41}
Suppose that $S_0$ is quasi-compact, and let $\sh{L}_\lambda$ be an invertible $\sh{O}_{X_\lambda}$-Module. For $\sh{L}$ to be ample for $f$ (resp.\ very ample for $f$), it is necessary and sufficient that there exist $\mu\geqslant\lambda$ such that $\sh{L}_\mu$ is ample for $f_\mu$ (resp.\ very ample for $f_\mu$).
\end{lemma}

\begin{proof}
The condition is plainly sufficient \sref[II]{II.4.4.10} and \sref[II]{II.4.6.13}. Conversely, since the $S_\lambda$ are quasi-compact and the $f_\lambda$ are of finite type, after replacing $\sh{L}$ by a suitable tensor power, one may reduce to the case where $\sh{L}$ is very ample \sref[II]{II.4.6.11}. The question is local on $S_0$ by \sref[II]{II.4.4.5} and the fact that $L$ is filtered, so one may suppose that $S_0$, and hence $S$, is affine. By \sref[II]{II.4.4.1},~(ii) and \sref[II]{II.4.1.2}, there is an $S$-immersion
\[
  j:X\longrightarrow\mathbf{P}^r_S=P
\]
such that $\sh{L}\simeq j^*(\sh{O}_P(1))$. By \sref{IV.8.8.2},~(i), case~(ii), and \sref[II]{II.4.1.3}, there exist $\mu\geqslant\lambda$ and an immersion $j_\mu:X_\mu\to\mathbf{P}^r_{S_\mu}=P_\mu$ such that $j=j_\mu\times1_S$. Then \sref[II]{II.4.1.3.2} and \sref{IV.8.5.2.5} give $\nu\geqslant\mu$ such that
\[
  \sh{L}_\nu\simeq j_\nu^*(\sh{O}_{P_\nu}(1)),
\]
which shows that $\sh{L}_\nu$ is very ample for $f_\nu$ \sref[II]{II.4.4.2}.
\end{proof}
\end{proof}

\subsection{Application to quasi-finite morphisms}
\label{subsection:IV.8.11}
\label{IV.8.11}

We propose in this section to prove the two following theorems:

\begin{theorem}[8.11.1]
\label{IV.8.11.1}
Let $f:X\to Y$ be a proper morphism, locally of finite presentation, and quasi-finite.
Then the morphism $f$ is finite.
\end{theorem}

\begin{theorem}[8.11.2]
\label{IV.8.11.2}
Let $f:X\to Y$ be a morphism locally of finite presentation, quasi-finite and separated.
Then the morphism $f$ is quasi-affine, and a fortiori quasi-projective.
\end{theorem}

\begin{remark}[8.11.3]
\label{IV.8.11.3}
We will see below that one can reduce, for the proof of \sref{IV.8.11.1} and \sref{IV.8.11.2}, to the case where $Y$ is locally Noetherian; one will note that in this case one will thus obtain another proof of the theorem of Chevalley \sref[III]{III.4.4.2}.
\end{remark}

\oldpage[IV-3]{41}

\begin{env}[8.11.4]
\label{IV.8.11.4}
The hypotheses and the conclusions of \sref{IV.8.11.1} and \sref{IV.8.11.2} are all local on $Y$.
Indeed, this follows from \sref[I]{I.1.6.1}, \sref[I]{I.1.2.6}, \sref[II]{II.5.1.1}, \sref[II]{II.5.4.1}, and \sref[II]{II.6.2.2}.
Therefore one can suppose $Y=\operatorname{Spec}(A)$ affine.
One knows that there then exists a subring $A_0$ of $A$, which is a $\bb{Z}$-algebra of finite type, and a morphism of finite type $f_0:X_0\to\operatorname{Spec}(A_0)$ such that $X$ identifies with $X_0\otimes_{A_0}A$ and $f$ with $f_0\times 1$ \sref{IV.8.9.1}.
Furthermore, $A$ can be considered as inductive limit of its subrings containing $A_0$ and which are $\bb{Z}$-algebras of finite type; using the method of \sref{IV.8.1.2},~$c$)) as well as \sref{IV.8.10.5},~(v),~(xi) and (xii)), one sees that it suffices to prove the theorems \sref{IV.8.11.1} and \sref{IV.8.11.2} for $f_0$.
Suppose henceforth $Y$ Noetherian; using now the method of \sref{IV.8.1.2},~$a$)) as well as \sref{IV.8.10.5},~(v),~(ix),~(x),~(xi) and (xii)) one can replace $Y$ by $\operatorname{Spec}(\sh{O}_y)$, where $y$ is an
\oldpage[IV-3]{42}
arbitrary point of $Y$, so one can finally suppose that $Y=\operatorname{Spec}(A)$, where $A$ is a \emph{local Noetherian ring}.
Let $\mathfrak{m}$ be the maximal ideal of $A$, and let $\hat{A}$ be the completion of $A$ for the $\mathfrak{m}$-adic topology. The ring $\hat{A}$ is Noetherian local and $\operatorname{Spec}(\hat{A})\to\operatorname{Spec}(A)$ is faithfully flat and quasi-compact \sref[0]{0.7.3.5}; applying \sref{IV.2.7.1},~(i), (vii), (xiv), (xv) and (xvi), we may therefore reduce to the case where $A$ is \emph{complete}.
It then follows from \sref[II]{II.6.2.6} that $X=X'\sqcup X''$, where $X'$ is a $Y$-scheme quasi-finite and $X''$ a $Y$-scheme quasi-finite such that $X''\cap f^{-1}(y)=\varnothing$.
\end{env}

Let us first place ourselves in the hypotheses of \sref{IV.8.11.1}; as $f$ is proper, $X'$, which is closed in $X$, is proper over $Y$ \sref[II]{II.5.4.10}), so $f(X'')$ is closed in $Y$; but $y$ is not contained in $f(X'')$, and moreover $f(X'')$ is adherent to every point of $Y$, so $f(X'')=\varnothing$, and consequently $X''$ is empty and $f$ is \emph{finite}.

Let us now place ourselves in the hypotheses of \sref{IV.8.11.2} and, restricting ourselves further (as one can do from what precedes) to the case where $Y=\operatorname{Spec}(A)$ is affine and Noetherian of finite dimension, let us moreover reason by induction on the dimension of $Y$.
Reducing as above to the case where $A$ is furthermore local and complete, one has $\dim(\sh{O}_y)=\dim(A)=\dim(Y)$ and for all $z\neq y$, $\dim(\sh{O}_z)<\dim(\sh{O}_y)$, so $\dim(Y-\{y\})<\dim(Y)$.
Now, by hypothesis one has $f(X'')\subset Y-\{y\}$ and the restriction of $f$ to $X''$ is evidently a quasi-finite and separated morphism; applying the hypothesis of induction to $Y-\{y\}$ and to $X''$, one sees that $X''$ is quasi-affine over $Y-\{y\}$; but the open subset $Y-\{y\}$ being quasi-affine over $Y$ since $Y$ is Noetherian, $X''$ is also quasi-affine over $Y$ \sref[II]{II.5.1.10},~(ii)); as moreover $X'$ is finite (and \emph{a fortiori} affine) over $Y$, $X$ is quasi-affine over $Y$ \sref[II]{II.4.6.17} and \sref[II]{II.5.1.2},~$c'$)).

\begin{proposition}[8.11.5]
\label{IV.8.11.5}
Let $f:X\to Y$ be a morphism of finite presentation.
The following properties are equivalent:
\begin{enumerate}[label=\emph{\alph*)}]
  \item $f$ is a closed immersion.
  \item $f$ is a proper monomorphism.
  \item $f$ is proper and for all $y\in Y$, $f^{-1}(y)$ is radicial and geometrically reduced over $k(y)$ (that is to say empty or $k(y)$-isomorphic to $\operatorname{Spec}(k(y))$).
\end{enumerate}
\end{proposition}

It is clear that \emph{a)} implies \emph{b)}.
To see that \emph{b)} implies \emph{c)}, let us note \sref[I]{I.3.3.12} that for all $y\in Y$, the morphism $f^{-1}(y)\to\operatorname{Spec}(k(y))$ deduced from $f$ by extension of the base is a monomorphism, hence injective, and consequently $f^{-1}(y)$ is empty or reduced to a point, and in every case affine.
Furthermore, if $A$ is the ring of $f^{-1}(y)$, the canonical homomorphism $A\otimes_k A\to A$ is bijective \sref[I]{I.5.3.8}.
This evidently implies that $A=k$, for otherwise there would exist an element $a\in A$ not in $k$ and the images of $a\otimes 1$ and of $1\otimes a$ in $A$ would both be equal to $a$, so $a\otimes 1=1\otimes a$ since $1$ and $a$ form a free system over $k$.

\oldpage[IV-3]{43}

It remains to prove that \emph{c)} implies \emph{a)}.
It follows first from \sref{IV.8.11.1} that $f$ is a \emph{finite} morphism, so $X=\operatorname{Spec}(\sh{A})$, where $\sh{A}$ is a finite $\sh{O}_Y$-Algebra.
It suffices therefore to prove that the canonical homomorphism $\sh{O}_Y\to\sh{A}$ is \emph{surjective} \sref[II]{II.1.4.10}), or equivalently that for all $y\in Y$, the homomorphism $\sh{O}_y\to\sh{A}_y$ is surjective.
But by hypothesis $f^{-1}(y)=\operatorname{Spec}(\sh{A}_y/\mathfrak{m}_y\sh{A}_y)$ \sref[II]{II.1.5.5}) is such that the corresponding homomorphism $k(y)=\sh{O}_y/\mathfrak{m}_y\to\sh{A}_y/\mathfrak{m}_y\sh{A}_y$ is bijective; as $\sh{A}_y$ is an $\sh{O}_y$-module of finite type, the lemma of Nakayama shows that $\sh{O}_y\to\sh{A}_y$ is surjective, which completes the proof.

\begin{remark}[8.11.5.1]
\label{IV.8.11.5.1}
One will note that the preceding reasoning proves that if $f$ is a \emph{monomorphism}, then, for all $y\in Y$, $f^{-1}(y)$ is empty or $k(y)$-isomorphic to $\operatorname{Spec}(k(y))$.
\end{remark}

\begin{proposition}[8.11.6]
\label{IV.8.11.6}
If a morphism $f:X\to Y$ of finite presentation is a universal homeomorphism, it is finite, surjective and radicial (the converse being true by \sref{IV.2.4.5},~(i))).
\end{proposition}

\begin{proof}
Indeed, $f$ being of finite type, universally closed, and separated by virtue of \sref{IV.2.4.4}), is proper by definition \sref[II]{II.5.4.1}).
As it is evidently quasi-finite \sref[II]{II.6.2.3}), it is finite by \sref{IV.8.11.1}.
One knows moreover that it is radicial \sref{IV.2.4.4}), and evidently surjective.
\end{proof}

\subsection{New proof and generalisation of the ``Main Theorem'' of Zariski}
\label{subsection:IV.8.12}
\label{IV.8.12}

\begin{lemma}[8.12.1]
\label{IV.8.12.1}
Let $f:X\to Y$ be a quasi-compact and quasi-separated morphism, $\sh{C}$ an $\sh{O}_Y$-Algebra quasi-coherent, $Z=\operatorname{Spec}(\sh{C})$, which is an affine $Y$-prescheme over $Y$.
Let $g:X\to Z$ be a $Y$-morphism, $\varphi=\sh{A}(g):\sh{C}\to f_*(\sh{O}_X)$ the corresponding $\sh{O}_Y$-homomorphism of $\sh{O}_Y$-Algebras \sref[II]{II.1.2.7}.
Suppose that $g$ is an immersion.
Then, for the closed image of $X$ by $g$ \sref[I]{I.9.5.3}) to equal $Z$, it is necessary and sufficient that $\varphi$ be injective; $g$ is then an open immersion.
\end{lemma}

\begin{proof}
The hypothesis implies that $f_*(\sh{O}_X)$ is an $\sh{O}_Y$-Algebra quasi-coherent \sref{IV.1.7.5}); furthermore, as the canonical morphism $h:Z\to Y$ is affine, hence quasi-compact and quasi-separated \sref{IV.1.2.2} and \sref{IV.1.1.2}), $g_*(\sh{O}_X)$ is an $\sh{O}_Z$-Algebra quasi-coherent \sref{IV.1.7.5}).
This being so, saying that the closed image of $X$ by $g$ equals $Z$ signifies \sref[I]{I.9.5.2}) that the canonical homomorphism $\theta:\sh{O}_Z\to g_*(\sh{O}_X)$ is injective.
But one has $h_*(\sh{O}_Z)=\sh{C}$ by definition of $Z$ \sref[II]{II.1.3.1}), and $h_*(g_*(\sh{O}_X))=f_*(\sh{O}_X)$.
As $Z$ is affine over $Y$, it amounts to the same to say that the homomorphism $\theta:\sh{O}_Z\to g_*(\sh{O}_X)$ is injective or that the corresponding homomorphism $\varphi=h_*(\theta):\sh{C}\to f_*(\sh{O}_X)$ is injective \sref[I]{I.1.3.9}).
The fact that $g$ is then an open immersion follows from \sref[I]{I.9.5.10}) and the hypothesis.
\end{proof}

\begin{lemma}[8.12.2]
\label{IV.8.12.2}
Let $Y$ be a prescheme quasi-compact and quasi-separated, $f:X\to Y$ a quasi-separated morphism of finite type, $\sh{C}$ an $\sh{O}_Y$-Algebra quasi-coherent, $Z=\operatorname{Spec}(\sh{C})$.
Let $g:X\to Z$ be a $Y$-morphism, $\varphi:\sh{C}\to f_*(\sh{O}_X)$ the corresponding $\sh{O}_Y$-homomorphism of $\sh{O}_Y$-Algebras.
Let $(\sh{C}_\lambda)$ be the filtered increasing family of sub-$\sh{O}_Y$-Algebras quasi-coherent of finite type of $\sh{C}$ (where $\sh{C}$ is the union \sref[I]{I.9.6.6} and \sref{IV.1.7.9})); set $Z_\lambda=\operatorname{Spec}(\sh{C}_\lambda)$ and let $g_\lambda$ be the morphism composed $X\xrightarrow{g}Z\to Z_\lambda$.
Then the following conditions are equivalent:
\begin{enumerate}[label=\emph{\alph*)}]
  \item $g$ is an immersion.
  \item There exists $\lambda$ such that $g_\lambda$ is an immersion.
\end{enumerate}
\oldpage[IV-3]{44}
Furthermore, when $g_\lambda$ is an immersion, the same is true of $g_\mu$ for $\mu\geqslant\lambda$.

It suffices to apply \sref[II]{II.3.8.4}) by replacing $\sh{L}$ by $\sh{O}_X$ and $\sh{S}$ by $\sh{C}[\mathrm{T}]$, and taking into account \sref[II]{II.3.1.7}).
\end{lemma}

\begin{proposition}[8.12.3]
\label{IV.8.12.3}
Let $Y$ be a prescheme quasi-compact and quasi-separated, $f:X\to Y$ a separated morphism of finite type.
Let $\sh{B}=f_*(\sh{O}_X)$, which is an $\sh{O}_Y$-Algebra quasi-coherent \sref[I]{I.9.2.2}); let $\sh{C}$ be the integral closure of $\sh{O}_Y$ in $\sh{B}$, which is an $\sh{O}_Y$-Algebra quasi-coherent \sref[II]{II.6.3.4}); set $Z=\operatorname{Spec}(\sh{C})$, and let $g:X\to Z$ be the $Y$-morphism corresponding to the injection $\sh{C}\to\sh{B}=f_*(\sh{O}_X)$, and for all $\lambda$, let $g_\lambda:X\to Z_\lambda=\operatorname{Spec}(\sh{C}_\lambda)$ be the $Y$-morphism corresponding to the injection $\sh{C}_\lambda\to\sh{B}=f_*(\sh{O}_X)$.
Then the following conditions are equivalent:
\begin{enumerate}[label=\emph{\alph*)}]
  \item There exists a factorisation of $f$ as
\[
X\xrightarrow{f'}Y'\xrightarrow{u}Y
\]
  \item[\rm a$'$)] There exists a factorisation $X\xrightarrow{f'}Y'\xrightarrow{u}Y$ of $f$, where $f'$ is an open immersion and $u$ a finite morphism.
  \item The morphism $g:X\to Z$ is an immersion.
  \item There exists $\lambda$ such that $g_\lambda:X\to Z_\lambda$ is an immersion.
\end{enumerate}
Furthermore, when $g_\lambda$ is an immersion, $g$ is an open immersion, $g(X)$ is dense in $Z$, and there exists $\lambda$ such that for $\mu\geqslant\lambda$, $g_\mu$ is an open immersion.
\end{proposition}

where $f'$ is an open immersion and $u$ a finite morphism.

As the homomorphism $\varphi:\sh{C}\to f_*(\sh{O}_X)$ is injective, it follows from \sref{IV.8.12.1} that if $g$ is an immersion, it is an open immersion and $g(X)$ is dense in $Z$, and the same is true for $g_\lambda$.
The fact that \emph{a)} implies \emph{a$'$)} follows also from \sref{IV.8.12.1}, applied by replacing $Z$ by $Y'$ and $g$ by $f'$ (since $Y'$ is finite over $Y$, hence affine over $Y$): indeed if $Y''$ is the closed image of $X$ by $f'$, $Y''$ is finite over $Y$ and $f'$ factorises as $X\xrightarrow{f''}Y''\xrightarrow{j}Y'$, where $j$ is the canonical injection, and $f''$ is an immersion \sref[I]{I.4.1.10}); it follows therefore from \sref{IV.8.12.1} that $f''$ is an open immersion.

The equivalence of \emph{b)} and \emph{c)} follows from \sref{IV.8.12.2} as well as the fact that $g_\lambda$ is then an immersion for $\lambda$ large enough.
It is clear that \emph{c)} implies \emph{a)}, since $Z_\lambda$ is finite over $Y$.
Let us finally show that \emph{a)} implies \emph{c)}.
One has seen above that $Y'$ is the closed image of $X$ by $f'$, and it then follows from \sref{IV.8.12.1} that if one sets $\sh{B}'=u_*(\sh{O}_{Y'})$, the homomorphism $\varphi':\sh{B}'\to\sh{B}=f_*(\sh{O}_X)$ is \emph{injective}.
But as by hypothesis $\sh{B}'$ is a finite $\sh{O}_Y$-algebra, it identifies by definition of $\sh{B}$ with one of the $\sh{O}_Y$-subalgebras $\sh{C}_\lambda$, which proves \emph{c)}.

We will say that a morphism $f:X\to Y$, where $Y$ is quasi-compact and quasi-separated, is \emph{pseudo-finite}, if it is of finite type and satisfies condition \emph{a)} of \sref{IV.8.12.3} (in which case it is necessarily separated).

\begin{corollary}[8.12.4]
\label{IV.8.12.4}
Let $Y$ be a prescheme quasi-compact and quasi-separated, $f:X\to Y$ a morphism.
\oldpage[IV-3]{45}
\begin{enumerate}[label=(\roman*)]
  \item Suppose $f$ pseudo-finite.
Then, for all morphisms $Y'\to Y$, where $Y'$ is quasi-compact and quasi-separated, $f_{(Y')}:X'=X_{(Y')}\to Y'$ is pseudo-finite.
  \item Let $(U_\lambda)$ be a covering of $Y$ by quasi-compact open subsets.
For $f$ to be pseudo-finite, it is necessary and sufficient that for all $\lambda$, the restriction $f_\lambda:f^{-1}(U_\lambda)\to U_\lambda$ of $f$ be a pseudo-finite morphism.
  \item Suppose furthermore $Y$ Noetherian, and $f$ of finite type.
Then, for $f$ to be pseudo-finite, it is necessary and sufficient that, for all $y\in Y$, the morphism $f_y=f\times 1:X_y=X\times_Y\operatorname{Spec}(\sh{O}_y)\to\operatorname{Spec}(\sh{O}_y)$ be pseudo-finite.
\end{enumerate}
\end{corollary}

\begin{proof}
(i) It is clear that $f_{(Y')}$ is of finite type \sref{IV.1.5.4}. Moreover, a factorisation $X\xrightarrow{g}Z\xrightarrow{u}Y$, with $g$ an immersion and $u$ finite, induces a factorisation $X'\xrightarrow{g'}Z'\xrightarrow{u'}Y'$ of $f_{(Y')}$, where $Z'=Z_{(Y')}$, $g'=g_{(Y')}$ and $u'=u_{(Y')}$. The morphism $g'$ is an immersion \sref[I]{I.4.3.2} and $u'$ is finite \sref[II]{II.6.1.5}; hence $f_{(Y')}$ is pseudo-finite.

(ii) The condition is necessary by virtue of (i), the $U_\lambda$ being quasi-separated since $Y$ is.
To see that it is sufficient, let us observe (with the notations of \sref{IV.8.12.3}) that if one sets $X_\lambda=f^{-1}(U_\lambda)$, we have $\sh{B}|U_\lambda=(f_\lambda)_*(\sh{O}_{X_\lambda})$, $\sh{C}|U_\lambda$ is the integral closure of $\sh{O}_{U_\lambda}$ in $\sh{B}|U_\lambda$, and consequently, if $h:Z\to Y$ is the canonical morphism, $Z_\lambda=\operatorname{Spec}(\sh{C}_\lambda)$ identifies with $h^{-1}(U_\lambda)=\operatorname{Spec}(\sh{C}|U_\lambda)$, the $Y$-morphism $g_\lambda:X_\lambda\to Z_\lambda=h^{-1}(U_\lambda)$ being the $U_\lambda$-morphism corresponding to $f_\lambda$.
Applying the hypothesis of induction, one sees that $X''$ is quasi-affine over $Y-\{y\}$; but the open subset $Y-\{y\}$ being quasi-affine over $Y$ since $Y$ is Noetherian, $X''$ is also quasi-affine over $Y$ \sref[II]{II.5.1.10},~(ii)); since on the other hand $X'$ is finite (and \emph{a fortiori} affine) over $Y$, $X$ is quasi-affine over $Y$.

(iii) It suffices, in virtue of (ii), to prove that for all $y\in Y$ there exists a neighbourhood $U$ of $y$ such that the restriction $f^{-1}(U)\to U$ of $f$ be pseudo-finite.
Designate by $(U_\lambda)$ the projective filtering system of affine open neighbourhoods of $y$, and apply the method of \sref{IV.8.1.2},~$a$)).
Since $Y$ is Noetherian, the restrictions $f_\lambda:f^{-1}(U_\lambda)\to U_\lambda$ of $f$ are of finite presentation and of the same nature as $f_y$.
By hypothesis $f_y$ factorises as $X_y\xrightarrow{g_y}Z_y\xrightarrow{u_y}\operatorname{Spec}(\sh{O}_y)$, where $u_y$ is finite and $g_y$ is an open immersion.
As $\sh{O}_y$ is Noetherian, it is of the same nature as $Z_y$, and as $Z_y$ is of finite presentation over $\operatorname{Spec}(\sh{O}_y)$, there exists $\lambda$ and a morphism of finite presentation $u_\lambda:Z_\lambda\to U_\lambda$ such that $Z_y$ identifies with $Z_\lambda\times_{U_\lambda}\operatorname{Spec}(\sh{O}_y)$; furthermore, one can suppose $\lambda$ taken such that $g_\lambda$ is an open immersion and $u_\lambda$ a finite morphism \sref{IV.8.10.5},~(ii) and (x)), which proves that $f_\lambda$ is pseudo-finite.
\end{proof}

\begin{env}[8.12.5]
\label{IV.8.12.5}
We can now give the ``Main theorem'' of Zariski \sref[III]{III.4.4.3}), a proof not using the ``global'' cohomological results of chap.~III, but rather appealing to the finer properties of local Noetherian rings; we will furthermore generalise the statement of the theorem by removing the Noetherian hypotheses:
\end{env}

\begin{theorem}[8.12.6]
\label{IV.8.12.6}
(``Main theorem'' of Zariski). --- Let $Y$ be a prescheme quasi-compact and quasi-separated.
If a morphism $f:X\to Y$ is quasi-finite, separated and of finite presentation, there exists a factorisation of $f$
\begin{equation}
\label{IV.8.12.6.1}
X\xrightarrow{f'}Y'\xrightarrow{u}Y \tag{8.12.6.1}
\end{equation}
where $f'$ is an open immersion and $u$ a finite morphism.
\end{theorem}

\oldpage[IV-3]{46}

\begin{proof}
By virtue of \sref{IV.8.12.4},~(ii)) and the local character (on $Y$) of the notions of quasi-finite morphism, separated and of finite presentation, one can restrict to the case where $Y=\operatorname{Spec}(A)$ is affine.
Applying \sref{IV.8.9.1}), one can suppose that there is a subring $A_0$ of $A$, which is a $\bb{Z}$-algebra of finite type, and an $A$-isomorphism $X_0\otimes_{A_0}A\xrightarrow{\sim}X$, $f$ identifying by this isomorphism with $f_0\times 1$, where $f_0:X_0\to\operatorname{Spec}(A_0)$ is a morphism \emph{of finite type}; furthermore \sref{IV.8.10.5},~(v) and (xi)) one can suppose that $f_0$ is separated and quasi-finite; if one proves that $f_0$ is pseudo-finite, the same will be true of $f$ by \sref{IV.8.12.4},~(i)).
As $A_0$ is then Noetherian and the notions of morphism of finite type, separated and quasi-finite are preserved by base change, it follows from \sref{IV.8.12.4},~(iii)) that one can even suppose that $A$ is a local ring, \emph{essentially of finite type over} $\bb{Z}$ \sref{IV.1.3.8}).
Set $n=\dim(A)$, and let us proceed by induction on $n$; for $n=0$, the theorem is evident, $A$ being a field and the morphism $f$ being already finite \sref[II]{II.6.2.2}).
Set $B=\Gamma(X,\sh{O}_X)$; denote by $C$ the integral closure of $A$ in $B$, set $Z=\operatorname{Spec}(C)$ and let $g:X\to Z$ be the $Y$-morphism corresponding to the canonical $A$-homomorphism injective $C\to B$; by virtue of \sref{IV.8.12.3}), it suffices to show that $g$ is an open immersion.
Let $a$ be the closed point of $Y$, and $U=Y-\{a\}$; $U$ is Noetherian and all its local rings are essentially of finite type over $\bb{Z}$ and of dimension $<n$; taking into account the hypothesis of induction, one sees that the restriction $f^{-1}(U)\to U$ of $f$ is a pseudo-finite morphism.
One concludes from \sref{IV.8.12.3}) that, if $h:Z\to Y$ is the structural morphism, the restriction $f^{-1}(U)\to h^{-1}(U)$ of $g$ is an open immersion.
Set $A'=\hat{A}$, $Y'=\operatorname{Spec}(A')$, $X'=X_{(Y')}$, $f'=f_{(Y')}:X'\to Y'$.
As the canonical morphism $u:Y'\to Y$ is flat, it follows from \sref{IV.2.3.1}) that $B'=\Gamma(X',\sh{O}_{X'})$ identifies with the $A'$-algebra $B\otimes_A A'$.
On the other hand, as $A$ is a local ring excellent \sref{IV.7.8.3}), the morphism $u:Y'\to Y$ is regular, and consequently \sref{IV.6.14.4}) \emph{the integral closure of $A'$ in $B'$ is equal to} $C'=C\otimes_A A'$.
One sees therefore that $Z'=\operatorname{Spec}(C')$ is equal to $Z_{(Y')}$ and the morphism $g':X'\to Z'$ coming from the canonical injection $C'\to B'$ is equal to $g_{(Y')}$.
As $u:Y'\to Y$ is faithfully flat and quasi-compact, to prove that $g$ is an open immersion, it suffices to prove that $g'$ is an open immersion \sref{IV.2.7.1},~(x)).
If $h':Z'\to Y'$ is the canonical morphism, the fact that the restriction $f'^{-1}(U')\to h'^{-1}(U')$ of $g$ is an open immersion implies that the same is true of the restriction $f'^{-1}(U')\to h'^{-1}(U')$ of $g'$.
Now let us note that $u^{-1}(a)$ is reduced to the closed point $a'$ of $Y'$ and consequently $U'=Y'-\{a'\}=u^{-1}(U)$.
If $h':Z'\to Y'$ is the canonical morphism, it follows from \sref[II]{II.6.2.4}); as $A'$ is \emph{complete}, one deduces from \sref[II]{II.6.2.6}) that $X'$ is $Y'$-isomorphic to a sum $X'_1\sqcup X'_2$, where the restriction $f'|X'_1:X'_1\to Y'$ is a morphism \emph{finite}, and $X'_2\subset f'^{-1}(U')$.
It follows from this that $B'$ is the direct product of the two $A'$-algebras $\Gamma(X'_1,\sh{O}_{X'_1})=B'_1$ and $\Gamma(X'_2,\sh{O}_{X'_2})=B'_2$; one concludes immediately that the integral closure $C'$ of $A'$ in $B'$ is the direct product of the integral closures $C'_1$, $C'_2$ of $A'$ in $B'_1$, $B'_2$ respectively, whence $Z'=Z'_1\sqcup Z'_2$, where $Z'_i=\operatorname{Spec}(C'_i)$ ($i=1,2$); and the canonical morphism $g':X'\to Z'$ is such that $g'|X'_i$ is the canonical morphism $g'_i:X'_i\to Z'_i$ ($i=1,2$).
But as $B'_1$ is already a finite $A'$-algebra, one has $C'_1=B'_1$, and $g'_1$ is therefore an isomorphism.
On the other hand, as $X'_2\subset f'^{-1}(U')$ and is open in $f'^{-1}(U')$, one
\oldpage[IV-3]{47}
already knows that $g'_2$ is an open immersion.
One concludes that $g'$ is indeed an open immersion.
\end{proof}

\begin{remark}[8.12.7]
\label{IV.8.12.7}
When, in \sref{IV.8.12.6}), one supposes that $Y$ is an affine \emph{scheme}, the proof by reduction to the Noetherian case shows that, in the factorisation \sref{IV.8.12.6.1}), the morphisms $f'$ and $u$ are also morphisms \emph{of finite presentation} \sref{IV.1.6.2}).
\end{remark}

\begin{corollary}[8.12.8]
\label{IV.8.12.8}
Let $Y$ be a quasi-compact scheme such that there exists an $\sh{O}_Y$-Module ample \sref[II]{II.4.5.3}), $f:X\to Y$ a quasi-finite and quasi-projective morphism.
Then there exists a factorisation of $f$ as
\[
X\xrightarrow{f'}Y'\xrightarrow{u}Y
\]
where $f'$ is an open immersion and $u$ a finite morphism.
\end{corollary}

\begin{proof}
The hypothesis implies that $X$ identifies with a sub-$Y$-scheme quasi-compact of a $Y$-scheme of the form $Z=\mathbf{P}^r_Y$ \sref[II]{II.5.3.3}).
There is therefore a quasi-compact open neighbourhood $U$ of $X$ in $Z$ such that $X$ is closed in $U$; as $Z$ is a scheme, the canonical injection $U\to Z$ is a morphism of finite presentation (\sref{IV.1.2.7} and \sref{IV.1.6.2})), hence the composite morphism $g:U\to Z\to Y$ is also a morphism of finite presentation (the fact that $\mathbf{P}^r_Y$ is of finite presentation over $Y$ following immediately from the definition \sref[II]{II.4.1.1})).
Let $\sh{J}$ be the quasi-coherent Ideal of $\sh{O}_U$ defining the closed sub-prescheme $X$; as $U$ is a quasi-compact scheme, $\sh{J}$ is the filtered inductive limit of its sub-Ideals quasi-coherent of finite type $\sh{J}_\lambda$ \sref[I]{I.9.4.9}).
If $X_\lambda$ is the closed sub-prescheme of $U$ defined by $\sh{J}_\lambda$, one therefore has $X=\bigcap X_\lambda$.
For all $y\in Y$, one therefore has $f^{-1}(y)=\bigcap(X_\lambda\cap g^{-1}(y))$, and as the sets $X_\lambda\cap g^{-1}(y)$ are closed in the Noetherian space $g^{-1}(y)$, there exists for all $y$ an index $\lambda(y)$ such that $f^{-1}(y)=X_{\lambda(y)}\cap g^{-1}(y)$.
Denote by $E_\lambda$ the set of $y\in Y$ such that the fibre $X_\lambda\cap g^{-1}(y)$ of the restriction of $g$ to $X_\lambda$ is a $k(y)$-prescheme finite.
The hypothesis that $f$ is quasi-finite implies, by virtue of what precedes, that $Y=\bigcup_\lambda E_\lambda$.
Now, each of the $X_\lambda$ is, by definition, \emph{of finite presentation} over $Y$; it follows therefore from \sref{IV.9.2.3} and \sref{IV.9.2.6})~($^1$) that the $E_\lambda$ are \emph{constructible} sets in the scheme $Y$; as they form a filtered increasing family, there exists an index $\lambda$ such that $E_\lambda=Y$ \sref{IV.1.9.9}), and for this $\lambda$, the morphism $f_\lambda:X_\lambda\to Y$, restriction of $g$ to $X_\lambda$, is of finite presentation, finite and separated, hence one can apply \sref{IV.8.12.6}), and $f_\lambda$ factorises as
\[
X_\lambda\xrightarrow{j_\lambda}Y'_\lambda\xrightarrow{u_\lambda}Y
\]
where $j_\lambda$ is an open immersion and $u_\lambda$ a finite morphism.
As $X$ is a closed sub-prescheme of $X_\lambda$, one has thus proved that $f$ possesses the property \sref{IV.8.12.3},~$a$)), whence the corollary by virtue of the equivalence of \sref{IV.8.12.3},~$a$)) and \sref{IV.8.12.3},~$a'$)).
\end{proof}

\footnotetext{($^1$) The reader will verify that the corollaries \sref{IV.8.12.8} to \sref{IV.8.12.11} are not used in \textsection~9.}

\begin{corollary}[8.12.9]
\label{IV.8.12.9}
\oldpage[IV-3]{48}
Let $f:X\to Y$ be a locally quasi-finite morphism $(\mathbf{Err}_{\text{III}},~20)$.
For all $x\in X$ there exist an open neighbourhood $U$ of $x$ in $X$, an open neighbourhood $V$ of $y=f(x)$ in $Y$, such that $f(U)\subset V$ and a factorisation
\[
U\xrightarrow{f'}V'\xrightarrow{u}V
\]
of the restriction of $f$ to $U$, where $f'$ is an open immersion and $u$ a finite morphism.
\end{corollary}

\begin{proof}
It evidently suffices to take for $V$ an affine neighbourhood of $y$ in $Y$, for $U$ an affine neighbourhood of $x$ in $X$ contained in $f^{-1}(V)$ and such that $f|U$ is quasi-finite.
The morphism $U\to V$ restriction of $f$ being then affine (hence quasi-projective), one can apply \sref{IV.8.12.8}).
\end{proof}

\begin{corollary}[8.12.10]
\label{IV.8.12.10}
Let $Y$ be an integral normal prescheme, $X$ an integral prescheme, and $f:X\to Y$ a birational, locally quasi-finite morphism $(\mathbf{Err}_{\text{III}},~20)$.
Then $f$ is an open immersion ouverte and furthermore it is necessary and sufficient that $f$ be separated.
\end{corollary}

\begin{proof}
The second assertion follows immediately from the first and from \sref[I]{I.8.2.8}).
To prove the first assertion, one may suppose $X$ and $Y$ affine and $f$ quasi-finite; consider the factorisation $f=u\circ f'$ of \sref{IV.8.12.8}, which identifies $X$, via $f'$, with the subprescheme induced on an open subset of $Y'$.
As $X$ is integral, on can, by virtue of \sref[I]{I.5.2.3}), replace $Y'$ by the reduced sub-prescheme of $Y'$ having underlying space $\overline{X}$, hence one can also suppose that $Y'$ is \emph{integral}.
Furthermore, since $f$ is birational, $u$ is also.
The conclusion then follows from the following lemma:
\end{proof}

\begin{lemma}[8.12.10.1]
\label{IV.8.12.10.1}
Let $Y'$ be an integral prescheme, $Y$ an integral prescheme and normal; then a morphism $u:Y'\to Y$ finite and birational is an isomorphism.
\end{lemma}

\begin{proof}
Set $\sh{A}=u_*(\sh{O}_{Y'})$; then $\sh{A}$ is a finite $\sh{O}_Y$-algebra and $Y'$ identifies with $\operatorname{Spec}(\sh{A})$ \sref[II]{II.1.3.6}.
If $R(Y)$ is the field of rational functions of $Y$, one therefore has, for all $y\in Y$, $\sh{O}_{Y,y}\subset\sh{A}_y\subset R(Y)$; but as the ring $\sh{O}_{Y,y}$ is by hypothesis integrally closed and $R(Y)$ a field of fractions, one necessarily has $\sh{A}_y=\sh{O}_{Y,y}$, whence the lemma.
\end{proof}

\begin{corollary}[8.12.11]
\label{IV.8.12.11}
Let $Y$ and $X$ be integral normal preschemes, and let $f:X\to Y$ be a dominant, locally quasi-finite morphism.
Let $K$ and $L$ (extension of $K$) be the fields of rational functions of $Y$ and $X$ respectively; and let $Y'$ be the integral closure of $Y$ relative to $L$ \sref[II]{II.6.3.4}); then $f$ factorises in a unique way as $f:X\xrightarrow{f'}Y'\xrightarrow{u}Y$, where $f'$ is birational, and correspond to the identity automorphism of $L$; $f'$ is then a local isomorphism, and for $f'$ to be an open immersion, it is necessary and sufficient that $f$ be separated.
\end{corollary}

\begin{proof}
The existence and uniqueness of the factorisation of $f$ follow from \sref[II]{II.6.3.9}).
It follows further from \sref[II]{II.6.2.4},~(v)), reducing to the affine case, that $f'$ is locally quasi-finite; furthermore, it follows from \sref[I]{I.5.5.1}) that, for $f'$ to be separated, it is necessary and sufficient that $f$ be, since $u$ is affine, hence separated; the last two assertions are therefore consequences of \sref{IV.8.12.10}) applied to $f'$.
\end{proof}

\subsection{Translation in terms of pro-objects}
\label{subsection:IV.8.13}

\oldpage[IV-3]{49}

The following proposition is essentially equivalent to \sref{IV.8.8.2},~(i)):

\begin{proposition}[8.13.1]
\label{IV.8.13.1}
Let $S$ be a prescheme, $(X_\lambda,v_{\lambda\mu})$ a projective filtering system of $S$-preschemes; suppose that there exists $\alpha$ such that $v_{\alpha\lambda}$ is an affine morphism for all $\lambda\geqslant\alpha$ (which implies \sref[II]{II.1.6.2}) that $v_{\lambda\mu}$ is affine for $\alpha\leqslant\lambda\leqslant\mu$), so that the projective limit $X=\varprojlim X_\lambda$ exists in the category of $S$-preschemes \sref{IV.8.2.3}).
Let $Y$ be an $S$-prescheme, and, for all $\lambda\geqslant\alpha$, let $e_\lambda:\operatorname{Hom}_S(X_\lambda,Y)\to\operatorname{Hom}_S(X,Y)$ be the map which, to every $S$-morphism $f_\lambda:X_\lambda\to Y$, makes correspond $f=f_\lambda\circ v_\lambda$, where $v_\lambda:X\to X_\lambda$ is the canonical morphism.
The family $(e_\lambda)$ is an inductive system of maps, which therefore defines a canonical map
\begin{equation}
\label{IV.8.13.1.1}
\varinjlim\operatorname{Hom}_S(X_\lambda,Y)\to\operatorname{Hom}_S(X,Y). \tag{8.13.1.1}
\end{equation}

Suppose $X_\alpha$ quasi-compact (resp. quasi-compact and quasi-separated), and the structural morphism $Y\to S$ locally of finite type (resp. locally of finite presentation).
Then the map \sref{IV.8.13.1.1} is injective (resp. bijective).
\end{proposition}

\begin{proof}
Setting, for $\lambda\geqslant\alpha$, $Z_\lambda=Y\times_S X_\lambda$, so that one has $Z_\lambda=Z_\alpha\times_{X_\alpha}X_\lambda$.
Setting similarly $Z=Y\times_S X=Z_\alpha\times_{X_\alpha}X$; one knows then \sref{IV.8.2.5}) that, if one also sets $w_{\lambda\mu}=1\times v_{\lambda\mu}:Z_\mu\to Z_\lambda$ for $\alpha\leqslant\lambda\leqslant\mu$ and $w_\lambda=1\times v_\lambda:Z\to Z_\lambda$ for $\alpha\leqslant\lambda$, $Z$ is the projective limit of the system $(Z_\lambda,w_{\lambda\mu})$ and the $w_\lambda$ are the canonical morphisms.
Let us note furthermore that the morphism $Z_\alpha\to X_\alpha$ is locally of finite type (resp. locally of finite presentation) \sref{IV.1.3.4} and \sref{IV.1.4.3}).
Finally, one knows that one has \sref[I]{I.3.3.14})
\[
\operatorname{Hom}_S(X_\lambda,Y)=\operatorname{Hom}_{X_\lambda}(X_\lambda,Z_\lambda) \quad\text{and}\quad \operatorname{Hom}_S(X,Y)=\operatorname{Hom}_X(X,Z)
\]
It suffices now to apply \sref{IV.8.8.2},~(i)) by setting $X_\lambda=S_\lambda$ and replacing $Y_\lambda$ by $Z_\lambda$.
\end{proof}

\begin{corollary}[8.13.2]
\label{IV.8.13.2}
With the notations of \sref{IV.8.13.1}), suppose $X_\alpha$ quasi-compact and quasi-separated, and the $v_{\alpha\lambda}$ affine for $\alpha\leqslant\lambda$; suppose furthermore that $Y=\varprojlim Y_\rho$, where $(Y_\rho,t_{\rho o})$ is a projective filtering system of $S$-preschemes such that, for each $\rho$, the structural morphism $Y_\rho\to S$ is locally of finite presentation.
Then one has a canonical bijection
\begin{equation}
\label{IV.8.13.2.1}
\operatorname{Hom}_S(X,Y)\xleftarrow{\sim}\varprojlim_\rho\left(\varinjlim_\lambda\operatorname{Hom}_S(X_\lambda,Y_\rho)\right). \tag{8.13.2.1}
\end{equation}
\end{corollary}

\begin{proof}
Indeed, the fact that $Y$ is the projective limit of the $Y_\rho$ implies in particular that the canonical map $\operatorname{Hom}_S(X,Y)\to\varprojlim\operatorname{Hom}_S(X,Y_\rho)$ is bijective; and on the other hand, the hypotheses imply, for each $\rho$, the existence of a canonical bijection $\operatorname{Hom}_S(X,Y_\rho)\xleftarrow{\sim}\varinjlim_\lambda\operatorname{Hom}_S(S_\lambda,Y_\rho)$ by virtue of \sref{IV.8.13.1}); whence the conclusion.
\end{proof}

\begin{env}[8.13.3]
\label{IV.8.13.3}
The preceding results allow one to interpret in the theory of preschemes the notions of ``pro-variety'' or of ``pro-scheme'' which arise in certain applications (for example in the theory of the class field body following the ideas of Serre~[39] or in the theory of N\'eron on the reduction of
\oldpage[IV-3]{50}
abelian varieties~[32]).
Let us recall here rapidly the notion of \emph{pro-object} of a category, referring to chap.~V for more ample developments (we will moreover not use before chap.~V the interpretation that follows, and the reader can therefore omit it until the end of this number).
Given a category $\mathcal{C}$, the category $\mathbf{Pro}(\mathcal{C})$ of \emph{pro-objects} of $\mathcal{C}$ has for objects the \emph{projective systems} (in the universe where one places oneself) $\mathbf{X}=(X_\mu)_{\mu\in M}$ of objects of $\mathcal{C}$ whose index sets (depending on the projective system considered) are supposed filtered preordered.
Given two such pro-objects $\mathbf{X}=(X_\mu)_{\mu\in M}$, $\mathbf{X}'=(X'_{\mu'})_{\mu'\in M'}$, the \emph{morphisms} from $\mathbf{X}$ to $\mathbf{X}'$ are by definition the elements of the set $\varprojlim_{\mu'}(\varinjlim_\mu\operatorname{Hom}(X_\mu,X_{\mu'}))$; the verification that one can take these sets as sets of morphisms is immediate, the composition of systems of morphisms $u''_\mu:X_\mu\to X'_{\mu'}$, $u''_{\mu'}:X'_{\mu'}\to X''_{\mu''}$ which are inductive and projective with respect to the lower index and upper index respectively, being carried out ``argument by argument'', in other words considering the system of $u''^{\mu'}_{\mu''}\circ u''^{\mu}_{\mu'}$.
\end{env}

\begin{env}[8.13.4]
\label{IV.8.13.4}
Consider then a prescheme quasi-compact and quasi-separated $S$, and denote by $\mathcal{C}$ the full subcategory of the category $(\mathbf{Sch})_{/S}$ of $S$-preschemes, formed of the $S$-preschemes $X$ having the following property: the structural morphism $X\to S$ factorises as $X\xrightarrow{g}X_0\xrightarrow{f}S$, where $g:X\to X_0$ is affine and $f:X_0\to S$ is \emph{of finite presentation}; we will say for short that the preschemes of $\mathcal{C}$ are \emph{essentially affine over} $S$.
Consider on the other hand the full subcategory $\mathcal{C}'_0$ of $(\mathbf{Sch})_{/S}$ formed of the $S$-preschemes \emph{of finite presentation}, and the category $\mathbf{Pro}(\mathcal{C}'_0)$ of pro-objects of $\mathcal{C}'_0$.
We will say that an object $\mathbf{X}=(X_\mu)_{\mu\in M}$ of $\mathbf{Pro}(\mathcal{C}'_0)$ is \emph{essentially affine} if there exists a $\gamma\in M$ such that for $\mu\geqslant\gamma$, the transition morphism $v_{\gamma\mu}:X_\mu\to X_\gamma$ is affine (which implies that for $\gamma\leqslant\nu\leqslant\mu$, $v_{\nu\mu}:X_\mu\to X_\nu$ is affine).
One will note that an object of $\mathbf{Pro}(\mathcal{C}'_0)$ isomorphic to an essentially affine object \emph{is not} necessarily itself essentially affine.
We will denote by $\mathcal{C}'$ the full subcategory of $\mathbf{Pro}(\mathcal{C}'_0)$ formed of pro-objects essentially affine.

This being so, it follows from \sref{IV.8.2.2} and \sref{IV.8.2.3}) that for all objects $\mathbf{X}=(X_\mu)_{\mu\in M}$ of $\mathcal{C}'$, the $S$-prescheme $X=\varprojlim_\mu X_\mu$ exists; furthermore, as, for $\mu$ large enough, the morphism $X\to X_\mu$ is affine \sref{IV.8.2.2}), $X$ is \emph{essentially affine over} $S$ by definition.
Set $X=L(\mathbf{X})$; let us show that one has thus defined a canonical \emph{functor}
\begin{equation}
\label{IV.8.13.4.1}
L:\mathcal{C}'\to\mathcal{C} \tag{8.13.4.1}
\end{equation}

\oldpage[IV-3]{51}

One has indeed, for two objects $\mathbf{X}=(X_\mu)$, $\mathbf{X}'=(X'_{\mu'})$ of $\mathcal{C}'$, a canonical map
\[
\varprojlim_{\mu'}\varinjlim_\mu\operatorname{Hom}_S(X_\mu,X'_{\mu'})\to\operatorname{Hom}_S(\varprojlim_\mu X_\mu,\varprojlim_{\mu'}X'_{\mu'})
\]
defined in \sref{IV.8.13.1.1}), and on the other hand, by definition of the projective limit, a canonical bijection
\[
\varprojlim_{\mu'}\operatorname{Hom}_S(\varprojlim_\mu X_\mu,X'_{\mu'})\xrightarrow{\sim}\operatorname{Hom}_S(\varprojlim_\mu X_\mu,\varprojlim_{\mu'}X'_{\mu'})
\]
\end{env}

whence a canonical map
\begin{equation}
\label{IV.8.13.4.2}
\varprojlim_{\mu'}\varinjlim_\mu\operatorname{Hom}_S(X_\mu,X'_{\mu'})\to\operatorname{Hom}_S(\varprojlim_\mu X_\mu,\varprojlim_{\mu'}X'_{\mu'}) \tag{8.13.4.2}
\end{equation}
evidently functorial in $\mathbf{X}$ and $\mathbf{X}'$, and which completes the definition of the functor $L$.

\begin{proposition}[8.13.5]
\label{IV.8.13.5}
The hypotheses and notations being those of \sref{IV.8.13.4}), the functor $L$ is pleinly faithful.
If furthermore $S$ is a Noetherian prescheme (which already implies that $S$ is quasi-compact and quasi-separated \sref{IV.1.2.8})), $L$ is \emph{an equivalence of categories}.
\end{proposition}

\begin{proof}
To say that $L$ is pleinly faithful signifies that the map \sref{IV.8.13.4.2}) is bijective, which is a particular case of \sref{IV.8.13.2}).
Indeed, the structural morphisms $X_\mu\to S$ being of finite presentation, are in particular quasi-compact and quasi-separated, hence the $X_\mu$ are quasi-compact and quasi-separated.

To show that when $S$ is Noetherian $L$ is an equivalence of categories, it suffices, since one already knows that $L$ is pleinly faithful, to prove that every $S$-prescheme essentially affine $X$ is $S$-\emph{isomorphic} to an object of the form $L(\mathbf{X})$ where $\mathbf{X}\in\mathcal{C}'$ $(\mathbf{0}_{\text{III}},~8.1.5)$.
Now, by hypothesis there is a factorisation $X\xrightarrow{g}X_0\xrightarrow{f}S$ of the structural morphism, $f$ being of finite presentation and $g$ affine.
One can therefore write $X=\operatorname{Spec}(\sh{A})$, where $\sh{A}$ is an $\sh{O}_{X_0}$-Algebra quasi-coherent \sref[II]{II.1.3.1}).
Now, as $X_0$ is Noetherian (since it is of finite type over $S$ and $S$ is Noetherian), $\sh{A}$ is the filtered inductive limit of the family $(\sh{A}_\lambda)$ of its sub-$\sh{O}_{X_0}$-Algebras quasi-coherent of finite type \sref[I]{I.9.6.6}).
Set $X_\lambda=\operatorname{Spec}(\sh{A}_\lambda)$; the morphisms $X_\lambda\to X_0$ are of finite type, hence of finite presentation since $X_0$ is Noetherian, and consequently the same is true of the composite morphisms $X_\lambda\to X_0\to S$ \sref{IV.1.6.2}); in other words, the $X_\lambda$ belong to $\mathcal{C}'_0$ and as the morphisms $X_\lambda\to X_0$ are affine, $\mathbf{X}=(X_\lambda)$ is an object of $\mathcal{C}'$ whose projective limit exists and is $S$-isomorphic to $X$ by virtue of \sref{IV.8.2.2}).
This completes the proof.
\end{proof}

\begin{remark}[8.13.6]
\label{IV.8.13.6}
It follows from \sref{IV.1.6.2} and \sref[II]{II.1.6.2}) that if $X$ and $Y$ are essentially affine over $S$, the same is true of $X\times_S Y$.
One concludes for example $(\mathbf{0}_{\text{III}},~8.2.5)$ that a $\mathcal{C}$-group is nothing other than a $(\mathbf{Sch})_{/S}$-group which is a prescheme essentially affine over $S$.
On the other hand, the finite products exist in the category $\mathcal{C}'$: indeed, if $\mathbf{X}=(X_\mu)_{\mu\in M}$, $\mathbf{Y}=(Y_\rho)_{\rho\in R}$ are two objects of $\mathcal{C}'$, the products $X_\mu\times_S Y_\rho$ are $S$-preschemes of finite presentation, and taking for transition morphisms $X_\nu\times_S Y_\sigma\to X_\mu\times_S Y_\rho$ the products of the transition morphisms $X_\nu\to X_\mu$ and $Y_\sigma\to Y_\rho$, one sees immediately that $(X_\mu\times_S Y_\rho)$ is a product of $\mathbf{X}$ and $\mathbf{Y}$ in $\mathbf{Pro}(\mathcal{C}'_0)$; furthermore \sref[II]{II.1.6.2}) the transition morphisms thus defined are affine for $\mu$ and $\rho$ large enough, hence the product $\mathbf{X}\times\mathbf{Y}$ also belongs to $\mathcal{C}'$.
One concludes then as above that a $\mathcal{C}'$-group is a $\mathbf{Pro}(\mathcal{C}'_0)$-group which is essentially affine.
One deduces from \sref{IV.8.13.5}) that the categories of the $\mathcal{C}'$-groups and the $\mathcal{C}$-groups are \emph{equivalent} when $S$ is \emph{Noetherian}.
It seems plausible that when $S$ is the spectrum of a field $k$, the category of $\mathcal{C}$-groups is \emph{equivalent} to that of the $\boldsymbol{K}$-\emph{groups}, where $\boldsymbol{K}$ is the category of $S$-preschemes \emph{quasi-compact}; in other words, every prescheme in groups quasi-compact over $k$ would be essentially affine.
On the other hand, if one denotes
\oldpage[IV-3]{52}
by $\mathcal{C}'_0$-$\boldsymbol{gr}$ the category of the $\mathcal{C}'_0$-groups, it is very likely that the category of the $\mathcal{C}'$-groups is equivalent to the full subcategory of $\mathbf{Pro}(\mathcal{C}'_0\text{-}\boldsymbol{gr})$ formed of the ``proalgebraic groups essentially affine'', that is to say the projective systems $\mathbf{G}=(G_\mu)_{\mu\in M}$, where the $G_\mu$ are algebraic groups over $k$ and the transition morphisms $G_\nu\to G_\mu$ are affine for $\mu$ large enough (which is what one can also express by saying that $\mathbf{G}$ is an extension of an algebraic group by a pro-algebraic group affine).
The conjunction of these two conjectures is moreover equivalent to the following: every prescheme in groups quasi-compact over $k$ (i.e.\ a prescheme in groups of finite type over $k$) is an extension by a prescheme in groups \emph{affine} over $k$.
The only pro-algebraic groups encountered in practice until now being in fact essentially affine, it would therefore be advantageous to substitute for the study of general pro-algebraic groups (introduced and studied by Serre~[40]) that of schemes in groups quasi-compact over $k$, whose definition is conceptually much simpler.
\end{remark}

\subsection{Characterisation of a prescheme locally of finite presentation over another, in terms of the functor it represents}
\label{subsection:IV.8.14}

\begin{env}[8.14.1]
\label{IV.8.14.1}
Given a prescheme $S$, we will say still, as in \sref{IV.8.13.4}), that a projective filtering system $(X_\lambda,v_{\lambda\mu})$ of $S$-preschemes is \emph{essentially affine} if there exists $\alpha$ such that $v_{\alpha\lambda}$ is an affine morphism for $\lambda\geqslant\alpha$.

The following statement, which will be especially useful in chap.~V, makes precise \sref{IV.8.8.2},~(i)) by providing it with a converse:
\end{env}

\begin{proposition}[8.14.2]
\label{IV.8.14.2}
Let $S$ be a prescheme, $f:X\to S$ a morphism.
For every $S$-prescheme $T$, set
\[
h_X(T)=\operatorname{Hom}_S(T,X)
\]
so that $h_X$ is a contravariant functor from the category $(\mathbf{Sch})_{/S}$ of $S$-preschemes to the category $\mathbf{Ens}$ of sets $(\mathbf{0}_{\text{III}},~8.1.1)$, and $X$ an object representing this functor $(\mathbf{0}_{\text{III}},~8.1.8)$.
The following conditions are equivalent:
\begin{enumerate}[label=\emph{\alph*)}]
  \item $f$ is locally of finite presentation.
  \item For every projective filtering system $(Z_\lambda)$ of $S$-preschemes, essentially affine \sref{IV.8.13.4}) and formed of preschemes quasi-compact and quasi-separated, the canonical map \sref{IV.8.13.1.1})
\begin{equation}
\label{IV.8.14.2.1}
\varinjlim h_X(Z_\lambda)\to h_X(\varprojlim Z_\lambda) \tag{8.14.2.1}
\end{equation}
is bijective.
  \item For every projective filtering system $(Z_\lambda)$ of $S$-preschemes, such that the $Z_\lambda$ are affine schemes, the map \sref{IV.8.14.2.1}) is bijective.
  \item[\rm c$'$)] For every affine open subset $U$ of $S$ and every projective filtering system $(Z_\lambda)$ of $U$-preschemes, such that the $Z_\lambda$ are affine schemes, the map \sref{IV.8.14.2.1}) is bijective.
\end{enumerate}
\end{proposition}

\begin{proof}
The fact that \emph{a)} implies \emph{b)} is nothing other than \sref{IV.8.13.1}); it is trivial that \emph{b)} implies \emph{c)} and that \emph{c)} implies \emph{c$'$)}.
It remains to see that \emph{c$'$)} implies \emph{a)}, and as the property \emph{a)} is local on $S$, one can restrict to the case where $S$ is \emph{affine}.
\oldpage[IV-3]{53}

Suppose first that $X$ is also affine; the assertion to prove is then equivalent to the following:

\begin{corollary}[8.14.2.2]
\label{IV.8.14.2.2}
Let $A$ be a ring, $B$ an $A$-algebra.
In order that, for every inductive filtering system $(C_\lambda)$ of $A$-algebras, the canonical map
\begin{equation}
\label{IV.8.14.2.3}
\varinjlim\operatorname{Hom}_{A\text{-}\mathrm{alg.}}(B,C_\lambda)\to\operatorname{Hom}_{A\text{-}\mathrm{alg.}}(B,\varinjlim C_\lambda) \tag{8.14.2.3}
\end{equation}
be bijective, it is necessary and sufficient that $B$ be an $A$-algebra of finite presentation.
\end{corollary}

It remains only to show that the condition is necessary.
First, for $(C_\lambda)$ take the inductive filtering system of sub-$A$-algebras of finite type of $B$, so that $\varinjlim C_\lambda=B$.
The fact that \sref{IV.8.14.2.3}) is bijective implies in particular that the identity map $1_B:B\to B$ factorises as $B\to C_\lambda\to B$ for a suitable $\lambda$, which implies that $C_\lambda=B$, hence $B$ is an $A$-algebra \emph{of finite type}.
Write then $B=C/\mathfrak{J}$, where $C=A[T_1,\ldots,T_n]$ and $\mathfrak{J}$ is an ideal of $C$; setting $C_\lambda=C/\mathfrak{J}_\lambda$, where $\mathfrak{J}_\lambda$ ranges over the ideals of finite type $\mathfrak{J}_\lambda\subset\mathfrak{J}$ of $C$; then $\mathfrak{J}$ is the filtered inductive limit of the ideals of finite type of $\mathfrak{J}$, and using the exactness of the inductive limit functor, one sees that $B$ is still isomorphic to the inductive limit of the inductive system $(C_\lambda)$.
There exists therefore a $\lambda$ and an $A$-homomorphism $u:B\to C_\lambda$ such that the composite $B\xrightarrow{u}C_\lambda\xrightarrow{p_\lambda}B$ (where $p_\lambda$ is the canonical homomorphism) is the identity.
Let $q_\lambda:C\to C_\lambda$ be the canonical homomorphism, and set
$t_i=p_\lambda(q_\lambda(T_i))$.
Then $p_\lambda(u(t_i))=p_\lambda(q_\lambda(T_i))$, in other words
$u(t_i)-q_\lambda(T_i)\in\mathfrak{J}/\mathfrak{J}_\lambda$.
There is therefore a $\mu\geqslant\lambda$ such that these $n$ elements belong to
$\mathfrak{J}_\mu/\mathfrak{J}_\lambda$ ($1\leqslant i\leqslant n$).
If $p_{\lambda\mu}:C_\lambda\to C_\mu$ is the canonical homomorphism, then
$p_{\lambda\mu}(u(t_i))=p_{\lambda\mu}(q_\lambda(T_i))=q_\mu(T_i)$.
Replacing $\lambda$ by $\mu$ and $u$ by $p_{\lambda\mu}\circ u$, we may therefore
suppose that $u(t_i)=q_\lambda(T_i)$ for every $i$; and if
$r=p_\lambda\circ q_\lambda$ is the canonical homomorphism
$C\to C/\mathfrak{J}=B$, we may write $u(r(T_i))=q_\lambda(T_i)$ for every $i$,
whence $q_\lambda=u\circ r$.
But this implies necessarily that $\mathfrak{J}=\mathfrak{J}_\lambda$, for if $z\in\mathfrak{J}$, one has $r(z)=0$; hence we have $B=C_\lambda$.

Let us now pass to the case where $S$ is affine and $X$ arbitrary; what amounts to proving is that an affine open subset $V$ of $X$ is of finite presentation over $S$, and by virtue of what has just been shown, it suffices to prove that for every projective filtering system $(Z_\lambda)$ of affine $S$-preschemes, the map
\begin{equation}
\label{IV.8.14.2.4}
\varinjlim\operatorname{Hom}_S(Z_\lambda,V)\to\operatorname{Hom}_S(\varprojlim Z_\lambda,V) \tag{8.14.2.4}
\end{equation}
is bijective.
It is immediate that this map is injective, for if $(v_\lambda)$, $(v'_\lambda)$ are two inductive systems of $S$-homomorphisms $v_\lambda:Z_\lambda\to V$, $v'_\lambda:Z_\lambda\to V$ such that the corresponding morphisms
\[
Z\xrightarrow{u_\lambda}Z_\lambda\xrightarrow{v_\lambda}V,\qquad
Z\xrightarrow{u_\lambda}Z_\lambda\xrightarrow{v'_\lambda}V
\]
are equal ($u_\lambda$ being the canonical morphism), then the morphisms
\[
Z\xrightarrow{u_\lambda}Z_\lambda\xrightarrow{v_\lambda}V\xrightarrow{j}X,\qquad Z\xrightarrow{u_\lambda}Z_\lambda\xrightarrow{v'_\lambda}V\xrightarrow{j}X
\]
(where $j$ is the canonical injection) are equal, which implies $j\circ v_\lambda=j\circ v'_\lambda$ by hypothesis for a suitable $\lambda$, hence $v_\lambda=v'_\lambda$.

It remains to prove that \sref{IV.8.14.2.4} is surjective.
Let $v:Z\to V$ be an $S$-morphism.
By hypothesis there are a $\lambda$ and an $S$-morphism
$w_\lambda:Z_\lambda\to X$ such that $j\circ v$ factors as
\[
 Z\xrightarrow{u_\lambda}Z_\lambda\xrightarrow{w_\lambda}X.
\]
It suffices to prove that there is a $\mu\geqslant\lambda$ such that
\[
 Z_\mu\xrightarrow{w_\lambda\circ u_{\lambda\mu}}X
\]
(where $u_{\lambda\mu}$ is the transition morphism) factors as
\[
 Z_\mu\xrightarrow{v_\mu}V\xrightarrow{j}X.
\]
Set $U_\lambda=w_\lambda^{-1}(V)$.
Then
\[
 u_\lambda^{-1}(U_\lambda)
 =u_\lambda^{-1}(w_\lambda^{-1}(V))
 =v^{-1}(V)=Z.
\]
The $Z_\lambda$ are quasi-compact and the $U_\lambda$, being open, are
ind-constructible \sref{IV.1.9.6}; hence \sref{IV.8.3.4} gives a
$\mu\geqslant\lambda$ for which $U_\mu=Z_\mu$.
\end{proof}

\begin{remark}[8.14.3]
\label{IV.8.14.3}
The fact that the map \sref{IV.8.14.2.1} is injective when $f$ is locally of finite
type \sref{IV.8.8.2},~(i) naturally leads one to ask whether this result has a converse.
It does not, even when $S$ and $X$ are affine, since there exist monomorphisms
$X\to S$ which are not of finite type \sref[I]{I.2.4.2}, and these disprove such a
conjecture.
\end{remark}
